% Centre of Mass of Uniform Rigid Bodies — Pure Mathematics with Mechanics Upper Sixth — Cours signé OnBuch+
% Official MINESEC syllabus. Compiler avec : tectonic PMM13-centre-of-mass-of-uniform-rigid-bodies.tex
\newcommand{\DOCMATIERE}{Pure Mathematics with Mechanics}
\newcommand{\DOCNIVEAU}{Upper Sixth}
\newcommand{\DOCMODULE}{Upper Sixth — Pure mathematics with mechanics}
\newcommand{\DOCLECON}{13}
\newcommand{\DOCTITRE}{Centre of Mass of Uniform Rigid Bodies}
% OnBuch+ — préambule « cours signés » ENRICHI (compilé avec tectonic / XeLaTeX)
% Système visuel « Pop » d'OnBuch+ : fond crème (jamais blanc), bordures encre
% épaisses, ombres « dures » décalées, titres Archivo Black en capitales.
% Version MATHÉMATIQUES : graphes (pgfplots/tikz), boîtes pédagogiques
% (définition, propriété, méthode, attention, le savais-tu, exercice résolu,
% exercice, TP, bilan), page de garde et signature OnBuch+ ancrée en bas.
%
% Macros attendues dans main.tex AVANT \input{preamble} :
%   \DOCMATIERE  \DOCNIVEAU  \DOCMODULE  \DOCLECON (numéro)  \DOCTITRE
\documentclass[11pt]{article}

\usepackage{fontspec}
\usepackage[english]{babel}
\usepackage[a4paper,margin=2cm,top=3.4cm,bottom=2.8cm,headheight=1.4cm,footskip=1.3cm]{geometry}
\usepackage{amsmath,amssymb,amsfonts}
\usepackage{mathtools} % \xmapsto, \coloneqq... souvent utilisés par les modèles
\usepackage{cancel} % simplifications barrées (bilans de forces)
% \abs{z} (module, valeur absolue) et \norm{u} (norme), utilisés par les modèles
\providecommand{\abs}{}\let\abs\relax\DeclarePairedDelimiter{\abs}{\lvert}{\rvert}
\providecommand{\norm}{}\let\norm\relax\DeclarePairedDelimiter{\norm}{\lVert}{\rVert}
\usepackage[table]{xcolor} % [table] : fournit aussi \rowcolors (alternance de lignes)
\usepackage{pagecolor}
\usepackage{tikz}
\usetikzlibrary{arrows.meta,positioning,calc,shapes.geometric,shapes.misc,shapes.arrows,decorations.pathmorphing,decorations.pathreplacing,decorations.markings,patterns,shadows,mindmap,trees,fit,backgrounds,matrix,intersections,angles,quotes}
\usepackage{pgfplots}
\usepackage[version=4]{mhchem} % équations nucléaires \ce{^{235}_{92}U}
\usepackage[european,siunitx]{circuitikz} % schémas électriques
\pgfplotsset{compat=1.18}
\usepgfplotslibrary{fillbetween,groupplots} % aires entre courbes (calcul intégral)
\usepackage{enumitem}
\usepackage{fancyhdr}
% [most] charge toutes les bibliothèques tcolorbox (listings, minted, raster,
% documentation...) et gonfle inutilement la mémoire TeX sur un document long
% (~300 boîtes) : on ne charge que ce qui est réellement utilisé.
\usepackage[skins,breakable]{tcolorbox}
\usepackage{graphicx}
\usepackage{colortbl}
\usepackage{booktabs}
\usepackage{tabularx}
\usepackage{diagbox} % case d'en-tête diagonale des tableaux à double entrée
% Colonnes extensibles centrée / gauche / droite (C, L, R), souvent utilisées par les modèles
\newcolumntype{C}{>{\centering\arraybackslash}X}
\newcolumntype{L}{>{\raggedright\arraybackslash}X}
\newcolumntype{R}{>{\raggedleft\arraybackslash}X}
\usepackage{array}
\usepackage{multicol}
\usepackage{siunitx}
% parse-numbers=false : accepte aussi \SI{2,5 \times 10^{-3}}{...}
\sisetup{locale=UK,output-decimal-marker={.},group-minimum-digits=4,parse-numbers=false}
\DeclareSIUnit\ohm{\ensuremath{\symup{\Omega}}} % U+2126 absent de la police
\DeclareSIUnit\division{div} % réglages d'oscilloscope (ms/div, V/div)
\DeclareSIUnit\tr{tr} % tours (vitesse de rotation, ex. \SI{1500}{\tr\per\minute})
\usepackage{qrcode}
% \degree : commande fréquemment utilisée par les modèles pour un angle ou une
% température en dehors de \SI{}{} (non fournie par les paquets chargés).
\providecommand{\degree}{^{\circ}}
% \qed (fin de démonstration), utilisé par les modèles sans amsthm
\providecommand{\qed}{\hfill$\square$}
% \barre : double barre des valeurs interdites dans un tableau de variations (tkz-tab)
\providecommand{\barre}{\ensuremath{\|}}
% environnement proof (amsthm non chargé) utilisé par les modèles
\ifdefined\proof\else\newenvironment{proof}[1][Proof]{\par\noindent\textit{#1.}\ }{\qed\par}\fi
\usepackage{lastpage}
\usepackage{hyperref}
\hypersetup{hidelinks}

% ============================================================
% Palette « Pop » OnBuch+ — mêmes hex que la classe P (plus_ui.dart)
% ============================================================
\definecolor{poporange}{HTML}{F59321}
\definecolor{popdark}{HTML}{E07A0C}
\definecolor{popcream}{HTML}{FFF6E8}
\definecolor{popink}{HTML}{1C1714}
\definecolor{poppurple}{HTML}{7A5AE0}
\definecolor{popgreen}{HTML}{2E9E6A}
\definecolor{poppink}{HTML}{E0457B}
\definecolor{popblue}{HTML}{2F7DE1}
\definecolor{popgold}{HTML}{C98A12}
\definecolor{popmuted}{HTML}{8A7F76}
\definecolor{popline}{HTML}{EADFD2}
\definecolor{popgray}{HTML}{8A7F76}
\colorlet{popgrey}{popgray}
% Alias tolérants (le modèle nomme parfois les couleurs différemment)
\colorlet{popwhite}{white}
\colorlet{popblack}{popink}
\colorlet{popred}{poppink}
\colorlet{popyellow}{popgold}
% Teintes claires pour fonds de tableaux / zones
\colorlet{popblueL}{popblue!10!white}
\colorlet{poporangeL}{poporange!14!white}
\colorlet{popgreenL}{popgreen!12!white}
\colorlet{poppurpleL}{poppurple!12!white}
\colorlet{poppinkL}{poppink!12!white}
\colorlet{popgoldL}{popgold!14!white}

\pagecolor{popcream}

% ============================================================
% Typographie
% ============================================================
\defaultfontfeatures{Ligatures=TeX}
\setmainfont{PlusJakartaSans-Medium.ttf}[Path=../../../fonts/,BoldFont=PlusJakartaSans-Bold.ttf]
% Maths en sans-serif (Fira Math) avec chiffres et lettres droites de Plus
% Jakarta, pour que les formules se fondent dans le texte.
\usepackage[math-style=ISO,bold-style=ISO]{unicode-math}
\setmathfont{FiraMath-Regular.otf}[Path=../../../fonts/]
\setmathfont{PlusJakartaSans-Medium.ttf}[Path=../../../fonts/,range={up/num,up/{Latin,latin},\mathup}]
% (icomma retiré : décimales avec point en anglais)
% Caractères mathématiques Unicode absents des polices de texte (Plus Jakarta,
% Archivo Black) : rendus par leur équivalent mathématique, en texte comme en
% formule — sinon ils s'affichent en carré vide (ex. « Arithmétique dans ℤ »).
\usepackage{newunicodechar}
\newunicodechar{ℝ}{\ensuremath{\mathbb{R}}}
\newunicodechar{ℤ}{\ensuremath{\mathbb{Z}}}
\newunicodechar{ℂ}{\ensuremath{\mathbb{C}}}
\newunicodechar{ℕ}{\ensuremath{\mathbb{N}}}
\newunicodechar{ℚ}{\ensuremath{\mathbb{Q}}}
\newunicodechar{𝔻}{\ensuremath{\mathbb{D}}}
\newunicodechar{α}{\ensuremath{\alpha}}
\newunicodechar{β}{\ensuremath{\beta}}
\newunicodechar{γ}{\ensuremath{\gamma}}
\newunicodechar{δ}{\ensuremath{\delta}}
\newunicodechar{ε}{\ensuremath{\varepsilon}}
\newunicodechar{θ}{\ensuremath{\theta}}
\newunicodechar{λ}{\ensuremath{\lambda}}
\newunicodechar{μ}{\ensuremath{\mu}}
\newunicodechar{σ}{\ensuremath{\sigma}}
\newunicodechar{τ}{\ensuremath{\tau}}
\newunicodechar{φ}{\ensuremath{\varphi}}
\newunicodechar{ψ}{\ensuremath{\psi}}
\newunicodechar{ω}{\ensuremath{\omega}}
\newunicodechar{Σ}{\ensuremath{\Sigma}}
% majuscules produites par \MakeUppercase dans les titres (α-aminés -> Α)
\newunicodechar{Α}{\ensuremath{\alpha}}
\newunicodechar{Β}{\ensuremath{\beta}}
\newunicodechar{Γ}{\ensuremath{\Gamma}}
\newunicodechar{Δ}{\ensuremath{\Delta}}
\newunicodechar{Ε}{\ensuremath{\varepsilon}}
\newunicodechar{Θ}{\ensuremath{\Theta}}
\newunicodechar{Λ}{\ensuremath{\Lambda}}
\newunicodechar{Μ}{\ensuremath{\mu}}
\newunicodechar{Τ}{\ensuremath{\tau}}
\newunicodechar{Φ}{\ensuremath{\Phi}}
\newunicodechar{Ψ}{\ensuremath{\Psi}}
\newunicodechar{Ω}{\ensuremath{\Omega}}
\newunicodechar{↦}{\ensuremath{\mapsto}}
\newunicodechar{∈}{\ensuremath{\in}}
\newunicodechar{∉}{\ensuremath{\notin}}
\newunicodechar{∧}{\ensuremath{\wedge}}
\newunicodechar{⇔}{\ensuremath{\Leftrightarrow}}
\newunicodechar{⟺}{\ensuremath{\Longleftrightarrow}}
\newunicodechar{⇒}{\ensuremath{\Rightarrow}}
\newunicodechar{⟹}{\ensuremath{\Longrightarrow}}
\newunicodechar{∘}{\ensuremath{\circ}}
\newunicodechar{∪}{\ensuremath{\cup}}
\newunicodechar{∩}{\ensuremath{\cap}}
\newunicodechar{≡}{\ensuremath{\equiv}}
\newunicodechar{⊂}{\ensuremath{\subset}}
\newunicodechar{⊥}{\ensuremath{\perp}}
\newunicodechar{∃}{\ensuremath{\exists}}
\newunicodechar{∀}{\ensuremath{\forall}}
\newunicodechar{∛}{\ensuremath{\sqrt[3]{\,}}}
\newunicodechar{∜}{\ensuremath{\sqrt[4]{\,}}}
\newunicodechar{⃗}{\ensuremath{{}^{\to}}}
\newunicodechar{ⁿ}{\ifmmode^{n}\else\textsuperscript{\ensuremath{n}}\fi}
\newunicodechar{ˣ}{\ifmmode^{x}\else\textsuperscript{\ensuremath{x}}\fi}
\newunicodechar{ᵇ}{\ifmmode^{b}\else\textsuperscript{\ensuremath{b}}\fi}
\newunicodechar{ʳ}{\ifmmode^{r}\else\textsuperscript{\ensuremath{r}}\fi}
\newunicodechar{ᵘ}{\ifmmode^{u}\else\textsuperscript{\ensuremath{u}}\fi}
\newunicodechar{ʸ}{\ifmmode^{y}\else\textsuperscript{\ensuremath{y}}\fi}
\newunicodechar{ᵃ}{\ifmmode^{a}\else\textsuperscript{\ensuremath{a}}\fi}
\newunicodechar{ᵉ}{\ifmmode^{e}\else\textsuperscript{\ensuremath{e}}\fi}
\newunicodechar{ᵏ}{\ifmmode^{k}\else\textsuperscript{\ensuremath{k}}\fi}
\newunicodechar{ᵖ}{\ifmmode^{p}\else\textsuperscript{\ensuremath{p}}\fi}
\newunicodechar{ᴺ}{\ifmmode^{N}\else\textsuperscript{\ensuremath{N}}\fi}
\newunicodechar{ᶜ}{\ifmmode^{c}\else\textsuperscript{\ensuremath{c}}\fi}
\newunicodechar{ʰ}{\ifmmode^{h}\else\textsuperscript{\ensuremath{h}}\fi}
\newunicodechar{ᵠ}{\ifmmode^{q}\else\textsuperscript{\ensuremath{q}}\fi}
\newunicodechar{ₙ}{\ifmmode_{n}\else\textsubscript{\ensuremath{n}}\fi}
\newunicodechar{ₐ}{\ifmmode_{a}\else\textsubscript{\ensuremath{a}}\fi}
\newunicodechar{ᵢ}{\ifmmode_{i}\else\textsubscript{\ensuremath{i}}\fi}
\newunicodechar{ᵣ}{\ifmmode_{r}\else\textsubscript{\ensuremath{r}}\fi}
\newunicodechar{ₖ}{\ifmmode_{k}\else\textsubscript{\ensuremath{k}}\fi}
\newunicodechar{ᵦ}{\ifmmode_{\beta}\else\textsubscript{\ensuremath{\beta}}\fi}
\newunicodechar{ₓ}{\ifmmode_{x}\else\textsubscript{\ensuremath{x}}\fi}
\newfontfamily\popdisplay{ArchivoBlack-Regular.ttf}[Path=../../../fonts/]
\newfontfamily\poplogo{SpaceGrotesk-Bold.ttf}[Path=../../../fonts/]
\newfontfamily\popsemi{PlusJakartaSans-SemiBold.ttf}[Path=../../../fonts/]
\newfontfamily\popxbold{PlusJakartaSans-ExtraBold.ttf}[Path=../../../fonts/]
\newfontfamily\popxboldls{PlusJakartaSans-ExtraBold.ttf}[Path=../../../fonts/,LetterSpace=3.0]

\setlist[enumerate]{leftmargin=1.7em,itemsep=5pt,topsep=4pt}
\setlist[itemize]{leftmargin=1.7em,itemsep=4pt,topsep=4pt,label={\color{poporange}\textbullet}}
\setlength{\parindent}{0pt}
\setlength{\parskip}{5pt}
\renewcommand{\arraystretch}{1.25}

% pgfplots : style Pop par défaut
\pgfplotsset{
  every axis/.append style={axis line style={popink,line width=1pt},tick style={popink},
    grid style={popline,line width=0.5pt},label style={font=\small},tick label style={font=\footnotesize}},
  popcurve/.style={poporange,line width=1.8pt,no markers,smooth},
}
% Même style disponible dans un \draw TikZ simple (hors pgfplots)
\tikzset{popcurve/.style={poporange,line width=1.8pt}}
\tikzset{popgrid/.style={popline,very thin}}
\colorlet{popcurve}{poporange} % le nom du style est parfois utilisé comme couleur

% ============================================================
% Pill / squiggle
% ============================================================
\newcommand{\poppill}[3]{%
  \tikz[baseline=-0.55ex]\node[draw=popink,line width=1.2pt,rounded corners=2.6pt,%
    fill=#2,inner xsep=6.5pt,inner ysep=3pt]{%
    \popxboldls\fontsize{7}{7}\selectfont\color{#3}\MakeUppercase{#1}};%
}
\newcommand{\popsquiggle}[1]{%
  \tikz[baseline=-0.4ex]{
    \draw[#1,line width=2.6pt,line cap=round]
      plot [smooth] coordinates {(0,0.09) (0.34,0.01) (0.68,0.21) (1.02,0.01) (1.36,0.10)};
  }%
}
% Mot-clé surligné (marqueur orange sous le texte)
\newcommand{\cle}[1]{{\popxbold\color{popdark}#1}}

% ============================================================
% En-tête / pied
% ============================================================
\pagestyle{fancy}
\fancyhf{}
\renewcommand{\headrulewidth}{0pt}
\renewcommand{\footrulewidth}{0pt}
\lhead{\raisebox{-0.15ex}{\poplogo\fontsize{13}{13}\selectfont\color{popink}OnBuch\color{poporange}+}}
\rhead{\poppill{\DOCMATIERE\ \textbullet\ \DOCNIVEAU\ \textbullet\ Chapter \DOCLECON}{white}{popink}}
\lfoot{\popxbold\fontsize{7.6}{7.6}\selectfont\color{popmuted}\MakeUppercase{Signed course OnBuch+}\ \textbullet\ {\popsemi\DOCTITRE}}
\rfoot{\tikz[baseline=-0.6ex]\node[rounded corners=8.5pt,fill=popink,minimum height=17pt,minimum width=17pt,inner xsep=5pt,inner sep=0]{\popxbold\fontsize{8}{8}\selectfont\color{popcream}\thepage\,/\,\pageref*{LastPage}};}

% ============================================================
% Titres
% ============================================================
\newcommand{\chaptitre}[2]{%
  \begin{tcolorbox}[enhanced,colback=poporange,colframe=popink,boxrule=2.6pt,arc=11pt,
      drop shadow=popink,left=15pt,right=15pt,top=12pt,bottom=11pt,before skip=3pt,after skip=18pt]
    \poppill{#2}{popink}{popcream}\par\vspace{9pt}
    {\raggedright\hyphenpenalty=10000\exhyphenpenalty=10000\popdisplay\fontsize{22}{24}\selectfont\color{popcream}\MakeUppercase{#1}\par}
  \end{tcolorbox}
}
% Section numérotée : pastille orange avec le numéro + titre Archivo
\newcounter{popsec}
\newcommand{\coursec}[1]{%
  \stepcounter{popsec}\setcounter{popsub}{0}%
  \par\vspace{16pt}\noindent
  \begin{minipage}{\linewidth}
  \tikz[baseline=-0.7ex]\node[draw=popink,line width=1.6pt,fill=poporange,rounded corners=4pt,minimum size=22pt,inner sep=0]{\popdisplay\fontsize{12}{12}\selectfont\color{popcream}\thepopsec};%
  \hspace{8pt}{\popdisplay\fontsize{15}{17}\selectfont\color{popink}\MakeUppercase{#1}}\par
  \vspace{4pt}\noindent{\color{poporange}\rule{36pt}{3.6pt}}
  \end{minipage}\par\nopagebreak\vspace{8pt}
}
\newcounter{popsub}
\newcommand{\courssub}[1]{%
  \stepcounter{popsub}%
  \par\vspace{9pt}\noindent{\popxbold\fontsize{11.5}{13}\selectfont\color{popink}{\color{poporange}\thepopsec.\thepopsub}\hspace{6pt}#1}\par\nopagebreak\vspace{4pt}
}

% ============================================================
% Boîtes pédagogiques — même recette : fond blanc, bordure épaisse colorée,
% ombre dure encre, badge plein. Titre optionnel [#1].
% ============================================================
\tcbset{popbase/.style={breakable,enhanced,colback=white,boxrule=2.3pt,arc=9pt,
  shadow={3pt}{-3pt}{0pt}{popink},left=14pt,right=14pt,top=11pt,bottom=11pt,before skip=11pt,after skip=15pt,
  fontupper=\popsemi\fontsize{10.6}{14.5}\selectfont\color{popink},
  fontlower=\popsemi\fontsize{10.6}{14.5}\selectfont\color{popink}}}
% Coche dessinée (le glyphe ✓ n'existe pas dans les polices du document)
\newcommand{\popcheck}[1]{\tikz[baseline=-0.5ex]\draw[#1,line width=1.6pt,line cap=round,line join=round](-0.13,0.0)--(-0.04,-0.09)--(0.13,0.1);}
\providecommand{\coche}{\popcheck{popgreen}}
\providecommand{\resultat}[1]{\par\smallskip\noindent\fcolorbox{popgreen}{popgreenL}{\parbox{\dimexpr\linewidth-2\fboxsep-2\fboxrule\relax}{\textbf{Result.} #1}}\par\smallskip}
\newcommand{\popboxhead}[3]{\poppill{#1}{#2}{white}\ifx\relax#3\relax\else\hspace{6pt}{\popxbold\fontsize{10.6}{12}\selectfont\color{popink}#3}\fi\par\vspace{7pt}}

\newtcolorbox{aretenir}[1][]{popbase,colframe=popgreen,before upper={\popboxhead{Key points}{popgreen}{#1}}}
\newtcolorbox{exemplebox}[1][]{popbase,colframe=poporange,before upper={\popboxhead{Exemple}{poporange}{#1}}}
\newtcolorbox{definition}[1][]{popbase,colframe=popblue,before upper={\popboxhead{Definition}{popblue}{#1}}}
\newtcolorbox{propriete}[1][]{popbase,colframe=poppurple,before upper={\popboxhead{Property}{poppurple}{#1}}}
\newtcolorbox{methode}[1][]{popbase,colframe=popgold,before upper={\popboxhead{Method}{popgold}{#1}}}
\newtcolorbox{attention}[1][]{popbase,colframe=poppink,before upper={\popboxhead{Warning}{poppink}{#1}}}
\newtcolorbox{experience}[1][]{popbase,colframe=popblue,colback=popblueL,before upper={\popboxhead{Experiment}{popblue}{#1}}}
% « Le savais-tu ? » : carte violette pleine (contexte, culture, Cameroun)
\newtcolorbox{savaistu}[1][]{popbase,colback=poppurple,colframe=popink,
  fontupper=\popsemi\fontsize{10.4}{14.2}\selectfont\color{white},
  before upper={\poppill{Did you know?}{white}{poppurple}\ifx\relax#1\relax\else\hspace{6pt}{\popxbold\color{white}#1}\fi\par\vspace{7pt}}}
% Exercice résolu : énoncé en haut, \tcblower puis la solution sur fond crème
\newcounter{popexo}\newcounter{popexr}
\newtcolorbox{exoresolu}[1][]{popbase,colframe=poporange,colbacklower=popcream,
  segmentation style={popink,line width=1.2pt,dashed},
  before upper={\stepcounter{popexr}\popboxhead{Worked example \thepopexr}{poporange}{#1}},
  before lower={\poppill{Solution}{popink}{popcream}\par\vspace{6pt}}}
% Exercice (non résolu en place) — la correction est dans la partie Corrigés
\newtcolorbox{exercice}[1][]{popbase,colframe=popink,
  before upper={\stepcounter{popexo}\popboxhead{Exercise \thepopexo}{popink}{#1}}}
% Corrigé
\newtcolorbox{corrige}[1][]{popbase,colframe=popgreen,colback=popgreenL,
  before upper={\popboxhead{Solution}{popgreen}{#1}}}
% Objectifs / prérequis / bilan
\newtcolorbox{objectifs}{popbase,colframe=popink,colback=poporangeL,
  before upper={\popboxhead{Chapter objectives}{popink}{}}}
\newtcolorbox{prerequis}{popbase,colframe=popink,colback=popblueL,
  before upper={\popboxhead{Prerequisites}{popblue}{}}}
\newtcolorbox{bilan}{popbase,colframe=popink,colback=white,boxrule=2.8pt,
  before upper={\popboxhead{Fiche bilan}{poporange}{L'essentiel à connaître pour l'examen}}}
% Figure encadrée avec légende
\newtcolorbox{popfigure}[1][]{enhanced,breakable=false,colback=white,colframe=popline,boxrule=1.4pt,arc=8pt,
  left=10pt,right=10pt,top=10pt,bottom=8pt,before skip=10pt,after skip=12pt,halign=center,
  lower separated=false,halign lower=center,
  fontlower=\popsemi\fontsize{9}{11.5}\selectfont\color{popmuted},
  #1}
\newcommand{\legende}[1]{\par\vspace{6pt}{\popsemi\fontsize{9}{11.5}\selectfont\color{popmuted}#1}}

% ============================================================
% Signature OnBuch+
% ============================================================
\newcommand{\poplogoinline}[1]{{\poplogo\fontsize{#1}{#1}\selectfont\color{popink}OnBuch\color{poporange}+}}
\newcommand{\signatureonbuch}{%
  \begin{tcolorbox}[enhanced,colback=popink,colframe=popink,boxrule=2.2pt,arc=10pt,
      drop shadow=poporange,left=14pt,right=14pt,top=10pt,bottom=10pt,before skip=0pt,after skip=0pt]
    \tikz[baseline=-0.7ex]\node[circle,fill=popgreen,draw=popcream,line width=1.3pt,minimum size=22pt,inner sep=0]{\popcheck{white}};%
    \hspace{9pt}\begin{minipage}[c]{0.62\linewidth}
      {\popxbold\fontsize{12}{13}\selectfont\color{popcream}Signed\ \textbullet\ {\poplogo OnBuch\color{poporange}+}}\par\vspace{2pt}
      {\raggedright\popsemi\fontsize{8}{10.5}\selectfont\color{popline}\MakeUppercase{OnBuch+ teaching team}\\\MakeUppercase{Follows the official MINESEC syllabus}\par}
    \end{minipage}\hfill
    {\poplogo\fontsize{22}{22}\selectfont\color{popcream}OnBuch\color{poporange}+}
  \end{tcolorbox}%
}

% ============================================================
% Page de garde
% #1 titre · #2 matière · #3 niveau · #4 module · #5 n° leçon · #6 durée ·
% #7 description / objectifs (texte ou itemize)
% ============================================================
\newcommand{\pagegarde}[7]{%
  \thispagestyle{empty}
  \newgeometry{a4paper,margin=1.7cm,top=1.6cm,bottom=1.4cm}%
  \begin{tikzpicture}[remember picture,overlay]
    \fill[poporange!18] ([xshift=-3.2cm,yshift=-3.6cm]current page.north east) circle (4.6cm);
    \fill[poppurple!14] ([xshift=2.4cm,yshift=7.6cm]current page.south west) circle (3.8cm);
  \end{tikzpicture}%
  {\poplogo\fontsize{30}{30}\selectfont\color{popink}OnBuch\color{poporange}+}\hfill
  \raisebox{0.6ex}{\poppill{Signed course \textbullet\ Premium}{popink}{popcream}}\par
  \vspace{1pt}\popsquiggle{poppurple}\par
  \vspace{8pt}
  \begin{tcolorbox}[enhanced,colback=poporange,colframe=popink,boxrule=2.8pt,arc=13pt,
      drop shadow=popink,left=18pt,right=18pt,top=12pt,bottom=13pt,after skip=10pt]
    \poppill{#2 \textbullet\ #3}{popink}{popcream}\hspace{5pt}\poppill{#4}{white}{popink}\par\vspace{8pt}
    {\popdisplay\fontsize{14}{14}\selectfont\color{popink}CHAPTER #5}\par\vspace{4pt}
    {\raggedright\hyphenpenalty=10000\exhyphenpenalty=10000\popdisplay\fontsize{25}{27}\selectfont\color{popcream}\MakeUppercase{#1}\par}
    \vspace{8pt}
    {\popxbold\fontsize{9.5}{9.5}\selectfont\color{popink}\MakeUppercase{Suggested time: #6}}
  \end{tcolorbox}
  \begin{tcolorbox}[enhanced,colback=white,colframe=popink,boxrule=2.2pt,arc=10pt,
      drop shadow=popink,left=16pt,right=16pt,top=10pt,bottom=10pt,after skip=8pt]
    \poppill{In this chapter}{poppurple}{white}\par\vspace{5pt}
    \popsemi\fontsize{9.3}{12.6}\selectfont\color{popink}#7
  \end{tcolorbox}
  \noindent\begin{minipage}[t]{0.487\linewidth}
    \begin{tcolorbox}[enhanced,colback=popgreen,colframe=popink,boxrule=2.2pt,arc=10pt,
        drop shadow=popink,left=11pt,right=11pt,top=10pt,bottom=10pt,height=3.1cm,valign=center]
      \tcbox[colback=white,colframe=popink,boxrule=1.3pt,arc=5pt,boxsep=0pt,nobeforeafter,
        left=3pt,right=3pt,top=3pt,bottom=3pt,valign=center]{\qrcode[height=1.9cm]{https://play.google.com/store/apps/details?id=com.luvvix.onbuch&hl=en}}%
      \hspace{8pt}\begin{minipage}[c]{0.5\linewidth}
        {\popdisplay\fontsize{11}{12.5}\selectfont\color{white}\MakeUppercase{Download OnBuch}}\par\vspace{4pt}
        {\raggedright\popsemi\fontsize{8.4}{11}\selectfont\color{white}Free \textbullet\ Google Play\\AI tutor, past papers, results\par}
      \end{minipage}
    \end{tcolorbox}
  \end{minipage}\hfill
  \begin{minipage}[t]{0.487\linewidth}
    \begin{tcolorbox}[enhanced,colback=poppurple,colframe=popink,boxrule=2.2pt,arc=10pt,
        drop shadow=popink,left=11pt,right=11pt,top=10pt,bottom=10pt,height=3.1cm,valign=center]
      \tikz[baseline=-0.7ex]\node[circle,fill=white,draw=popink,line width=1.3pt,minimum size=1.6cm,inner sep=0]{\popdisplay\fontsize{24}{24}\selectfont\color{poppurple}+};%
      \hspace{8pt}\begin{minipage}[c]{0.6\linewidth}
        {\popdisplay\fontsize{11}{12.5}\selectfont\color{white}\MakeUppercase{Go OnBuch+}}\par\vspace{4pt}
        {\raggedright\popsemi\fontsize{8.4}{11}\selectfont\color{white}All signed courses, unlimited AI tutor, PDF sheets — OnBuch+ tab in the app\par}
      \end{minipage}
    \end{tcolorbox}
  \end{minipage}
  \vspace{8pt}
  \popcontenu
  \vfill
  \signatureonbuch
  \restoregeometry
  \newpage
}

% Rangée « Ce cours contient » (page de garde)
\newcommand{\poptile}[3]{% #1 couleur #2 gros texte #3 légende
  \begin{tcolorbox}[enhanced,colback=#1,colframe=popink,boxrule=2pt,arc=9pt,shadow={2.5pt}{-2.5pt}{0pt}{popink},
      left=6pt,right=6pt,top=9pt,bottom=9pt,halign=center,valign=center,height=1.75cm,width=\linewidth,nobeforeafter]
    {\popdisplay\fontsize{12.5}{13}\selectfont\color{white}\MakeUppercase{#2}}\par\vspace{3pt}
    {\centering\popsemi\fontsize{7.8}{9.5}\selectfont\color{white}#3\par}
  \end{tcolorbox}}
\newcommand{\popcontenu}{%
  \par\noindent{\popxboldls\fontsize{7.5}{7.5}\selectfont\color{popmuted}THIS SIGNED COURSE CONTAINS}\par\vspace{7pt}
  \noindent\begin{minipage}[t]{0.235\linewidth}\poptile{popblue}{Course}{complete, illustrated, step by step}\end{minipage}\hfill
  \begin{minipage}[t]{0.235\linewidth}\poptile{poporange}{Methods}{and worked examples}\end{minipage}\hfill
  \begin{minipage}[t]{0.235\linewidth}\poptile{poppink}{Exercises}{exam style + solutions}\end{minipage}\hfill
  \begin{minipage}[t]{0.235\linewidth}\poptile{popgreen}{Summary sheet}{the essentials to remember}\end{minipage}\par}

% Sommaire
\newtcolorbox{sommaire}{enhanced,colback=white,colframe=popink,boxrule=2.2pt,arc=10pt,
  drop shadow=popink,left=18pt,right=18pt,top=13pt,bottom=13pt,before skip=6pt,after skip=20pt,
  before upper={\poppill{Contents}{poppurple}{white}\par\vspace{9pt}\popsemi\fontsize{10.6}{14.5}\selectfont\color{popink}}}

% Signature de fin de cours — poussée tout en bas de la dernière page
\newcommand{\findecours}{%
  \par\vfill\nopagebreak
  \begin{center}\popsquiggle{poporange}\end{center}\vspace{4pt}
  \signatureonbuch
}
\AtBeginDocument{\def\llbracket{\mathopen{[\mkern-3.5mu[}}\def\rrbracket{\mathclose{]\mkern-3.5mu]}}} % intervalles d'entiers (glyphe ⟦ absent de la police)
% Glyphes absents de Fira Math : redéfinis à partir de glyphes disponibles
\AtBeginDocument{%
  \def\setminus{\mathbin{\backslash}}\let\smallsetminus\setminus
  \def\vdots{\mathord{\vbox{\baselineskip3.2pt\lineskiplimit0pt\kern4pt\hbox{.}\hbox{.}\hbox{.}}}}%
  \def\ddots{\mathinner{\mkern1mu\raise6.5pt\hbox{.}\mkern2mu\raise3.6pt\hbox{.}\mkern2mu\raise0.7pt\hbox{.}\mkern1mu}}%
  \def\triangle{\mathord{\scalebox{0.9}{$\symup{\Delta}$}}}%
  \def\nabla{\mathord{\raisebox{\depth}{\scalebox{1}[-1]{$\symup{\Delta}$}}}}%
  \def\checkmark{\popcheck{popdark}}%
  \def\star{\mathbin{\ast}}\def\bigstar{\mathord{\ast}}%
  \def\diamond{\mathbin{\scalebox{0.6}{$\Diamond$}}}%
  \def\bigcup{\mathop{\mathchoice{\raisebox{-0.3em}{\scalebox{1.7}{$\cup$}}}{\scalebox{1.25}{$\cup$}}{\cup}{\cup}}\displaylimits}%
  \def\bigcap{\mathop{\mathchoice{\raisebox{-0.3em}{\scalebox{1.7}{$\cap$}}}{\scalebox{1.25}{$\cap$}}{\cap}{\cap}}\displaylimits}%
  \def\lesseqqgtr{\mathrel{\lessgtr}}\def\bigoplus{\mathop{\oplus}\displaylimits}%
}
\providecommand{\ssi}{if and only if} % abréviation fréquente
\AtBeginDocument{\providecommand{\wideparen}{\overparen}} % arc de cercle

%% FIGFIX-BEGIN v5 — correctifs globaux des figures (docs/onbuch-generation/tools/figfix)
\usepackage{adjustbox}\usepackage{etoolbox}\tcbuselibrary{raster}
% toute figure TikZ/pgfplots plus large que la ligne est réduite à la largeur disponible
\BeforeBeginEnvironment{tikzpicture}{\begin{adjustbox}{max width=\linewidth}}
\AfterEndEnvironment{tikzpicture}{\end{adjustbox}}
% légendes pgfplots lisibles par-dessus les courbes
\pgfplotsset{every axis/.append style={legend style={fill=white,fill opacity=0.92,text opacity=1,draw=black!15,inner sep=3pt}}}
% étiquettes d'arêtes (syntaxe edge["..."]) avec fond blanc
\tikzset{every edge quotes/.append style={fill=white,inner sep=1.2pt}}
%% FIGFIX-END
\begin{document}

\pagegarde{Centre of Mass of Uniform Rigid Bodies}{Pure Mathematics with Mechanics}{Upper Sixth}{Upper Sixth — Pure mathematics with mechanics}{13}{18 periods}{This chapter develops the concept of centre of mass, from systems of particles to uniform rigid bodies, laminae and composite figures. Students learn to locate centres of mass using symmetry, coordinate formulae, position vectors and calculus, and apply these ideas to equilibrium, toppling and real engineering situations.
\vspace{4pt}
\begin{multicols}{2}\small
\begin{itemize}
\item By the end of this chapter I can explain the concept of centre of mass and its role in the motion of a body.
\item By the end of this chapter I can find the centre of mass of a system of particles in one, two and three dimensions using coordinates.
\item By the end of this chapter I can use position vectors to locate the centre of mass of a system of particles.
\item By the end of this chapter I can use symmetry arguments to locate the centre of mass of uniform rods, laminae and standard solids.
\item By the end of this chapter I can find the centre of mass of composite laminae by the addition and subtraction (negative mass) methods.
\item By the end of this chapter I can determine the centroid of a region bounded by curves using integration.
\end{itemize}\end{multicols}}

\begin{sommaire}
\begin{multicols}{2}
\begin{enumerate}
\item Concept of Centre of Mass and Motion of the Mass Centre
\item Centre of Mass of a Uniform Rigid Body
\item Centre of Mass Using Coordinates
\item Position Vectors and the Centre of Mass Formula
\item Uniform Laminae, Symmetry and Composite Laminae
\item Applications of Centre of Mass of Laminae and Rigid Bodies
\item Centroids of Regions Bounded by Curves
\item Centroids of Composite Areas and Applications of Centroids
\item Integration activity
\item Methods and worked examples
\item Exercises
\item Solutions to the exercises
\item Summary sheet
\end{enumerate}
\end{multicols}
\end{sommaire}

\begin{objectifs}
\begin{itemize}
\item By the end of this chapter I can explain the concept of centre of mass and its role in the motion of a body.
\item By the end of this chapter I can find the centre of mass of a system of particles in one, two and three dimensions using coordinates.
\item By the end of this chapter I can use position vectors to locate the centre of mass of a system of particles.
\item By the end of this chapter I can use symmetry arguments to locate the centre of mass of uniform rods, laminae and standard solids.
\item By the end of this chapter I can find the centre of mass of composite laminae by the addition and subtraction (negative mass) methods.
\item By the end of this chapter I can determine the centroid of a region bounded by curves using integration.
\item By the end of this chapter I can find the centroid of composite areas by combining standard results.
\item By the end of this chapter I can apply centre of mass to problems of equilibrium, suspension, balancing and toppling of rigid bodies.
\end{itemize}
\end{objectifs}

\begin{prerequis}
\begin{itemize}
\item Vectors in two and three dimensions, including position vectors (Lower Sixth, Further Pure/PM).
\item Moments of forces and conditions for equilibrium of a rigid body (Lower Sixth Mechanics).
\item Definite integration, including areas under and between curves (Pure Mathematics).
\item Coordinate geometry: midpoints, distance formula, equations of lines and curves.
\item Density, mass and uniform density of a body (basic physics/earlier mechanics).
\end{itemize}
\end{prerequis}

\newpage

\coursec{Concept of Centre of Mass and Motion of the Mass Centre}

At the Mokolo market in Yaoundé, a trader balances a huge basket of plantains on her head and walks steadily through the crowd — she has, without any formula, placed the basket exactly above its \cle{centre of mass}. On the Douala–Bafoussam road, an overloaded lorry topples over on a bend because its load was stacked too high and too far to one side. Carpenters in Buea cut table tops so that they never wobble, and engineers check that a water tank on a tower will not tip in the wind. All these situations obey the same mathematical idea: every rigid body behaves, for balance and motion, as if its whole mass were concentrated at one single point. This chapter gives you the tools to locate that point for particles, rods, laminae and solids — by symmetry, by coordinates, by position vectors and by integration.

\courssub{The Balance Point: an Intuitive Introduction}

Take a uniform wooden ruler and lay it flat on the edge of your finger. If your finger is near one end, the ruler tips and falls. Slide your finger towards the middle: at one special position, the ruler stays perfectly horizontal. That position is the \emph{balance point} of the ruler. Now imagine a plank with a heavy toolbox resting near one end. To balance the plank on a single support, you must move the support towards the toolbox — the balance point has shifted towards the heavier side. Finally, think of a loaded wheelbarrow: the gardener lifts the handles and the whole load pivots about the wheel; the load feels ``light'' or ``heavy'' depending on where the mass is concentrated.

In each case there is a single point with a remarkable property: if the body is supported exactly at (or below) that point, gravity produces no turning effect, and the body balances. That point is called the \cle{centre of mass} of the body.

\begin{experience}[Balancing a ruler and a loaded plank]
Place a uniform metre rule on a horizontal knife edge. It balances at the \SI{50}{cm} mark. Now tape a lump of clay of mass comparable to the rule near the \SI{80}{cm} mark: the balance point moves towards the clay, to about the \SI{60}{cm} mark. The balance point always shifts \emph{towards the added mass}, and the shift is larger when the added mass is larger or placed further away. This simple experiment contains the whole idea of the mass-weighted average.
\end{experience}

\begin{popfigure}
\begin{tikzpicture}[scale=1.0, >=stealth]
% pivot
\fill[poppurple] (0,0) -- (-0.45,-0.8) -- (0.45,-0.8) -- cycle;
\draw[popdark, thick] (0,0) -- (-0.45,-0.8) -- (0.45,-0.8) -- cycle;
\draw[popdark] (-1.4,-0.8) -- (1.4,-0.8);
% plank
\draw[fill=popblueL, draw=popdark, thick] (-3.2,0) rectangle (3.2,0.35);
% centre of mass
\fill[poppink] (0,0.175) circle (0.09);
\node[popdark, above] at (0,0.42) {\small $G$};
% weight arrow
\draw[->, poppink, very thick] (0,0.175) -- (0,-0.75) node[right, popdark, pos=0.7]{\small $M\vec{g}$};
% reaction
\draw[->, popgreen, very thick] (0,-0.05) -- (0,1.05) node[right, popdark, pos=0.75]{\small $\vec{R}$};
% uniform label
\node[popmuted] at (-2.2,0.18) {\small uniform plank};
% dimension
\draw[<->, popmuted] (-3.2,-1.15) -- (3.2,-1.15);
\node[popmuted, below] at (0,-1.15) {\small length $L$};
\end{tikzpicture}
\legende{A uniform plank balanced on a pivot. The centre of mass $G$ is at the midpoint; the weight $M\vec{g}$ acts through $G$ and is exactly opposed by the reaction $\vec{R}$ of the pivot.}
\end{popfigure}

\courssub{Definition of the Centre of Mass of a System of Particles}

The experiment above suggests that the balance point is an \emph{average} position in which heavy particles ``count more'' than light ones. This is made precise as follows.

\begin{definition}[Centre of mass of a system of particles]
Consider a system of $n$ particles of masses $m_1, m_2, \dots, m_n$ situated at points with coordinates $(x_1,y_1,z_1), (x_2,y_2,z_2), \dots, (x_n,y_n,z_n)$ in a frame $(O;\vec{i},\vec{j},\vec{k})$. The \cle{centre of mass} (or \emph{mass centre}) $G$ of the system is the point whose coordinates are
\[
\bar{x}=\frac{\sum_{i=1}^{n} m_i x_i}{\sum_{i=1}^{n} m_i},\qquad
\bar{y}=\frac{\sum_{i=1}^{n} m_i y_i}{\sum_{i=1}^{n} m_i},\qquad
\bar{z}=\frac{\sum_{i=1}^{n} m_i z_i}{\sum_{i=1}^{n} m_i}.
\]
Writing $M=\sum_{i=1}^{n} m_i$ for the total mass, this is the \textbf{mass-weighted average} of the positions:
\[
M\bar{x}=\sum_{i=1}^{n} m_i x_i,\qquad M\bar{y}=\sum_{i=1}^{n} m_i y_i,\qquad M\bar{z}=\sum_{i=1}^{n} m_i z_i .
\]
\end{definition}

For two particles on a line, the formula reduces to a very intuitive statement: $G$ divides the segment joining the particles in the \emph{inverse ratio} of the masses. Indeed, if $m_1$ is at $x_1$ and $m_2$ at $x_2$, then
\[
\bar{x}=\frac{m_1x_1+m_2x_2}{m_1+m_2}
\quad\Longrightarrow\quad
m_1(\bar{x}-x_1)=m_2(x_2-\bar{x}),
\]
so the heavier particle is the \emph{closer} one to $G$ — exactly what the loaded plank experiment showed.

\begin{exemplebox}[Two masses on a rod]
A light rod $AB$ of length \SI{1.2}{m} carries a mass of \SI{3}{kg} at $A$ and a mass of \SI{5}{kg} at $B$. Taking $A$ as origin and the $x$-axis along the rod,
\[
\bar{x}=\frac{3(0)+5(1.2)}{3+5}=\frac{6}{8}=0.75\ \mathrm{m}.
\]
The centre of mass lies \SI{0.75}{m} from $A$, nearer to the heavier mass at $B$. The rod balances on a knife edge placed at that point.
\end{exemplebox}

\courssub{Centre of Mass and Centre of Gravity}

In everyday language the two terms are used interchangeably, but a careful student must distinguish them.

\begin{definition}[Centre of gravity]
The \cle{centre of gravity} of a body is the point through which the \emph{resultant of the weights} of all its particles acts, whatever the orientation of the body. Each particle of mass $m_i$ experiences a weight $m_i\vec{g}_i$; the centre of gravity is the point about which the total moment of these weights is zero.
\end{definition}

If the gravitational field is \emph{uniform} — the same vector $\vec{g}$ at every particle — then the total moment about $G$ (the centre of mass) is
\[
\sum_{i} m_i\,\vec{GA}_i\wedge\vec{g}=\left(\sum_{i}m_i\,\vec{GA}_i\right)\wedge\vec{g}=\vec{0}\wedge\vec{g}=\vec{0},
\]
because $\sum_i m_i\vec{GA}_i=M\vec{GG}=\vec{0}$ by the very definition of $G$. Hence:

\begin{propriete}[Coincidence of centre of mass and centre of gravity]
In a \textbf{uniform gravitational field}, the centre of gravity coincides with the centre of mass. For every body of ordinary size on or near the Earth's surface, the two points are the same, and we may treat the weight $M\vec{g}$ as a single force acting at $G$. They would differ only for a body so large that $\vec{g}$ varies appreciably across it (for example a very tall tower or an orbiting station).
\end{propriete}

\begin{attention}[The centre of mass need not lie inside the material]
The centre of mass is a \emph{mathematical} point; it does not have to be inside the body. The centre of mass of a uniform ring is at its geometric centre — where there is no metal at all. The centre of mass of a boomerang lies in the empty space between its arms; that of a cup lies somewhere in the hollow. Never reject an answer merely because $G$ ``is not on the object''.
\end{attention}

\courssub{Concentration of the Mass at the Centre of Mass}

Why is this point so important? Because for \emph{translational} questions, the whole body can be replaced by a single particle of mass $M$ placed at $G$.

\begin{propriete}[Motion of the mass centre — proof examinable]
Let a system of particles of total mass $M$ be subject to external forces $\vec{F}_1^{\,\mathrm{ext}},\vec{F}_2^{\,\mathrm{ext}},\dots$ and internal forces between the particles. Then the centre of mass $G$ moves as a single particle of mass $M$ under the resultant of the \textbf{external} forces only:
\[
M\,\ddot{\vec{r}}_G=\sum_{i}\vec{F}_i^{\,\mathrm{ext}}.
\]
\emph{Proof.} For each particle, Newton's second law gives $m_i\ddot{\vec{r}}_i=\vec{F}_i^{\,\mathrm{ext}}+\vec{F}_i^{\,\mathrm{int}}$. Summing over all particles,
\[
\sum_i m_i\ddot{\vec{r}}_i=\sum_i\vec{F}_i^{\,\mathrm{ext}}+\sum_i\vec{F}_i^{\,\mathrm{int}}.
\]
By Newton's third law, internal forces occur in equal and opposite pairs, so $\sum_i\vec{F}_i^{\,\mathrm{int}}=\vec{0}$. Differentiating $M\vec{r}_G=\sum_i m_i\vec{r}_i$ twice gives $M\ddot{\vec{r}}_G=\sum_i m_i\ddot{\vec{r}}_i$, and the result follows. $\square$
\end{propriete}

\begin{aretenir}[Key consequences]
\begin{itemize}
\item The \textbf{whole mass} of a body may be considered concentrated at $G$ for translational motion.
\item \textbf{Internal forces} (explosions, pushes between parts, muscles) cannot alter the motion of $G$; only external forces can.
\item If the resultant external force is zero, $G$ moves with \textbf{constant velocity} — in particular, if $G$ starts at rest, it stays at rest even while the parts move.
\end{itemize}
\end{aretenir}

\begin{exemplebox}[A thrown spanner]
A spanner is thrown across a workshop. It spins and tumbles, and each point of it follows a complicated path. But the only external force (neglecting air resistance) is its weight $M\vec{g}$, so
\[
M\ddot{\vec{r}}_G=M\vec{g}\quad\Longrightarrow\quad \ddot{\vec{r}}_G=\vec{g}.
\]
The centre of mass $G$ — and $G$ alone — follows a perfect \textbf{parabola}, exactly as a single projectile of mass $M$ would. The rotation happens \emph{about} $G$ and does not disturb the path of $G$.
\end{exemplebox}

\begin{popfigure}
\begin{tikzpicture}[scale=1.05, >=stealth]
% parabolic path of G
\draw[thick, poppink, domain=0:6, samples=60] plot (\x, {2.2*(\x)*(1-\x/6)*1.35});
\node[popdark] at (5.6,2.6) {\small path of $G$};
% spanner at three positions (rotating)
\begin{scope}[shift={(0.6,0.62)}, rotate=15]
\draw[line width=2.2pt, popblue] (-0.55,0) -- (0.55,0);
\draw[popblue, thick] (-0.55,0) circle (0.14);
\draw[popblue, thick] (0.55,0) circle (0.14);
\end{scope}
\begin{scope}[shift={(3,2.23)}, rotate=100]
\draw[line width=2.2pt, popblue] (-0.55,0) -- (0.55,0);
\draw[popblue, thick] (-0.55,0) circle (0.14);
\draw[popblue, thick] (0.55,0) circle (0.14);
\end{scope}
\begin{scope}[shift={(5.4,0.62)}, rotate=185]
\draw[line width=2.2pt, popblue] (-0.55,0) -- (0.55,0);
\draw[popblue, thick] (-0.55,0) circle (0.14);
\draw[popblue, thick] (0.55,0) circle (0.14);
\end{scope}
% mark G at each position
\fill[poppink] (0.6,0.62) circle (0.06);
\fill[poppink] (3,2.23) circle (0.06);
\fill[poppink] (5.4,0.62) circle (0.06);
\node[popdark, above right] at (3,2.3) {\small $G$};
% ground
\draw[popdark] (-0.5,0) -- (6.6,0);
\fill[poppurple!25] (-0.5,0) rectangle (6.6,-0.18);
\end{tikzpicture}
\legende{A spanner thrown through the air rotates about its centre of mass $G$, but $G$ itself describes a parabola (pink curve), exactly like a point projectile of the same total mass.}
\end{popfigure}

\begin{exoresolu}[Two connected particles on a smooth table]
Two particles $P$ and $Q$, of masses \SI{2}{kg} and \SI{4}{kg} respectively, rest on a smooth horizontal table, joined by a light inextensible string of length \SI{3}{m}. Initially $P$ is at the origin and $Q$ at $(3,0)$. An external agent pulls the string so that $P$ moves to $(0,2)$ and $Q$ moves to a point on the $x$-axis such that the string remains taut. Find the final position of the centre of mass and explain why it moves. Then consider the case where the particles move only under their mutual attraction (internal forces) and show that $G$ remains fixed.
\tcblower
\textbf{Initial position of $G$:}
\[
\bar{x}=\frac{2(0)+4(3)}{2+4}=\frac{12}{6}=2,\qquad \bar{y}=\frac{2(0)+4(0)}{6}=0,
\]
so $G=(2,0)$ initially.

\textbf{Final position with external pull:} The string length is \SI{3}{m}, so if $P$ is at $(0,2)$, $Q$ must be at $(x,0)$ with $\sqrt{x^2+2^2}=3$, giving $x=\sqrt{5}$ (positive since $Q$ is to the right). Then
\[
\bar{x}=\frac{2(0)+4\sqrt{5}}{6}=\frac{2\sqrt{5}}{3}\approx 1.49,\qquad \bar{y}=\frac{2(2)+4(0)}{6}=\frac{2}{3}\approx 0.67.
\]
The centre of mass \emph{has moved} from $(2,0)$ to $(\frac{2\sqrt{5}}{3},\frac{2}{3})$. This is correct: the pull was applied by an \textbf{external} agent, so the resultant external force is not zero and $G$ accelerates.

\textbf{Motion under internal forces only:} Suppose instead the particles are released from rest with the string taut, and they move \emph{only} under their mutual tension (internal forces). There is no external horizontal force on the system. Then the horizontal position of $G$ must remain constant:
\[
2\,x_P+4\,x_Q = 6\,\bar{x} = \text{constant} = 12.
\]
As $P$ and $Q$ move towards each other, $G$ stays fixed at $x=2$. This illustrates the principle: \emph{internal forces alone never displace the centre of mass}.
\end{exoresolu}

\courssub{Everyday Applications of the Mass Centre}

\textbf{The high-jumper.} A skilled high-jumper arches her body over the bar so that her \emph{body} passes over the bar while her \emph{centre of mass} actually passes \emph{under} it. Since the energy needed depends on how high $G$ must be lifted, this technique allows a higher clearance for the same effort — a beautiful use of the fact that $G$ need not lie inside the body.

\textbf{Stability of a loaded lorry.} A vehicle topples when the vertical line through its centre of mass falls outside the base of support (the polygon joining the contact points of the wheels). Stacking a load high and to one side raises and shifts $G$; on a bend, the vertical through $G$ can then pass outside the wheels and the lorry overturns. This is why loads must be kept \emph{low} and \emph{centred}.

\textbf{The trader's basket.} The basket is stable on the trader's head only when the vertical through $G$ passes through the support area. She adjusts the load instinctively until this is so — she has located the centre of mass by feel.

\begin{savaistu}[Did you know?]
The same principle governs the ``Leaning Tower'' problem: a uniform stack of blocks on a table edge remains stable as long as the vertical through the centre of mass of the whole stack passes through the supporting block below. With infinitely many blocks, the top block can overhang the table edge by an arbitrarily large distance — the overhangs grow like the harmonic series $1+\tfrac12+\tfrac13+\cdots$, which diverges!
\end{savaistu}

\begin{attention}[Common mistakes]
\begin{itemize}
\item Confusing the \emph{geometric centre} with the centre of mass when the body is \textbf{not} uniform: symmetry arguments only work for uniform bodies (see later sections).
\item Forgetting that the formula uses \textbf{masses}, not weights, volumes or numbers of particles — although in a uniform field weighting by weights gives the same point.
\item Believing internal forces (an explosion, a push between two parts) can move $G$: they cannot; only the \textbf{resultant external force} governs $\ddot{\vec{r}}_G$.
\item Assuming $G$ must be inside the material of the body: false for rings, cups, boomerangs and arched jumpers.
\end{itemize}
\end{attention}

\begin{exercice}[Mass centre of three collinear particles]
Three particles of masses \SI{1}{kg}, \SI{2}{kg} and \SI{3}{kg} are placed on a straight line at positions $x=0$, $x=4$ and $x=10$ (in metres). Find the position of the centre of mass, and verify that it lies closer to the heaviest particle than the midpoint of the extreme particles.
\end{exercice}

\begin{corrige}[Exercise 1]
\[
\bar{x}=\frac{1(0)+2(4)+3(10)}{1+2+3}=\frac{0+8+30}{6}=\frac{38}{6}=\frac{19}{3}\approx 6.33\ \mathrm{m}.
\]
The midpoint of the extreme particles is at $x=5$; since $6.33>5$, $G$ is indeed shifted towards the heaviest mass at $x=10$, as expected.
\end{corrige}

\begin{exercice}[Motion of the mass centre]
Two men of masses \SI{70}{kg} and \SI{50}{kg} stand at the ends of a light plank of length \SI{6}{m} lying on smooth ice. The heavier man walks \SI{2}{m} towards the other. Using the fact that no external horizontal force acts on the system, determine how far the plank slides on the ice.
\end{exercice}

\begin{corrige}[Exercise 2]
The centre of mass of the system (men $+$ light plank) remains fixed in space because there is no external horizontal force. Take the initial position of the heavy man as origin; the light man is at $x=6$. Initially,
\[
\bar{x}=\frac{70(0)+50(6)}{120}=\frac{300}{120}=2.5\ \mathrm{m}.
\]
Suppose the plank slides a distance $d$ in the direction opposite to the walk (i.e. negative $x$ direction). The heavy man, who walked \SI{2}{m} along the plank towards the light man, is then at $x=2-d$, and the light man at $x=6-d$. Since $G$ does not move,
\[
\frac{70(2-d)+50(6-d)}{120}=2.5
\;\Longrightarrow\;
140-70d+300-50d=300
\;\Longrightarrow\;
120d=140,
\]
so $d=\dfrac{7}{6}\approx 1.17\ \mathrm{m}$. The plank slides about \SI{1.17}{m} in the direction opposite to the man's walk — a direct consequence of the conservation of the position of the mass centre.
\end{corrige}

\begin{aretenir}[To remember]
\begin{itemize}
\item The centre of mass is the \textbf{mass-weighted average} position: $M\bar{x}=\sum m_ix_i$ (and similarly for $y$, $z$).
\item In a uniform gravitational field, centre of mass and centre of gravity \textbf{coincide}.
\item For translational motion, the whole body behaves as a particle of mass $M$ at $G$: $M\ddot{\vec{r}}_G=\sum\vec{F}^{\,\mathrm{ext}}$.
\item Internal forces never move $G$; if the resultant external force is zero, $G$ keeps a constant velocity.
\item $G$ may lie \textbf{outside} the material of the body.
\end{itemize}
\end{aretenir}

\coursec{Centre of Mass of a Uniform Rigid Body}

In the previous section we defined the centre of mass of a system of particles and saw that the motion of any body, however complicated, can be split into the motion of its mass centre plus a rotation about it. We now apply these ideas to the objects we actually meet in mechanics problems: rods, plates, rings, cones and spheres. The key simplifying assumption is that these bodies are \cle{uniform}.

\courssub{Uniform Bodies: Mass Proportional to Size}

\begin{definition}[Uniform rigid body]
A \cle{uniform rigid body} is a body whose \cle{density is constant} throughout, and whose parts are rigidly fixed relative to one another. Consequently:
\begin{itemize}
\item for a uniform \textbf{wire or rod} of line density $\lambda$ (mass per unit length), the mass of a piece of length $l$ is $m=\lambda l$: mass is proportional to \textbf{length};
\item for a uniform \textbf{lamina} of surface density $\sigma$ (mass per unit area), the mass of a piece of area $A$ is $m=\sigma A$: mass is proportional to \textbf{area};
\item for a uniform \textbf{solid} of density $\rho$ (mass per unit volume), the mass of a piece of volume $V$ is $m=\rho V$: mass is proportional to \textbf{volume}.
\end{itemize}
\end{definition}

This proportionality has a powerful consequence. When we combine or compare parts of a uniform body, we may replace masses by lengths, areas or volumes in the centre of mass formula, because the constant density cancels:
\[
\bar{x}=\frac{\sum m_i x_i}{\sum m_i}=\frac{\sum \rho V_i x_i}{\sum \rho V_i}=\frac{\sum V_i x_i}{\sum V_i}.
\]
This is why, for uniform bodies, the centre of mass depends only on the \emph{geometry} of the body. For a uniform body the centre of mass coincides with the geometric \cle{centroid}.

\courssub{The Uniform Rod}

\begin{experience}[Balancing a ruler]
Place a uniform 30 cm ruler horizontally on one finger. It balances when your finger is under the 15 cm mark — the midpoint. Slide the finger a few centimetres either way and the ruler tips. The balance point is the centre of mass.
\end{experience}

\begin{propriete}[Centre of mass of a uniform rod]
The centre of mass of a uniform rod $AB$ is at its \cle{midpoint}.
\end{propriete}

\textbf{Justification by symmetry.} For every small element of the rod at distance $d$ to the left of the midpoint, there is an identical element of equal mass at distance $d$ to the right. Their contributions to the moment about the midpoint cancel in pairs, so the total moment about the midpoint is zero: the centre of mass is the midpoint.

\textbf{Justification by summation.} Let the rod have length $L$ and line density $\lambda$, and measure $x$ from one end. Divide the rod into $n$ equal small pieces of length $\Delta x = L/n$, the $i$-th piece having mass $m_i=\lambda\Delta x$ and position $x_i$. Then
\[
\bar{x}=\frac{\sum_{i=1}^{n} \lambda\,\Delta x\, x_i}{\sum_{i=1}^{n}\lambda\,\Delta x}
=\frac{1}{L}\sum_{i=1}^{n}x_i\,\Delta x
\;\xrightarrow[n\to\infty]{}\;
\frac{1}{L}\int_0^L x\,dx=\frac{1}{L}\cdot\frac{L^2}{2}=\frac{L}{2}.
\]
The centre of mass is at the midpoint, as expected. (The integral step previews the calculus method studied later; the symmetry argument alone is sufficient at this stage.)

\courssub{Symmetry Arguments for Standard Bodies}

Symmetry is the fastest tool for locating centres of mass. The guiding principle is:

\begin{aretenir}[Symmetry principle]
If a uniform body has a line (or plane) of symmetry, its centre of mass lies on that line (or plane). If it has two or more axes of symmetry, the centre of mass is at their intersection.
\end{aretenir}

\begin{itemize}
\item \textbf{Uniform circular ring:} every diameter is an axis of symmetry, so the centre of mass is at the \textbf{centre} of the ring — even though no material is there!
\item \textbf{Uniform circular disc:} same argument; centre of mass at the \textbf{centre}.
\item \textbf{Uniform rectangular lamina:} the two midlines (and the diagonals) are symmetry-related bisectors; the centre of mass is at the \textbf{intersection of the diagonals}.
\item \textbf{Uniform cylinder:} centre of mass at the \textbf{midpoint of its axis}.
\item \textbf{Uniform solid sphere:} centre of mass at its \textbf{centre}.
\end{itemize}

\courssub{The Uniform Triangular Lamina}

\begin{propriete}[Centre of mass of a uniform triangular lamina]
The centre of mass of a uniform triangular lamina coincides with the \cle{centroid} of the triangle: the point of intersection of the three \textbf{medians}. It lies on each median \textbf{one third of the way from the midpoint of the corresponding side} (equivalently, two thirds of the way from the vertex).
\end{propriete}

\textbf{Justification.} Divide the triangle into thin strips parallel to side $BC$. Each strip is approximately a uniform rod, so its centre of mass is at its midpoint. The midpoints of all these strips lie on the median from $A$ to the midpoint $M$ of $BC$. Hence the centre of mass of the whole lamina lies on the median $AM$. Repeating the argument with strips parallel to the other sides shows the centre of mass lies on all three medians; it is therefore their intersection, the centroid $G$. The median division ratio $AG:GM = 2:1$ gives the "one third from the base" statement.

\begin{popfigure}
\begin{center}
\begin{tikzpicture}[scale=1.5, line cap=round]
\coordinate (A) at (0,2.6);
\coordinate (B) at (-1.8,0);
\coordinate (C) at (2.2,0);
\coordinate (M) at (0.2,0);
\coordinate (N) at (1.1,1.3);
\coordinate (P) at (-0.9,1.3);
\coordinate (G) at (0.133,0.867);
\fill[popblueL] (A) -- (B) -- (C) -- cycle;
\draw[thick, popink] (A) -- (B) -- (C) -- cycle;
\draw[dashed, poporange, thick] (A) -- (M);
\draw[dashed, popgreen, thick] (B) -- (N);
\draw[dashed, poppurple, thick] (C) -- (P);
\fill[popdark] (G) circle (0.045);
\node[popdark, right] at (G) {$G$ (centroid)};
\node[above] at (A) {$A$};
\node[below left] at (B) {$B$};
\node[below right] at (C) {$C$};
\node[below] at (M) {$M$};
\node[poporange, right, align=left] at (0.55,0.35) {$GM=\tfrac{1}{3}AM$};
\end{tikzpicture}
\legende{The medians of a uniform triangular lamina meet at the centroid $G$, with $AG:GM=2:1$.}
\end{center}
\end{popfigure}

\courssub{The Uniform Semicircular Lamina}

\begin{propriete}[Centre of mass of a uniform semicircular lamina]
The centre of mass of a uniform semicircular lamina of radius $r$ lies on the axis of symmetry, at distance
\[
\bar{x}=\frac{4r}{3\pi}
\]
from the centre of the diameter (i.e. from the flat edge).
\end{propriete}

This result is \emph{quoted} here; its proof requires integration and will be given when we study centroids of regions bounded by curves. Note that $\frac{4}{3\pi}\approx 0.424$, so the centre of mass is a little less than halfway from the flat edge to the arc — reasonable, since more material is concentrated near the diameter.

\courssub{Standard Uniform Solids}

\begin{propriete}[Centres of mass of standard solids]
For uniform solids:
\begin{itemize}
\item \textbf{Solid sphere} of radius $r$: at the centre.
\item \textbf{Solid hemisphere} of radius $r$: on the axis of symmetry, at distance $\dfrac{3r}{8}$ from the centre of the flat face.
\item \textbf{Solid cone or pyramid} of height $h$: on the line joining the apex to the centre of mass of the base, at distance $\dfrac{h}{4}$ from the \textbf{base} (equivalently $\dfrac{3h}{4}$ from the apex).
\item \textbf{Solid cylinder} of height $h$: at the midpoint of its axis, distance $\dfrac{h}{2}$ from either end.
\end{itemize}
\end{propriete}

\begin{attention}[Cone: measure from the base, not the vertex]
The centre of mass of a uniform solid cone is at $\dfrac{h}{4}$ from the \textbf{base}, i.e. $\dfrac{3h}{4}$ from the vertex. Candidates frequently quote "$h/4$ from the vertex" and lose the marks. Remember: a cone has most of its mass near the base, so the centre of mass must be \emph{closer to the base than to the vertex} — indeed $h/4 < h/2$. The same check applies to the hemisphere: $\frac{3r}{8}<\frac{r}{2}$, so the centre of mass is nearer the flat face than the pole, as expected.
\end{attention}

\courssub{Table of Standard Results}

\begin{propriete}[Standard centres of mass to learn for the GCE]
The following results are quoted freely in the examination; learn them.
\begin{center}
\rowcolors{2}{popblueL}{white}
\begin{tabularx}{\linewidth}{|l|X|}
\hline
\rowcolor{popblueL}\textbf{Uniform body} & \textbf{Position of centre of mass} \\
\hline
Rod of length $L$ & Midpoint, $\frac{L}{2}$ from either end \\
Circular ring, radius $r$ & Centre of the ring \\
Circular disc, radius $r$ & Centre of the disc \\
Rectangular lamina & Intersection of the diagonals \\
Triangular lamina & Centroid: $\frac{1}{3}$ of the median from the base \\
Semicircular lamina, radius $r$ & $\frac{4r}{3\pi}$ from the centre, on the axis of symmetry \\
Solid sphere, radius $r$ & Centre \\
Solid hemisphere, radius $r$ & $\frac{3r}{8}$ from the flat face, on the axis \\
Solid cone / pyramid, height $h$ & $\frac{h}{4}$ from the base, on the axis \\
Solid cylinder, height $h$ & Midpoint of the axis, $\frac{h}{2}$ from each end \\
\hline
\end{tabularx}
\end{center}
\end{propriete}

\begin{popfigure}
\begin{center}
\begin{tikzpicture}[scale=0.9, line cap=round]
% Rod
\draw[very thick, popink] (0,0) -- (3,0);
\fill[poporange] (1.5,0) circle (0.06);
\node[below, poporange] at (1.5,-0.05) {$G$};
\node[above] at (1.5,0.05) {rod};
% Rectangle
\draw[thick, popink, fill=popblueL] (4,-0.5) rectangle (6.4,0.7);
\draw[popmuted, thin] (4,-0.5) -- (6.4,0.7);
\draw[popmuted, thin] (4,0.7) -- (6.4,-0.5);
\fill[poporange] (5.2,0.1) circle (0.06);
\node[below, poporange] at (5.2,0.02) {$G$};
\node[above] at (5.2,0.75) {rectangle};
% Triangle
\draw[thick, popink, fill=popgreenL] (7.4,-0.5) -- (9.8,-0.5) -- (8.2,1.0) -- cycle;
\fill[poporange] (8.47,-0.17) circle (0.06);
\node[right, poporange] at (8.55,-0.2) {$G$};
\node[above] at (8.6,1.05) {triangle};
% Semicircle
\draw[thick, popink, fill=popblueL] (0.6,-2.2) arc (180:0:0.9) -- cycle;
\fill[poporange] (1.5,-1.82) circle (0.06);
\node[right, poporange] at (1.58,-1.85) {$G$};
\node[below] at (1.5,-2.3) {semicircle: $OG=\frac{4r}{3\pi}$};
% Cone
\draw[thick, popink, fill=popgreenL] (3.6,-2.4) ellipse (0.8 and 0.22);
\draw[thick, popink] (2.8,-2.4) -- (3.6,-0.9) -- (4.4,-2.4);
\draw[dashed, popmuted] (3.6,-0.9) -- (3.6,-2.4);
\fill[poporange] (3.6,-2.03) circle (0.06);
\node[right, poporange] at (3.68,-2.05) {$G$};
\node[below] at (3.6,-2.75) {cone: $\frac{h}{4}$ from base};
% Hemisphere
\draw[thick, popink, fill=popblueL] (6.0,-2.4) arc (180:0:0.85) -- cycle;
\draw[thick, popink] (6.0,-2.4) ellipse (0.85 and 0.2);
\fill[poporange] (6.85,-2.08) circle (0.06);
\node[right, poporange] at (6.93,-2.1) {$G$};
\node[below] at (6.85,-2.75) {hemisphere: $\frac{3r}{8}$ from face};
\end{tikzpicture}
\legende{Standard uniform bodies with their centres of mass marked at $G$.}
\end{center}
\end{popfigure}

\begin{methode}[Using standard results instead of re-deriving them]
\begin{enumerate}
\item \textbf{Identify} the shape of each uniform part of the body (rod, disc, triangle, cone, ...).
\item \textbf{Quote} the standard position of its centre of mass from the table above — no proof is required in the examination unless the question explicitly asks for one.
\item \textbf{Model} each part as a particle of the same total mass placed at its centre of mass.
\item \textbf{Apply} the centre of mass formula $\bar{x}=\dfrac{\sum m_i x_i}{\sum m_i}$ to the resulting system of particles, using lengths, areas or volumes in place of masses (density cancels).
\item \textbf{Check} plausibility: the answer must lie within the body (or on its axis of symmetry) and closer to the heavier or larger part.
\end{enumerate}
\end{methode}

\begin{exoresolu}[Balance point of a blacksmith's hammer]
A uniform hammer head is modelled as a solid cylinder of mass $1.2\ \mathrm{kg}$, and its uniform wooden handle as a rod of length $30\ \mathrm{cm}$ and mass $0.3\ \mathrm{kg}$. The handle is attached to the centre of the head. Find the distance of the centre of mass of the hammer from the centre of the head.
\tcblower
\textbf{Solution.} Take the centre of the head as origin, with the $x$-axis along the handle.
\begin{itemize}
\item Head (uniform cylinder): $m_1 = 1.2\ \mathrm{kg}$, centre of mass at $x_1 = 0$.
\item Handle (uniform rod): $m_2 = 0.3\ \mathrm{kg}$, centre of mass at its midpoint, $x_2 = 15\ \mathrm{cm}$.
\end{itemize}
Apply the formula:
\[
\bar{x}=\frac{m_1x_1+m_2x_2}{m_1+m_2}
=\frac{(1.2)(0)+(0.3)(15)}{1.2+0.3}
=\frac{4.5}{1.5}=3\ \mathrm{cm}.
\]
The hammer balances at a point $3\ \mathrm{cm}$ from the centre of the head along the handle. The small value is plausible: the head is four times as heavy as the handle, so the balance point lies close to the head.
\end{exoresolu}

\begin{exercice}[Balancing a composite tool]
A uniform metal rod $AB$ of length $80\ \mathrm{cm}$ and mass $2\ \mathrm{kg}$ has a uniform circular disc of mass $1\ \mathrm{kg}$ and radius $10\ \mathrm{cm}$ fixed at end $A$, with the centre of the disc at $A$. Find the distance of the centre of mass of the tool from $A$, and state where a support must be placed for the tool to rest horizontally.
\end{exercice}

\begin{corrige}[Exercise 1]
Measure $x$ from $A$. Disc: $m_1=1\ \mathrm{kg}$ at $x_1=0$. Rod: $m_2=2\ \mathrm{kg}$ at its midpoint, $x_2=40\ \mathrm{cm}$.
\[
\bar{x}=\frac{(1)(0)+(2)(40)}{1+2}=\frac{80}{3}\approx 26.7\ \mathrm{cm}.
\]
The support must be placed $26.7\ \mathrm{cm}$ from $A$, directly below the centre of mass, for the tool to balance horizontally.
\end{corrige}

\begin{savaistu}[Balance in everyday Cameroonian life]
A porter carrying a load on a head-pan, or a carpenter testing a new hoe, instinctively finds the centre of mass: the tool feels comfortable only when it is gripped or supported near its balance point. Market traders in Douala and Bamenda who weigh goods with a simple beam balance are applying exactly the principle of moments about the centre of mass that you have just used.
\end{savaistu}

\begin{attention}[Common mistakes]
\begin{itemize}
\item Do not average the \emph{positions} of the parts; average them \emph{weighted by mass} (or by length, area, volume for uniform bodies).
\item For a uniform body, never include the density $\rho$ in your final calculation — it cancels. Including it only creates opportunities for arithmetic errors.
\item The centre of mass of a ring or hollow body can lie in empty space; do not reject an answer merely because no material is present at $G$.
\item When a body has an axis of symmetry, use it: choose it as a coordinate axis and the perpendicular coordinate of $G$ is automatically zero.
\end{itemize}
\end{attention}

\begin{aretenir}[Key points]
\begin{itemize}
\item \textbf{Uniform} means constant density: mass $\propto$ length, area or volume, so the centre of mass of a uniform body is a purely geometric point (the centroid).
\item Rod: midpoint; ring, disc, sphere: centre; rectangle: intersection of diagonals; triangle: centroid at $\frac{1}{3}$ of the median from the base.
\item Semicircular lamina: $\frac{4r}{3\pi}$ from the centre; hemisphere: $\frac{3r}{8}$ from the flat face; cone or pyramid: $\frac{h}{4}$ from the \textbf{base}.
\item Quote standard results directly, model each part as a particle at its centre of mass, then apply $\bar{x}=\dfrac{\sum m_ix_i}{\sum m_i}$.
\end{itemize}
\end{aretenir}

\coursec{Centre of Mass Using Coordinates}

In the previous section we located the centre of mass of a uniform rigid body by exploiting its symmetry: the centre of mass of a uniform rod lies at its midpoint, that of a uniform disc at its geometric centre, and so on. Symmetry arguments, however, have a limit. When the masses are \emph{not} uniformly distributed — a set of point masses scattered along a line, in a plane, or in space — we need a systematic, algebraic tool. That tool is the \cle{coordinate method}: we attach numbers (coordinates) to the positions of the particles, and the centre of mass becomes a simple weighted average of those numbers. This section develops the method in one, two and three dimensions, and shows how to use it both forwards (finding the centre of mass) and backwards (finding an unknown mass or position).

\courssub{Particles on a Line: the One-Dimensional Formula}

\textbf{Intuition from the balance point.}\par
Imagine a light horizontal rod resting on a knife-edge, with a bag of cement of mass \SI{25}{kg} hung at one point and a bag of rice of mass \SI{10}{kg} hung at another. Experience tells us the rod balances when the knife-edge is placed \emph{closer to the heavier bag}. The balance point is exactly the centre of mass, and its position depends on two things: \emph{where} each mass is, and \emph{how heavy} each mass is. The heavier mass "pulls" the balance point towards itself.

Place an origin $O$ on the rod and measure positions along an axis $Ox$. Suppose particles of masses $m_1, m_2, \dots, m_n$ are placed at positions $x_1, x_2, \dots, x_n$. The balance condition (equality of moments about the centre of mass $\bar{x}$) gives
\[
\sum_{i=1}^{n} m_i (x_i - \bar{x}) = 0
\quad\Longrightarrow\quad
\bar{x}\sum_{i=1}^{n} m_i = \sum_{i=1}^{n} m_i x_i .
\]

\begin{definition}[Centre of mass of particles on a line]
For particles of masses $m_1, m_2, \dots, m_n$ placed at coordinates $x_1, x_2, \dots, x_n$ on a line, the centre of mass $G$ has coordinate
\[
\bar{x} = \frac{\displaystyle\sum_{i=1}^{n} m_i x_i}{\displaystyle\sum_{i=1}^{n} m_i} = \frac{\sum m x}{M}, \qquad \text{where } M = \sum_{i=1}^{n} m_i .
\]
\end{definition}

The numerator $\sum mx$ is the \cle{total moment of mass} about the origin; the denominator is the \cle{total mass}. The formula says: \emph{the whole system behaves, for the purpose of locating $G$, as a single particle of mass $M$ placed at $\bar{x}$, with the same total moment about $O$ as the original system.}

\begin{exemplebox}[Two bags on a rod]
A \SI{25}{kg} bag of cement is at $x = 0$ and a \SI{10}{kg} bag of rice is at $x = \SI{1.4}{m}$ on a light rod. Then
\[
\bar{x} = \frac{25(0) + 10(1.4)}{25 + 10} = \frac{14}{35} = \SI{0.4}{m}.
\]
The balance point is \SI{0.4}{m} from the cement — closer to the heavier mass, exactly as intuition predicts.
\end{exemplebox}

\courssub{Particles in a Plane: Two Dimensions}

When particles lie in a plane, we choose axes $Ox$ and $Oy$ and give each particle coordinates $(x_i, y_i)$. The centre of mass is found by applying the one-dimensional idea \emph{separately} to each coordinate: the $x$-coordinate of $G$ is the weighted average of the $x$-coordinates of the particles, and similarly for the $y$-coordinate. Equivalently, if $\vec{r}_i = (x_i, y_i)$ is the position vector of the $i$-th particle, the vector formula $\displaystyle \vec{R}_G = \frac{\sum m_i \vec{r}_i}{\sum m_i}$ yields the two scalar equations below.

\begin{definition}[Centre of mass of particles in a plane]
For particles of masses $m_i$ at points $(x_i, y_i)$, $i = 1, \dots, n$, the centre of mass $G$ has coordinates
\[
\bar{x} = \frac{\sum m_i x_i}{\sum m_i}, \qquad \bar{y} = \frac{\sum m_i y_i}{\sum m_i}.
\]
\end{definition}

\begin{methode}[Tabulating before summing]
For any centre-of-mass calculation by coordinates:
\begin{enumerate}
\item Choose a convenient origin and axes; write down the coordinates of every particle.
\item Draw a table with columns $m$, $x$, $y$, $mx$, $my$ (add $z$ and $mz$ in 3D).
\item Fill in one row per particle, then sum each column.
\item Compute $\bar{x} = \dfrac{\sum mx}{\sum m}$ and $\bar{y} = \dfrac{\sum my}{\sum m}$.
\item Check: $G$ must lie "inside" the cloud of particles, biased towards the heavier masses.
\end{enumerate}
\end{methode}

\begin{exoresolu}[Particles at the corners of a rectangle]
Particles of masses $2\ \mathrm{kg}$, $3\ \mathrm{kg}$, $1\ \mathrm{kg}$ and $4\ \mathrm{kg}$ are placed at the corners $A(0,0)$, $B(6,0)$, $C(6,4)$ and $D(0,4)$ of a rectangle respectively. Find the centre of mass.
\tcblower
\textbf{Solution.} Tabulate:
\begin{center}
\begin{tabular}{|c|c|c|c|c|}
\hline
\rowcolor{popblueL} $m$ & $x$ & $y$ & $mx$ & $my$ \\
\hline
$2$ & $0$ & $0$ & $0$ & $0$ \\
\hline
$3$ & $6$ & $0$ & $18$ & $0$ \\
\hline
$1$ & $6$ & $4$ & $6$ & $4$ \\
\hline
$4$ & $0$ & $4$ & $0$ & $16$ \\
\hline
\rowcolor{popblueL} $\sum m = 10$ & & & $\sum mx = 24$ & $\sum my = 20$ \\
\hline
\end{tabular}
\end{center}
Hence
\[
\bar{x} = \frac{24}{10} = 2.4, \qquad \bar{y} = \frac{20}{10} = 2.
\]
The centre of mass is $G(2.4,\ 2)$. Note that $G$ lies to the \emph{left} of the rectangle's geometric centre $(3,2)$ because the total mass on the left edge ($2+4=6\ \mathrm{kg}$) is greater than that on the right edge ($3+1=4\ \mathrm{kg}$). The $y$-coordinate is exactly the mid-height $y=2$ because the total mass on the bottom edge ($2+3=5\ \mathrm{kg}$) equals that on the top edge ($1+4=5\ \mathrm{kg}$).
\end{exoresolu}

\begin{popfigure}
\begin{center}
\begin{tikzpicture}[scale=0.85]
\draw[step=1cm, gray!25, very thin] (-0.5,-0.5) grid (6.5,4.5);
\draw[->, thick] (-0.5,0) -- (6.8,0) node[right] {$x$};
\draw[->, thick] (0,-0.5) -- (0,4.8) node[above] {$y$};
\draw[popmuted, thick] (0,0) rectangle (6,4);
\foreach \x in {1,...,6} \draw (\x,0.05) -- (\x,-0.05) node[below, font=\small] {\x};
\foreach \y in {1,...,4} \draw (0.05,\y) -- (-0.05,\y) node[left, font=\small] {\y};
\fill[popblue] (0,0) circle (0.14) node[below left, font=\small] {$A$: $2$ kg};
\fill[popgreen] (6,0) circle (0.14) node[below right, font=\small] {$B$: $3$ kg};
\fill[poporange] (6,4) circle (0.14) node[above right, font=\small] {$C$: $1$ kg};
\fill[poppurple] (0,4) circle (0.14) node[above left, font=\small] {$D$: $4$ kg};
\fill[popink] (2.4,2) circle (0.12);
\draw[popink] (2.4,2) node[above right=2pt, font=\small\bfseries] {$G(2.4,\,2)$};
\draw[popink, dashed] (2.4,2) -- (2.4,0);
\draw[popink, dashed] (2.4,2) -- (0,2);
\end{tikzpicture}
\end{center}
\legende{Four particles at the corners of a $6 \times 4$ rectangle and their centre of mass $G(2.4,\,2)$.}
\end{popfigure}

\begin{attention}[The most common error: forgetting to divide by $\sum m$]
The quantities $\sum mx$ and $\sum my$ are \emph{moments of mass}, not coordinates. A candidate who writes $\bar{x} = \sum mx$ obtains an answer with the wrong units (kg$\cdot$m instead of m) and usually an absurd position. Always finish the calculation: $\bar{x} = \dfrac{\sum mx}{\sum m}$. A quick sanity check: $\bar{x}$ must lie between the smallest and the largest of the $x_i$; if your answer is outside that range, you have almost certainly forgotten the division (or mis-copied a coordinate).
\end{attention}

\courssub{Particles in Space: Three Dimensions}

The extension to three dimensions is immediate: introduce a third axis $Oz$ and average the $z$-coordinates in exactly the same way.

\begin{definition}[Centre of mass of particles in space]
For particles of masses $m_i$ at points $(x_i, y_i, z_i)$, $i = 1, \dots, n$, the centre of mass $G$ has coordinates
\[
\bar{x} = \frac{\sum m_i x_i}{\sum m_i}, \qquad
\bar{y} = \frac{\sum m_i y_i}{\sum m_i}, \qquad
\bar{z} = \frac{\sum m_i z_i}{\sum m_i}.
\]
\end{definition}

\begin{exoresolu}[Three particles in space]
Particles of masses $2\ \mathrm{kg}$, $5\ \mathrm{kg}$ and $3\ \mathrm{kg}$ are placed at the points $P(1,0,2)$, $Q(3,4,0)$ and $R(0,2,6)$ respectively. Find the coordinates of the centre of mass.
\tcblower
\textbf{Solution.} Tabulate:
\begin{center}
\begin{tabular}{|c|c|c|c|c|c|c|}
\hline
\rowcolor{popblueL} $m$ & $x$ & $y$ & $z$ & $mx$ & $my$ & $mz$ \\
\hline
$2$ & $1$ & $0$ & $2$ & $2$ & $0$ & $4$ \\
\hline
$5$ & $3$ & $4$ & $0$ & $15$ & $20$ & $0$ \\
\hline
$3$ & $0$ & $2$ & $6$ & $0$ & $6$ & $18$ \\
\hline
\rowcolor{popblueL} $10$ & & & & $17$ & $26$ & $22$ \\
\hline
\end{tabular}
\end{center}
Therefore
\[
\bar{x} = \frac{17}{10} = 1.7, \qquad \bar{y} = \frac{26}{10} = 2.6, \qquad \bar{z} = \frac{22}{10} = 2.2.
\]
The centre of mass is $G(1.7,\ 2.6,\ 2.2)$.
\end{exoresolu}

\begin{popfigure}
\begin{center}
\begin{tikzpicture}[scale=1.0, x={(1cm,0cm)}, y={(0.55cm,0.35cm)}, z={(0cm,1cm)}]
\draw[->, thick] (0,0,0) -- (4.2,0,0) node[right] {$x$};
\draw[->, thick] (0,0,0) -- (0,5,0) node[above right, font=\small] {$y$};
\draw[->, thick] (0,0,0) -- (0,0,6.6) node[above] {$z$};
\draw[gray!40, dashed] (1,0,0) -- (1,0,2) -- (0,0,2);
\draw[gray!40, dashed] (3,0,0) -- (3,4,0) -- (0,4,0);
\draw[gray!40, dashed] (0,2,0) -- (0,2,6) -- (0,0,6);
\fill[popblue] (1,0,2) circle (0.09) node[left, font=\small] {$P$: $2$ kg};
\fill[popgreen] (3,4,0) circle (0.09) node[right, font=\small] {$Q$: $5$ kg};
\fill[poporange] (0,2,6) circle (0.09) node[above left, font=\small] {$R$: $3$ kg};
\fill[popink] (1.7,2.6,2.2) circle (0.10);
\node[popink, font=\small\bfseries] at (2.6,3.1,2.6) {$G(1.7,\,2.6,\,2.2)$};
\draw[popink, dashed] (1.7,2.6,0) -- (1.7,2.6,2.2);
\draw[popink, dashed] (1.7,0,0) -- (1.7,2.6,0) -- (0,2.6,0);
\end{tikzpicture}
\end{center}
\legende{Three particles in space and their centre of mass $G$, with dashed projection lines onto the coordinate planes.}
\end{popfigure}

\courssub{Choosing a Convenient Origin and Axes}

The centre of mass is a \emph{physical} point: it does not depend on the coordinate system we use to describe it. We are therefore free to choose the origin and axes so as to minimise the work.

\begin{methode}[Choosing axes wisely]
\begin{enumerate}
\item Put the origin \emph{on one of the particles} (or at a corner of the figure): that particle then contributes $0$ to $\sum mx$ and $\sum my$.
\item Align the axes with the sides of the figure (rectangle, rod, plate) so that most coordinates are $0$ or simple numbers.
\item Use any line of symmetry as an axis: the centre of mass always lies on every axis of symmetry of the mass distribution, so the coordinate perpendicular to that axis needs no computation.
\end{enumerate}
\end{methode}

\begin{exemplebox}[Masses along a rod]
A rod $AB$ of length \SI{80}{cm} carries masses of $1\ \mathrm{kg}$, $2\ \mathrm{kg}$ and $3\ \mathrm{kg}$ at distances $20\ \mathrm{cm}$, $50\ \mathrm{cm}$ and $80\ \mathrm{cm}$ from $A$. Taking $A$ as origin and the rod as the $x$-axis:
\[
\bar{x} = \frac{1(20) + 2(50) + 3(80)}{1+2+3} = \frac{20 + 100 + 240}{6} = \frac{360}{6} = \SI{60}{cm}.
\]
The centre of mass is \SI{60}{cm} from $A$ — beyond the midpoint, because the heaviest mass sits at the far end $B$.
\end{exemplebox}

\courssub{Working Backwards: Finding an Unknown Mass or Position}

In many GCE questions the centre of mass is \emph{given} and one mass or one coordinate is unknown. The same formula is used, but now it becomes an equation to solve.

\begin{methode}[Working backwards from a given centre of mass]
\begin{enumerate}
\item Write $\bar{x} = \dfrac{\sum mx}{\sum m}$ with the unknown symbol ($m$, $x$, \dots) included in the sums.
\item Cross-multiply: $\bar{x}\,\sum m = \sum mx$.
\item Solve the resulting linear equation for the unknown.
\item Check the physical sense of the answer (a mass must be positive; a coordinate must lie in a plausible range).
\end{enumerate}
\end{methode}

\begin{exoresolu}[Finding an unknown mass]
Particles of mass $4\ \mathrm{kg}$ at $(1,2)$ and $2\ \mathrm{kg}$ at $(5,4)$ are joined by a third particle of mass $m$ kg at $(7,1)$. The $x$-coordinate of the centre of mass of the system is $4$. Find $m$ and the $y$-coordinate of $G$.
\tcblower
\textbf{Solution.} Use the $x$-coordinate formula:
\[
\bar{x} = \frac{4(1) + 2(5) + m(7)}{4 + 2 + m} = \frac{14 + 7m}{6 + m} = 4.
\]
Cross-multiplying:
\[
14 + 7m = 4(6 + m) = 24 + 4m \quad\Longrightarrow\quad 3m = 10 \quad\Longrightarrow\quad m = \tfrac{10}{3}\ \mathrm{kg}.
\]
Now find $\bar{y}$ with this value of $m$:
\[
\bar{y} = \frac{4(2) + 2(4) + \tfrac{10}{3}(1)}{6 + \tfrac{10}{3}} = \frac{8 + 8 + \tfrac{10}{3}}{\tfrac{28}{3}} = \frac{\tfrac{58}{3}}{\tfrac{28}{3}} = \frac{58}{28} = \frac{29}{14} \approx 2.07.
\]
Thus $G = \left(4,\ \frac{29}{14}\right)$.
\end{exoresolu}

\begin{exoresolu}[Finding an unknown coordinate]
Three particles, each of mass $2\ \mathrm{kg}$, are placed at $(0,0)$, $(6,0)$ and $(a, 3)$. The centre of mass lies on the line $x = 4$. Find $a$ and the coordinates of $G$.
\tcblower
\textbf{Solution.} Since the masses are equal,
\[
\bar{x} = \frac{0 + 6 + a}{3} = 4 \quad\Longrightarrow\quad 6 + a = 12 \quad\Longrightarrow\quad a = 6.
\]
Then
\[
\bar{y} = \frac{0 + 0 + 3}{3} = 1,
\]
so $G = (4, 1)$. Note that with equal masses the centre of mass is simply the \emph{average} of the coordinates — the centroid of the triangle formed by the three points.
\end{exoresolu}

\begin{exercice}[Practice]
\begin{enumerate}
\item Particles of masses $3\ \mathrm{kg}$, $5\ \mathrm{kg}$ and $2\ \mathrm{kg}$ are placed at $(0,0)$, $(4,0)$ and $(4,6)$. Find the centre of mass.
\item Four particles of masses $1\ \mathrm{kg}$, $2\ \mathrm{kg}$, $3\ \mathrm{kg}$ and $m\ \mathrm{kg}$ are at $(0,0)$, $(2,0)$, $(2,2)$ and $(0,2)$. Given that $\bar{x} = 1$, find $m$ and $\bar{y}$.
\item Particles of mass $2\ \mathrm{kg}$ at $(1,1,1)$, $3\ \mathrm{kg}$ at $(3,0,2)$ and $5\ \mathrm{kg}$ at $(0,2,4)$. Find $G$.
\end{enumerate}
\end{exercice}

\begin{corrige}[Exercise 1]
\begin{enumerate}
\item $\bar{x} = \dfrac{3(0)+5(4)+2(4)}{10} = \dfrac{28}{10} = 2.8$;\quad $\bar{y} = \dfrac{3(0)+5(0)+2(6)}{10} = \dfrac{12}{10} = 1.2$. So $G = (2.8,\ 1.2)$.
\item $\bar{x} = \dfrac{0 + 4 + 6 + 0}{6 + m} = \dfrac{10}{6+m} = 1 \Rightarrow m = 4\ \mathrm{kg}$. Then $\bar{y} = \dfrac{0+0+6+8}{10} = 1.4$.
\item $\bar{x} = \dfrac{2+9+0}{10} = 1.1$;\quad $\bar{y} = \dfrac{2+0+10}{10} = 1.2$;\quad $\bar{z} = \dfrac{2+6+20}{10} = 2.8$. So $G = (1.1,\ 1.2,\ 2.8)$.
\end{enumerate}
\end{corrige}

\begin{savaistu}[Centre of mass and the weighted mean in statistics]
The formula $\bar{x} = \dfrac{\sum m_i x_i}{\sum m_i}$ is exactly the \cle{weighted mean} you meet in statistics, where the masses play the role of frequencies or weights. The mean mark of a class, $\bar{x} = \dfrac{\sum f_i x_i}{\sum f_i}$, is the "centre of mass" of the marks: each mark $x_i$ is counted with weight $f_i$. Physics and statistics share the same mathematics here — a weighted average is a balance point. (This cross-curricular remark is enrichment, not examinable material.)
\end{savaistu}

\begin{aretenir}[Key points of the section]
\begin{itemize}
\item On a line: $\bar{x} = \dfrac{\sum mx}{\sum m}$; in a plane add $\bar{y} = \dfrac{\sum my}{\sum m}$; in space add $\bar{z} = \dfrac{\sum mz}{\sum m}$.
\item Each coordinate of $G$ is computed \emph{independently} as a weighted average of the corresponding coordinates of the particles.
\item Tabulate $m$, $x$, $y$, $(z)$, $mx$, $my$, $(mz)$ before summing — this prevents sign and arithmetic errors.
\item Never forget to divide by the total mass $M = \sum m$; check that $G$ lies within the spread of the particles, biased towards the heavier ones.
\item Choose the origin on a particle and axes along lines of symmetry to simplify the sums.
\item Given $\bar{x}$, find an unknown mass or coordinate by cross-multiplying $\bar{x}\sum m = \sum mx$ and solving the linear equation.
\end{itemize}
\end{aretenir}

\coursec{Position Vectors and the Centre of Mass Formula}

In the previous section we located the centre of mass of a system of particles by giving the \emph{coordinates} $(\bar{x},\bar{y},\bar{z})$ of the point $G$ relative to a chosen set of axes. That method works perfectly, but it forces us to handle the $x$-, $y$- and $z$-coordinates separately, with three distinct formulae. Vector notation allows us to compress all of this into \emph{one single formula}, which is shorter to write, easier to manipulate algebraically, and extremely powerful in examination questions — especially when a particle is added, removed, or when the origin itself is moving. This section develops that vector approach.

\courssub{Revision: Position Vectors in Two and Three Dimensions}

Imagine you are standing at the central roundabout of a town (your \emph{origin} $O$) and you want to tell a friend exactly where the post office is. You could say ``go \SI{300}{m} East and \SI{400}{m} North''. That instruction is a \emph{vector}: it starts at $O$ and ends at the post office $P$. In mechanics we call it the \cle{position vector} of $P$ relative to $O$.

\begin{definition}[Position vector of a point]
Let $O$ be a fixed origin and let $\vec{i}$, $\vec{j}$ (and $\vec{k}$ in space) be unit vectors along the coordinate axes. The \textbf{position vector} of a point $P(x,y)$ in the plane is
\[
\vec{r} = \overrightarrow{OP} = x\,\vec{i} + y\,\vec{j},
\]
and the position vector of a point $P(x,y,z)$ in space is
\[
\vec{r} = \overrightarrow{OP} = x\,\vec{i} + y\,\vec{j} + z\,\vec{k}.
\]
The coefficients $x$, $y$, $z$ are exactly the coordinates of $P$; the position vector is simply a compact way of carrying them all at once.
\end{definition}

Recall the two operations we shall need:
\begin{itemize}
\item \textbf{Scalar multiplication:} $m(x\vec{i}+y\vec{j}+z\vec{k}) = mx\,\vec{i}+my\,\vec{j}+mz\,\vec{k}$ — multiply each component by $m$.
\item \textbf{Addition:} add vectors component by component:
$(x_1\vec{i}+y_1\vec{j})+(x_2\vec{i}+y_2\vec{j}) = (x_1+x_2)\vec{i}+(y_1+y_2)\vec{j}$.
\end{itemize}

\begin{attention}[Position vector versus the point itself]
A position vector always depends on the \emph{chosen origin}. If the origin moves, the coordinates of $P$ change, hence $\vec{r}$ changes — even though the physical point $P$ has not moved at all. Keep this in mind: it will matter when we discuss the centre of mass.
\end{attention}

\courssub{The Centre of Mass Formula in Vector Form}

Consider $n$ particles of masses $m_1, m_2, \dots, m_n$ placed at points with position vectors $\vec{r}_1, \vec{r}_2, \dots, \vec{r}_n$ relative to an origin $O$. From the coordinate section, the centre of mass $G$ has coordinates
\[
\bar{x} = \frac{\sum m_i x_i}{\sum m_i}, \qquad
\bar{y} = \frac{\sum m_i y_i}{\sum m_i}, \qquad
\bar{z} = \frac{\sum m_i z_i}{\sum m_i}.
\]
Now form the position vector of $G$:
\[
\bar{\vec{r}} = \bar{x}\,\vec{i} + \bar{y}\,\vec{j} + \bar{z}\,\vec{k}
= \frac{\sum m_i x_i}{\sum m_i}\,\vec{i} + \frac{\sum m_i y_i}{\sum m_i}\,\vec{j} + \frac{\sum m_i z_i}{\sum m_i}\,\vec{k}
= \frac{\sum m_i (x_i\vec{i}+y_i\vec{j}+z_i\vec{k})}{\sum m_i}.
\]
The numerator is exactly $\sum m_i \vec{r}_i$. Hence the three coordinate formulae merge into one.

\begin{propriete}[Position vector of the centre of mass]
For a system of particles of masses $m_i$ at positions $\vec{r}_i$ ($i = 1,\dots,n$) relative to an origin $O$, the position vector of the centre of mass $G$ is
\[
\boxed{\;\bar{\vec{r}} = \frac{\displaystyle\sum_{i=1}^{n} m_i \vec{r}_i}{\displaystyle\sum_{i=1}^{n} m_i} = \frac{m_1\vec{r}_1 + m_2\vec{r}_2 + \cdots + m_n\vec{r}_n}{m_1+m_2+\cdots+m_n}\;}
\]
The total mass is often written $M = \sum m_i$, so $M\bar{\vec{r}} = \sum m_i \vec{r}_i$. This formula is \textbf{examinable}: you must be able to state it and to derive it from the coordinate formulae, as done above.
\end{propriete}

\begin{aretenir}[One formula, three coordinates]
The vector formula $\bar{\vec{r}} = \dfrac{\sum m_i \vec{r}_i}{\sum m_i}$ is \emph{equivalent} to the three coordinate formulae: the $\vec{i}$-component of $\bar{\vec{r}}$ is $\bar{x}$, the $\vec{j}$-component is $\bar{y}$, and the $\vec{k}$-component is $\bar{z}$. Use whichever form the question suggests.
\end{aretenir}

\begin{methode}[Computing the centre of mass with position vectors]
\begin{enumerate}
\item \textbf{List} the data in a small table: mass $m_i$ and position vector $\vec{r}_i$ of each particle.
\item \textbf{Compute} each product $m_i\vec{r}_i$ (multiply every component by $m_i$).
\item \textbf{Add} the vectors $\sum m_i\vec{r}_i$ component by component, and compute the total mass $M=\sum m_i$.
\item \textbf{Divide} each component of the sum by $M$ to obtain $\bar{\vec{r}}$.
\item \textbf{Check} the plausibility: $G$ must lie ``between'' the particles, closer to the heavier ones.
\end{enumerate}
\end{methode}

\begin{popfigure}
\begin{center}
\begin{tikzpicture}[scale=1.05, >=stealth]
\draw[->, gray] (-0.4,0) -- (7.2,0) node[below left] {$x$};
\draw[->, gray] (0,-0.4) -- (0,4.6) node[below left] {$y$};
\fill (0,0) circle (1.5pt) node[below left] {$O$};
% particles
\fill[popblue] (1,3) circle (3pt) node[above left, popdark] {$m_1$ at $P_1$};
\fill[poporange] (6,1) circle (4.5pt) node[above right, popdark] {$m_2$ at $P_2$};
\fill[poppurple] (4,3.5) circle (3pt) node[above, popdark] {$m_3$ at $P_3$};
% position vectors
\draw[->, thick, popblue] (0,0) -- (1,3) node[midway, left, popdark] {$\vec{r}_1$};
\draw[->, thick, poporange] (0,0) -- (6,1) node[midway, below, popdark] {$\vec{r}_2$};
\draw[->, thick, poppurple] (0,0) -- (4,3.5) node[midway, above right, popdark] {$\vec{r}_3$};
% centre of mass (example: m1=1, m2=2, m3=1 -> G=( (1+12+4)/4, (3+2+3.5)/4 )=(4.25,2.125))
\fill[popgreen] (4.25,2.125) circle (3.5pt) node[below right, popdark] {$G$};
\draw[->, very thick, popgreen] (0,0) -- (4.25,2.125) node[midway, below left, popdark] {$\bar{\vec{r}}$};
\end{tikzpicture}
\end{center}
\legende{Position vectors $\vec{r}_1, \vec{r}_2, \vec{r}_3$ of three particles from the origin $O$, and the position vector $\bar{\vec{r}}$ of their centre of mass $G$. Note that $G$ lies inside the triangle of the particles, pulled towards the heaviest mass.}
\end{popfigure}

\courssub{Worked Examples in Two Dimensions}

\begin{exoresolu}[Two-dimensional system]
Particles of masses \SI{2}{kg}, \SI{3}{kg} and \SI{5}{kg} are placed at the points with position vectors $\vec{r}_1 = \vec{i} + 2\vec{j}$, $\vec{r}_2 = 4\vec{i} - \vec{j}$ and $\vec{r}_3 = -2\vec{i} + 3\vec{j}$ respectively (distances in metres). Find the position vector of the centre of mass.
\tcblower
\textbf{Solution.} Total mass: $M = 2+3+5 = \SI{10}{kg}$.
Compute the moment sum:
\begin{align*}
\sum m_i\vec{r}_i &= 2(\vec{i}+2\vec{j}) + 3(4\vec{i}-\vec{j}) + 5(-2\vec{i}+3\vec{j})\\
&= (2+12-10)\vec{i} + (4-3+15)\vec{j}\\
&= 4\vec{i} + 16\vec{j}.
\end{align*}
Therefore
\[
\bar{\vec{r}} = \frac{4\vec{i}+16\vec{j}}{10} = 0.4\,\vec{i} + 1.6\,\vec{j}\ \text{(m)}.
\]
\emph{Check:} the heaviest particle (\SI{5}{kg}) is at $(-2,3)$; the centre of mass $(0.4,\,1.6)$ is indeed pulled somewhat towards it, and lies within the triangle of the three points. \hfill$\blacksquare$
\end{exoresolu}

\begin{exemplebox}[Reading positions from a diagram]
Three particles lie at the vertices of a triangle $A(0,0)$, $B(6,0)$, $C(0,3)$ with masses \SI{1}{kg}, \SI{2}{kg}, \SI{3}{kg} respectively. Then
\[
\sum m_i\vec{r}_i = 1(0) + 2(6\vec{i}) + 3(3\vec{j}) = 12\vec{i}+9\vec{j}, \qquad M=6,
\]
so $\bar{\vec{r}} = 2\vec{i} + 1.5\vec{j}$. The centre of mass is the point $(2,\,1.5)$ — inside the triangle, nearest to the heaviest vertex $C$ in the $y$-direction.
\end{exemplebox}

\courssub{Worked Examples in Three Dimensions}

Nothing new is needed: the same formula applies, with a $\vec{k}$-component.

\begin{exoresolu}[Three-dimensional system]
Four particles are placed in space:
\begin{itemize}
\item $m_1 = \SI{1}{kg}$ at $\vec{r}_1 = \vec{i} + \vec{j} + \vec{k}$,
\item $m_2 = \SI{2}{kg}$ at $\vec{r}_2 = 2\vec{i} - \vec{k}$,
\item $m_3 = \SI{3}{kg}$ at $\vec{r}_3 = -\vec{i} + 2\vec{j}$,
\item $m_4 = \SI{4}{kg}$ at $\vec{r}_4 = \vec{j} + 2\vec{k}$.
\end{itemize}
Find the position vector of the centre of mass.
\tcblower
\textbf{Solution.} Total mass: $M = 1+2+3+4 = \SI{10}{kg}$.
\begin{align*}
\sum m_i\vec{r}_i &= 1(\vec{i}+\vec{j}+\vec{k}) + 2(2\vec{i}-\vec{k}) + 3(-\vec{i}+2\vec{j}) + 4(\vec{j}+2\vec{k})\\
&= (1+4-3+0)\vec{i} + (1+0+6+4)\vec{j} + (1-2+0+8)\vec{k}\\
&= 2\vec{i} + 11\vec{j} + 7\vec{k}.
\end{align*}
Hence
\[
\bar{\vec{r}} = \frac{2\vec{i}+11\vec{j}+7\vec{k}}{10} = 0.2\,\vec{i} + 1.1\,\vec{j} + 0.7\,\vec{k}.
\]
Each component was computed independently — exactly the coordinate method, but written in one line. \hfill$\blacksquare$
\end{exoresolu}

\begin{attention}[Do not mix components]
When summing $m_i\vec{r}_i$, group the $\vec{i}$-terms together, the $\vec{j}$-terms together, and the $\vec{k}$-terms together. A frequent error at the GCE is to add an $\vec{i}$-coefficient to a $\vec{j}$-coefficient, or to forget the sign of a negative component. Write one bracket per component, as in the solutions above.
\end{attention}

\courssub{Applications of the Position Vector Method}

\textbf{Adding or removing a particle.} Suppose a system already has total mass $M$ and centre of mass at $\bar{\vec{r}}$, so that $\sum m_i\vec{r}_i = M\bar{\vec{r}}$. If a new particle of mass $m$ is added at position $\vec{r}$, the new centre of mass is
\[
\bar{\vec{r}}_{\text{new}} = \frac{M\bar{\vec{r}} + m\vec{r}}{M+m}.
\]
You do \emph{not} need to redo the whole sum: the old system behaves like a single particle of mass $M$ placed at $G$. Similarly, removing a particle of mass $m$ at $\vec{r}$ gives $\bar{\vec{r}}_{\text{new}} = \dfrac{M\bar{\vec{r}} - m\vec{r}}{M-m}$.

\begin{exemplebox}[Adding a particle]
A system of total mass \SI{6}{kg} has its centre of mass at $\bar{\vec{r}} = \vec{i}+2\vec{j}$. A particle of mass \SI{4}{kg} is added at $\vec{r} = 6\vec{i}+7\vec{j}$. The new centre of mass is
\[
\bar{\vec{r}}_{\text{new}} = \frac{6(\vec{i}+2\vec{j}) + 4(6\vec{i}+7\vec{j})}{6+4}
= \frac{30\vec{i}+40\vec{j}}{10} = 3\vec{i}+4\vec{j}.
\]
\end{exemplebox}

\textbf{Change of origin.} If the origin is moved from $O$ to a new origin $O'$ whose position vector relative to $O$ is $\vec{a}$, then every position vector transforms as $\vec{r}\,' = \vec{r} - \vec{a}$, and consequently
\[
\bar{\vec{r}}\,' = \frac{\sum m_i(\vec{r}_i-\vec{a})}{\sum m_i} = \bar{\vec{r}} - \vec{a}.
\]
The centre of mass transforms exactly like any other point.

\begin{attention}[The physical point does not move]
The \emph{position vector} $\bar{\vec{r}}$ of the centre of mass depends on the chosen origin, but the \emph{physical point} $G$ — the balance point of the system — is fixed by the particles themselves. Two candidates using two different origins will obtain two different vectors $\bar{\vec{r}}$, yet both vectors point to the \emph{same} point in space. Never conclude that two answers disagree simply because the coordinates differ: check whether the origins are the same.
\end{attention}

\courssub{Exam-Style Problem: Finding a Missing Particle}

A classic GCE question gives you the centre of mass and all but one particle, and asks for the missing mass or position. The key is to write $M\bar{\vec{r}} = \sum m_i\vec{r}_i$ and solve the resulting vector equation component by component.

\begin{exoresolu}[Determining a missing mass]
Three particles of masses \SI{2}{kg}, \SI{3}{kg} and $m$ kg are placed at the points with position vectors $2\vec{i}+\vec{j}$, $-\vec{i}+4\vec{j}$ and $4\vec{i}-\vec{j}$ respectively. The centre of mass of the system is at $\bar{\vec{r}} = \vec{i}+2\vec{j}$. Find the value of $m$.
\tcblower
\textbf{Solution.} Apply $M\bar{\vec{r}} = \sum m_i\vec{r}_i$ with $M = 2+3+m = 5+m$:
\[
(5+m)(\vec{i}+2\vec{j}) = 2(2\vec{i}+\vec{j}) + 3(-\vec{i}+4\vec{j}) + m(4\vec{i}-\vec{j}).
\]
The right-hand side is
\[
(4-3+4m)\vec{i} + (2+12-m)\vec{j} = (1+4m)\vec{i} + (14-m)\vec{j}.
\]
Equating components:
\begin{itemize}
\item $\vec{i}$-component: $5+m = 1+4m \;\Rightarrow\; 4 = 3m \;\Rightarrow\; m = \dfrac{4}{3}$.
\item $\vec{j}$-component (check): $2(5+m) = 14-m \;\Rightarrow\; 10+2m = 14-m \;\Rightarrow\; 3m = 4 \;\Rightarrow\; m = \dfrac{4}{3}$.
\end{itemize}
Both components give the same value, confirming consistency. Hence $m = \dfrac{4}{3}$ kg. \hfill$\blacksquare$
\end{exoresolu}

\begin{exoresolu}[Determining a missing position vector]
Particles of mass \SI{1}{kg}, \SI{2}{kg} and \SI{3}{kg} are at $\vec{r}_1 = \vec{i}+\vec{j}$, $\vec{r}_2 = 3\vec{i}-\vec{j}$ and at an unknown position $\vec{r}_3$ respectively. The centre of mass is at $\bar{\vec{r}} = 2\vec{i}+\vec{j}$. Find $\vec{r}_3$.
\tcblower
\textbf{Solution.} Total mass $M = 6$, so
\[
6(2\vec{i}+\vec{j}) = 1(\vec{i}+\vec{j}) + 2(3\vec{i}-\vec{j}) + 3\vec{r}_3.
\]
\[
12\vec{i}+6\vec{j} = 7\vec{i}-\vec{j} + 3\vec{r}_3
\;\Rightarrow\;
3\vec{r}_3 = 5\vec{i}+7\vec{j}
\;\Rightarrow\;
\vec{r}_3 = \tfrac{5}{3}\vec{i}+\tfrac{7}{3}\vec{j}.
\]
The third particle is at $\left(\dfrac{5}{3},\,\dfrac{7}{3}\right)$. \hfill$\blacksquare$
\end{exoresolu}

\begin{savaistu}[Barycentric coordinates and astronomy]
The vector formula $\bar{\vec{r}} = \frac{\sum m_i\vec{r}_i}{\sum m_i}$ is the mathematical foundation of \emph{barycentric coordinates}. In astronomy, the centre of mass of the Earth–Moon system (the \emph{barycentre}) lies about \SI{4670}{km} from Earth's centre — inside the Earth but not at its centre. Both bodies orbit this point. For the Sun–Jupiter system, the barycentre lies just outside the Sun's surface! This concept is crucial for detecting exoplanets via the ``wobble'' of their host stars.
\end{savaistu}

\begin{experience}[Balancing a lamina on a pin]
\textbf{Materials:} A piece of stiff cardboard (irregular shape), a pin, a plumb line (string with a weight), a ruler.
\textbf{Procedure:}
\begin{enumerate}
\item Make three small holes near the edge of the cardboard.
\item Suspend the cardboard from one hole using the pin so it swings freely.
\item Hang the plumb line from the same pin; draw a vertical line along the string on the cardboard.
\item Repeat from the other two holes.
\item The three lines intersect at a single point — the centre of mass $G$.
\item Balance the cardboard horizontally on a fingertip placed at $G$; it should stay level.
\end{enumerate}
This demonstrates that the centre of mass is the unique point where the weight of the whole body can be considered to act.
\end{experience}

\courssub{Consolidation Exercises}

\begin{exercice}[Practice]
\begin{enumerate}
\item Particles of masses \SI{1}{kg}, \SI{4}{kg} and \SI{2}{kg} are at $\vec{i}+3\vec{j}$, $2\vec{i}-\vec{j}$ and $5\vec{i}+4\vec{j}$. Find $\bar{\vec{r}}$.
\item Three particles of mass \SI{2}{kg} each are at $\vec{i}+\vec{k}$, $2\vec{j}-\vec{k}$ and $3\vec{i}+\vec{j}+2\vec{k}$. Find the centre of mass.
\item A system of total mass \SI{8}{kg} has centre of mass at $2\vec{i}+\vec{j}$. A particle of mass \SI{2}{kg} at $7\vec{i}+6\vec{j}$ is added. Find the new centre of mass.
\item Particles of masses \SI{3}{kg} and \SI{5}{kg} are at $\vec{i}+2\vec{j}$ and $4\vec{i}-\vec{j}$. A third particle of mass \SI{2}{kg} is placed so that the centre of mass is at $3\vec{i}+\vec{j}$. Find its position vector.
\end{enumerate}
\end{exercice}

\begin{corrige}[Exercise 1]
\textbf{1.} $M=7$; $\sum m_i\vec{r}_i = (1+8+10)\vec{i}+(3-4+8)\vec{j} = 19\vec{i}+7\vec{j}$, so $\bar{\vec{r}} = \dfrac{19}{7}\vec{i}+\vec{j}$.

\textbf{2.} $M=6$; $\sum m_i\vec{r}_i = 2\big[(1+0+3)\vec{i}+(0+2+1)\vec{j}+(1-1+2)\vec{k}\big] = 8\vec{i}+6\vec{j}+4\vec{k}$, so $\bar{\vec{r}} = \dfrac{4}{3}\vec{i}+\vec{j}+\dfrac{2}{3}\vec{k}$.

\textbf{3.} $\bar{\vec{r}}_{\text{new}} = \dfrac{8(2\vec{i}+\vec{j})+2(7\vec{i}+6\vec{j})}{10} = \dfrac{30\vec{i}+20\vec{j}}{10} = 3\vec{i}+2\vec{j}$.

\textbf{4.} $M=10$: $10(3\vec{i}+\vec{j}) = 3(\vec{i}+2\vec{j})+5(4\vec{i}-\vec{j})+2\vec{r}_3 = 23\vec{i}+\vec{j}+2\vec{r}_3$, so $2\vec{r}_3 = 7\vec{i}+9\vec{j}$ and $\vec{r}_3 = \dfrac{7}{2}\vec{i}+\dfrac{9}{2}\vec{j}$.
\end{corrige}

\begin{aretenir}[Key points of the section]
\begin{itemize}
\item The position vector $\vec{r} = x\vec{i}+y\vec{j}+z\vec{k}$ packages the coordinates of a point into a single object relative to a chosen origin.
\item The centre of mass is given by $\bar{\vec{r}} = \dfrac{\sum m_i\vec{r}_i}{\sum m_i}$, equivalent to the three coordinate formulae.
\item Work component by component: compute $\sum m_i\vec{r}_i$, divide by the total mass $M$.
\item An existing system behaves like a single particle of mass $M$ at $G$ when particles are added or removed.
\item $\bar{\vec{r}}$ depends on the origin; the physical balance point $G$ does not.
\item For a missing particle, write $M\bar{\vec{r}} = \sum m_i\vec{r}_i$ and equate components.
\end{itemize}
\end{aretenir}

\coursec{Uniform Laminae, Symmetry and Composite Laminae}

In the previous section we developed the vector formula $\vec{R} = \frac{\sum m_i \vec{r}_i}{\sum m_i}$ for a system of particles. A \emph{lamina} is a continuous distribution of mass over a two‑dimensional region. Because it is made of infinitely many particles, the sums become integrals. However, when the lamina is \emph{uniform} (constant surface density) and can be split into simple geometric shapes whose centres of mass are already known, we can avoid integration entirely by treating each part as a particle located at its own centre of mass. This is the powerful \emph{composite lamina} method.

\courssub{Concept of a Uniform Lamina}

\begin{definition}[Uniform Lamina]
A \cle{uniform lamina} is a thin, flat plate of constant thickness made of a homogeneous material. Its \cle{surface density} (mass per unit area) $\sigma$ is constant everywhere. Consequently, the mass of any portion of the lamina is directly proportional to its area:
\[
m = \sigma \times A.
\]
Since $\sigma$ is the same for all parts, it cancels out in the centre‑of‑mass formulas, and we can work with \emph{areas} instead of masses.
\end{definition}

\begin{aretenir}[Why areas replace masses]
For a uniform lamina divided into parts with areas $A_1, A_2, \dots$ and centres of mass $(\bar{x}_1,\bar{y}_1), (\bar{x}_2,\bar{y}_2), \dots$, the coordinates of the overall centre of mass are
\[
\bar{x} = \frac{\sum A_i \bar{x}_i}{\sum A_i}, \qquad
\bar{y} = \frac{\sum A_i \bar{y}_i}{\sum A_i}.
\]
This is exactly the same formula as for particles, with $m_i$ replaced by $A_i$.
\end{aretenir}

\courssub{Use of Symmetry}

Symmetry is the fastest way to locate the centre of mass of many standard laminae without any calculation.

\begin{propriete}[Centre of Mass and Axes of Symmetry]
If a uniform lamina has an axis of symmetry, its centre of mass lies on that axis. If it has two (or more) non‑parallel axes of symmetry, the centre of mass is at their intersection.
\end{propriete}

\begin{exemplebox}[Symmetry in common shapes]
\begin{itemize}
\item \textbf{Rectangle:} Two axes of symmetry (mid‑lines) $\Rightarrow$ centre of mass at the intersection of the diagonals (geometric centre).
\item \textbf{Circle / Disc:} Infinitely many axes through the centre $\Rightarrow$ centre of mass at the geometric centre.
\item \textbf{Equilateral triangle:} Three axes of symmetry (medians) $\Rightarrow$ centre of mass at the centroid, which is the intersection of the medians, located at a distance $\frac{h}{3}$ from the base (or $\frac{2h}{3}$ from the vertex).
\item \textbf{Regular polygon:} Centre of mass at the geometric centre (intersection of symmetry axes).
\item \textbf{L‑shape:} No single axis of symmetry for the whole shape $\Rightarrow$ must use composite method.
\end{itemize}
\end{exemplebox}

\begin{popfigure}
\begin{tikzpicture}[scale=0.8]
% Rectangle
\begin{scope}[shift={(0,0)}]
\draw[popblue, thick] (0,0) rectangle (3,2);
\draw[popmuted, dashed] (1.5,0) -- (1.5,2);
\draw[popmuted, dashed] (0,1) -- (3,1);
\fill[popblue] (1.5,1) circle (2pt) node[above right] {G};
\node at (1.5,-0.5) {Rectangle};
\end{scope}
% Circle
\begin{scope}[shift={(5,0)}]
\draw[poporange, thick] (0,0) circle (1.5);
\draw[popmuted, dashed] (-1.5,0) -- (1.5,0);
\draw[popmuted, dashed] (0,-1.5) -- (0,1.5);
\draw[popmuted, dashed] (-1.06,-1.06) -- (1.06,1.06);
\draw[popmuted, dashed] (-1.06,1.06) -- (1.06,-1.06);
\fill[poporange] (0,0) circle (2pt) node[above right] {G};
\node at (0,-2) {Circle};
\end{scope}
% Equilateral triangle
\begin{scope}[shift={(10,0)}]
\coordinate (A) at (0,0);
\coordinate (B) at (3,0);
\coordinate (C) at (1.5,2.6);
\draw[popgreen, thick] (A)--(B)--(C)--cycle;
\draw[popmuted, dashed] (A) -- (1.5,1.3);
\draw[popmuted, dashed] (B) -- (0.75,1.3);
\draw[popmuted, dashed] (C) -- (1.5,0);
\fill[popgreen] (1.5,0.87) circle (2pt) node[above right] {G};
\node at (1.5,-0.5) {Equilateral triangle};
\end{scope}
% L-shape
\begin{scope}[shift={(0,-4)}]
\draw[poppurple, thick] (0,0) -- (3,0) -- (3,1) -- (1,1) -- (1,3) -- (0,3) -- cycle;
\fill[poppurple] (1.2,1.2) circle (2pt) node[above right] {G?};
\node at (1.5,-0.5) {L-shape (no global symmetry)};
\end{scope}
\end{tikzpicture}
\legende{Axes of symmetry (dashed) and centre of mass G for standard laminae. The L‑shape has no global axis of symmetry.}
\end{popfigure}

\courssub{Composite Laminae by Addition}

When a lamina consists of several standard shapes joined together, we use the \emph{addition method}.

\begin{methode}[Addition Method for Composite Laminae]
\begin{enumerate}
\item \textbf{Choose a convenient origin} and axes (usually along straight edges of the lamina).
\item \textbf{Divide} the lamina into simple parts (rectangles, triangles, circles, etc.) whose centres of mass are known.
\item \textbf{Make a table} with columns: Part, Area $A_i$, $\bar{x}_i$, $\bar{y}_i$, $A_i\bar{x}_i$, $A_i\bar{y}_i$.
\item \textbf{Compute totals} $\sum A_i$, $\sum A_i\bar{x}_i$, $\sum A_i\bar{y}_i$.
\item \textbf{Apply formulas:}
\[
\bar{x} = \frac{\sum A_i\bar{x}_i}{\sum A_i}, \qquad
\bar{y} = \frac{\sum A_i\bar{y}_i}{\sum A_i}.
\]
\item \textbf{Check} that the result lies inside the lamina (or on its boundary for concave shapes).
\end{enumerate}
\end{methode}

\begin{attention}[Common Mistake: Inconsistent Coordinates]
All $\bar{x}_i$ and $\bar{y}_i$ must be measured from the \emph{same origin}. Do not mix coordinates taken from different reference points for different parts.
\end{attention}

\courssub{Composite Laminae by Subtraction (Negative Mass Method)}

If a lamina has a hole or a cut‑out, we can treat the missing part as a \emph{negative area}. This is often simpler than splitting the remaining shape into many pieces.

\begin{methode}[Negative Mass (Subtraction) Method]
\begin{enumerate}
\item Consider the lamina as a \textbf{complete solid shape} (without the hole) plus the \textbf{hole} taken with \textbf{negative area}.
\item In the table, enter the area of the hole as $-A_{\text{hole}}$.
\item The coordinates of the hole's centre of mass are its geometric centre (measured from the same origin).
\item Apply the same formulas; the negative area will automatically shift the overall centre of mass away from the hole.
\end{enumerate}
\end{methode}

\begin{aretenir}[Key Idea]
\[
\bar{x} = \frac{A_{\text{whole}}\bar{x}_{\text{whole}} + (-A_{\text{hole}})\bar{x}_{\text{hole}}}{A_{\text{whole}} - A_{\text{hole}}}
\]
The same holds for $\bar{y}$.
\end{aretenir}

\courssub{Worked Examples}

\begin{exoresolu}[L‑shaped Lamina]
A uniform lamina is made by joining a $6\,\text{cm} \times 4\,\text{cm}$ rectangle to a $4\,\text{cm} \times 2\,\text{cm}$ rectangle to form an L‑shape as shown. Find the distance of the centre of mass from the edges $AB$ and $AD$.

\begin{popfigure}
\begin{tikzpicture}[scale=0.7]
% L-shape: horizontal rectangle 6x4 (0,0)-(6,4), vertical rectangle 4x2 (0,4)-(4,6)
\draw[popblue, thick] (0,0) -- (6,0) -- (6,4) -- (4,4) -- (4,6) -- (0,6) -- cycle;
\draw[popmuted, dashed] (4,0) -- (4,4);
\draw[popmuted, dashed] (0,4) -- (4,4);
% Dimensions
\node at (3,-0.4) {6 cm};
\node at (-0.4,3) {4 cm};
\node at (5,2) {4 cm};
\node at (2,5) {2 cm};
\node at (4.4,5) {2 cm};
\node at (2,4.4) {4 cm};
% Labels
\node at (-0.3,6.3) {A};
\node at (6.3,0) {B};
\node at (0,-0.3) {D};
\node at (6.3,4) {C};
% Centres of parts
\fill[popblue] (3,2) circle (2pt) node[below left] {$G_1$};
\fill[popblue] (2,5) circle (2pt) node[below left] {$G_2$};
\fill[poporange] (2.75,2.75) circle (3pt) node[above right] {G};
\end{tikzpicture}
\legende{L‑shaped lamina split into two rectangles: horizontal $6\times4$ and vertical $4\times2$.}
\end{popfigure}

\textbf{Solution:}
Choose origin at $A$, $x$-axis along $AB$, $y$-axis along $AD$.

\begin{center}
\begin{tabular}{|c|c|c|c|c|c|}\hline
Part & Area $A_i$ (cm$^2$) & $\bar{x}_i$ (cm) & $\bar{y}_i$ (cm) & $A_i\bar{x}_i$ & $A_i\bar{y}_i$ \\ \hline
Rectangle 1 (horizontal) & $6\times4 = 24$ & $3$ & $2$ & $72$ & $48$ \\
Rectangle 2 (vertical) & $4\times2 = 8$ & $2$ & $5$ & $16$ & $40$ \\ \hline
Total & $32$ & & & $88$ & $88$ \\ \hline
\end{tabular}
\end{center}

\[
\bar{x} = \frac{88}{32} = 2.75\,\text{cm}, \qquad
\bar{y} = \frac{88}{32} = 2.75\,\text{cm}.
\]
The centre of mass is $2.75\,\text{cm}$ from $AD$ and $2.75\,\text{cm}$ from $AB$. It lies on the line $y=x$ because the chosen L‑shape is symmetric about that line.
\tcblower
\textbf{Detailed solution:}
The lamina is divided into two rectangles. Rectangle 1 (horizontal) has centre at $(3,2)$; Rectangle 2 (vertical) has centre at $(2,5)$. The table sums the first moments. The final coordinates are obtained by dividing by total area. The result lies inside the L‑shape, as expected.
\end{exoresolu}

\begin{exoresolu}[Rectangular Lamina with a Circular Hole]
A uniform rectangular lamina $ABCD$ has $AB = 10\,\text{cm}$, $BC = 6\,\text{cm}$. A circular hole of radius $1\,\text{cm}$ is cut out with its centre at the point $4\,\text{cm}$ from $AD$ and $3\,\text{cm}$ from $AB$. Find the centre of mass of the remaining lamina.

\begin{popfigure}
\begin{tikzpicture}[scale=0.6]
\draw[popblue, thick] (0,0) rectangle (10,6);
\draw[poporange, thick, fill=white] (4,3) circle (1);
\draw[popmuted, dashed] (4,0) -- (4,6);
\draw[popmuted, dashed] (0,3) -- (10,3);
% Dimensions
\node at (5,-0.5) {10 cm};
\node at (-0.5,3) {6 cm};
\node at (4,1.5) {1 cm};
% Labels
\node at (-0.3,6.3) {A};
\node at (10.3,0) {B};
\node at (0,-0.3) {D};
\node at (10.3,6) {C};
% Centres
\fill[popblue] (5,3) circle (2pt) node[above left] {$G_{\text{rect}}$};
\fill[poporange] (4,3) circle (2pt) node[below right] {$G_{\text{hole}}$};
\fill[popgreen] (5.07,3) circle (3pt) node[above] {$G$};
\end{tikzpicture}
\legende{Rectangle with a circular hole. The hole is treated as negative area.}
\end{popfigure}

\textbf{Solution:}
Origin at $A$, $x$ along $AB$, $y$ along $AD$.
\begin{itemize}
\item Whole rectangle: $A_1 = 60\,\text{cm}^2$, centre $(5,3)$.
\item Hole (circle): $A_2 = -\pi(1)^2 = -\pi\,\text{cm}^2$, centre $(4,3)$.
\end{itemize}

\begin{center}
\begin{tabular}{|c|c|c|c|c|c|}\hline
Part & Area $A_i$ & $\bar{x}_i$ & $\bar{y}_i$ & $A_i\bar{x}_i$ & $A_i\bar{y}_i$ \\ \hline
Rectangle & $60$ & $5$ & $3$ & $300$ & $180$ \\
Hole & $-\pi$ & $4$ & $3$ & $-4\pi$ & $-3\pi$ \\ \hline
Total & $60-\pi$ & & & $300-4\pi$ & $180-3\pi$ \\ \hline
\end{tabular}
\end{center}

\[
\bar{x} = \frac{300-4\pi}{60-\pi} \approx \frac{300-12.57}{60-3.14} = \frac{287.43}{56.86} \approx 5.06\,\text{cm},
\]
\[
\bar{y} = \frac{180-3\pi}{60-\pi} = \frac{3(60-\pi)}{60-\pi} = 3\,\text{cm}.
\]
The centre of mass is at $(5.06\,\text{cm}, 3\,\text{cm})$ from $A$. It lies on the horizontal midline ($y=3$) because the hole is centred on that line, and it is shifted slightly to the right (away from the hole).
\tcblower
\textbf{Detailed solution:}
The negative mass method avoids splitting the remaining shape into awkward pieces. The hole's area is subtracted. Note that $\bar{y}$ remains exactly $3$ because the hole is symmetrically placed about the horizontal midline. The $x$-coordinate moves toward the larger remaining area (right side).
\end{exoresolu}

\begin{exoresolu}[Triangle Joined to a Rectangle]
A uniform lamina consists of a rectangle $8\,\text{cm} \times 4\,\text{cm}$ with a right‑angled triangle of base $4\,\text{cm}$ and height $4\,\text{cm}$ attached along one $4\,\text{cm}$ side. Find the centre of mass relative to the rectangle's corner.

\begin{popfigure}
\begin{tikzpicture}[scale=0.7]
\draw[popblue, thick] (0,0) rectangle (8,4);
\draw[popblue, thick] (8,0) -- (12,0) -- (8,4) -- cycle;
\draw[popmuted, dashed] (8,0) -- (8,4);
% Dimensions
\node at (4,-0.4) {8 cm};
\node at (-0.4,2) {4 cm};
\node at (10,-0.4) {4 cm};
\node at (8.5,2) {4 cm};
% Centres
\fill[popblue] (4,2) circle (2pt) node[below] {$G_1$};
\fill[popblue] (28/3, 4/3) circle (2pt) node[above right] {$G_2$};
\fill[poporange] (5.2, 1.8) circle (3pt) node[above] {$G$};
\end{tikzpicture}
\legende{Rectangle with attached right triangle.}
\end{popfigure}

\textbf{Solution:}
Origin at bottom‑left corner of rectangle. $x$ to the right, $y$ upward.
\begin{itemize}
\item Rectangle: $A_1 = 32$, centre $(4,2)$.
\item Triangle (right angle at $(8,0)$): centroid is at $\left(8+\frac{4}{3}, \frac{4}{3}\right) = \left(\frac{28}{3}, \frac{4}{3}\right)$. Area $A_2 = \frac{1}{2}\times4\times4 = 8$.
\end{itemize}

\begin{center}
\begin{tabular}{|c|c|c|c|c|c|}\hline
Part & $A_i$ & $\bar{x}_i$ & $\bar{y}_i$ & $A_i\bar{x}_i$ & $A_i\bar{y}_i$ \\ \hline
Rectangle & $32$ & $4$ & $2$ & $128$ & $64$ \\
Triangle & $8$ & $28/3$ & $4/3$ & $224/3$ & $32/3$ \\ \hline
Total & $40$ & & & $608/3$ & $224/3$ \\ \hline
\end{tabular}
\end{center}

\[
\bar{x} = \frac{608/3}{40} = \frac{608}{120} = \frac{76}{15} \approx 5.07\,\text{cm},
\quad
\bar{y} = \frac{224/3}{40} = \frac{224}{120} = \frac{28}{15} \approx 1.87\,\text{cm}.
\]
\tcblower
\textbf{Detailed solution:}
The triangle's centroid is at one‑third of the height from the base and one‑third of the base from the right angle. The table uses exact fractions to avoid rounding errors. The final centre of mass lies inside the composite shape, closer to the rectangle because it has larger area.
\end{exoresolu}

\courssub{Applications: Predicting the Balance Point}

The centre of mass of a lamina is the point where it will balance horizontally on a sharp support. This has direct practical uses.

\begin{experience}[Balancing a Signboard]
A rectangular metal signboard $120\,\text{cm} \times 80\,\text{cm}$ has a circular logo (hole) of radius $10\,\text{cm}$ cut out near the top right corner. The sign is to be hung from a single hook on a wall. Where should the hook be placed so that the sign hangs with its plane vertical and the bottom edge horizontal?
\textbf{Answer:} The hook must be placed vertically above the centre of mass of the lamina. Using the negative mass method, compute $(\bar{x},\bar{y})$ relative to the bottom‑left corner. The hook position on the wall is then directly above that point. This ensures no net torque about the hook.
\end{experience}

\begin{savaistu}[Centre of Mass in Engineering]
In Cameroon, artisans who make metal gates, signboards, or decorative panels often need to find the balance point to design the mounting brackets. The composite lamina method allows them to calculate it exactly from the design drawings, without trial and error.
\end{savaistu}

\courssub{Summary of Key Points}

\begin{aretenir}[Summary: Uniform Laminae and Composite Bodies]
\begin{itemize}
\item A \textbf{uniform lamina} has constant surface density $\sigma$; mass $\propto$ area.
\item \textbf{Symmetry:} Centre of mass lies on every axis of symmetry. Intersection of axes gives the centre directly.
\item \textbf{Addition method:} Split into known shapes $\rightarrow$ table of areas and moments $\rightarrow$ $\bar{x} = \frac{\sum A_i\bar{x}_i}{\sum A_i}$.
\item \textbf{Subtraction (negative mass) method:} Treat holes as negative areas in the same table.
\item \textbf{Always} use a common origin for all coordinates.
\item The centre of mass is the \textbf{balance point} for a horizontal lamina supported at a single point.
\end{itemize}
\end{aretenir}

\courssub{Exercises}

\begin{exercice}[Practice Problems]
\begin{enumerate}
\item A uniform lamina consists of a square of side $10\,\text{cm}$ with an equilateral triangle of side $10\,\text{cm}$ attached to one side. Find the distance of the centre of mass from the base of the square.
\item A rectangular lamina $20\,\text{cm} \times 12\,\text{cm}$ has a semicircular hole of radius $3\,\text{cm}$ cut out with its diameter on the longer side, centred on that side. Find the centre of mass of the remaining lamina.
\item An L‑shaped lamina is made from three rectangles: $8\,\text{cm}\times2\,\text{cm}$, $2\,\text{cm}\times6\,\text{cm}$, and $4\,\text{cm}\times2\,\text{cm}$, arranged to form a symmetrical L (the lamina is symmetric about the line $y=x$). Use symmetry to simplify the calculation of the centre of mass.
\item A uniform lamina is bounded by the $x$-axis, the lines $x=0$, $x=4$, and the curve $y = x^2$ for $0\le x\le 2$, and $y=4$ for $2\le x\le 4$. Explain why the composite method cannot be used directly and what would be needed.
\end{enumerate}
\end{exercice}

\courssub{Corrections}

\begin{corrige}[Exercise 1]
Place the square with its base on the $x$-axis from $x=0$ to $10$, and the triangle attached on top (base from $x=0$ to $10$ at $y=10$). 
Square: area $A_1 = 100\,\text{cm}^2$, centre $(5,5)$.
Equilateral triangle: side $a=10$, height $h = \frac{\sqrt{3}}{2}a = 5\sqrt{3}\,\text{cm}$, area $A_2 = \frac{\sqrt{3}}{4}a^2 = 25\sqrt{3}\,\text{cm}^2$. Its centroid is at distance $\frac{h}{3} = \frac{5\sqrt{3}}{3}$ from its base, so $\bar{y}_2 = 10 + \frac{5\sqrt{3}}{3}$, $\bar{x}_2 = 5$ (by symmetry).
Moments about $x$-axis: $A_1\bar{y}_1 = 100\times5 = 500$, $A_2\bar{y}_2 = 25\sqrt{3}\left(10+\frac{5\sqrt{3}}{3}\right) = 250\sqrt{3} + 125$.
Total area $= 100+25\sqrt{3}$.
\[
\bar{y} = \frac{500 + 250\sqrt{3} + 125}{100+25\sqrt{3}} = \frac{625+250\sqrt{3}}{25(4+\sqrt{3})} = \frac{25+10\sqrt{3}}{4+\sqrt{3}}.
\]
Rationalising: $\bar{y} = \frac{(25+10\sqrt{3})(4-\sqrt{3})}{13} = \frac{100-25\sqrt{3}+40\sqrt{3}-30}{13} = \frac{70+15\sqrt{3}}{13} \approx 7.38\,\text{cm}$.
The centre of mass is $\frac{70+15\sqrt{3}}{13}\,\text{cm}$ above the base of the square.
\end{corrige}

\begin{corrige}[Exercise 2]
Place the rectangle with its bottom left corner at the origin, width $20\,\text{cm}$ along $x$, height $12\,\text{cm}$ along $y$. The semicircular hole has its diameter on the bottom edge ($y=0$), centred at $(10,0)$, radius $3\,\text{cm}$.
Rectangle: $A_1 = 240\,\text{cm}^2$, centre $(10,6)$.
Semicircle (hole): area $A_2 = -\frac{1}{2}\pi(3)^2 = -\frac{9\pi}{2}\,\text{cm}^2$. Centroid of a semicircle is at distance $\frac{4r}{3\pi} = \frac{4}{\pi}\,\text{cm}$ from the diameter, so its centre is at $(10, \frac{4}{\pi})$.
Moments:
$x$-moments: $240\times10 + (-\frac{9\pi}{2})\times10 = 2400 - 45\pi$; total area $= 240 - \frac{9\pi}{2}$.
By symmetry, $\bar{x} = 10\,\text{cm}$.
$y$-moments: $240\times6 + (-\frac{9\pi}{2})\times\frac{4}{\pi} = 1440 - 18 = 1422$.
\[
\bar{y} = \frac{1422}{240 - \frac{9\pi}{2}} = \frac{1422}{240 - 4.5\pi} \approx \frac{1422}{225.86} \approx 6.30\,\text{cm}.
\]
The centre of mass is at $(10\,\text{cm}, 6.30\,\text{cm})$ from the bottom left corner.
\end{corrige}

\begin{corrige}[Exercise 3]
The lamina is symmetric about the line $y=x$, so its centre of mass lies on that line: $\bar{x} = \bar{y}$.
Arrange the three rectangles as follows (see figure):
\begin{itemize}
\item Rectangle A: $8\times2$, vertices $(0,0)$, $(8,0)$, $(8,2)$, $(0,2)$. Area $16$, centre $(4,1)$.
\item Rectangle B: $2\times6$, vertices $(0,2)$, $(2,2)$, $(2,8)$, $(0,8)$. Area $12$, centre $(1,5)$.
\item Rectangle C: $4\times2$, vertices $(2,6)$, $(6,6)$, $(6,8)$, $(2,8)$. Area $8$, centre $(4,7)$.
\end{itemize}
Total area $= 36\,\text{cm}^2$.
Sum of $A_i\bar{x}_i = 16\times4 + 12\times1 + 8\times4 = 64+12+32 = 108$.
\[
\bar{x} = \frac{108}{36} = 3\,\text{cm}, \quad \bar{y} = 3\,\text{cm}.
\]
The centre of mass is at $(3\,\text{cm}, 3\,\text{cm})$ from the corner of the L.
\end{corrige}

\begin{corrige}[Exercise 4]
The region is not composed of standard geometric shapes (rectangles, triangles, circles) whose centroids are known by elementary formulas. The curved boundary $y=x^2$ requires integration to find the area and first moments. This leads to the centroid formulas using integrals:
\[
\bar{x} = \frac{\int x\,y\,dx}{\int y\,dx}, \qquad \bar{y} = \frac{\int \frac{1}{2}y^2\,dx}{\int y\,dx},
\]
which are covered in the next section (Centroids of Regions Bounded by Curves).
\end{corrige}

\coursec{Applications of Centre of Mass of Laminae and Rigid Bodies}

In the previous sections we learnt how to locate the centre of mass of a uniform lamina, using symmetry for simple shapes and the composite-body formula for laminae made of several parts. Knowing \emph{where} the centre of mass is, however, is only half the story. The real power of the concept appears when we ask practical questions: at what angle will a suspended sign hang? Will a lorry topple over on a steep hill, or will it slide first? Where must a support be placed so that a loaded plank balances horizontally? All of these questions are answered by one simple idea: \textbf{a body behaves, for the purposes of weight and balance, as if its whole mass were concentrated at its centre of mass}. This section develops that idea into a complete problem-solving toolkit.

\courssub{Freely Suspended Laminae}

\textbf{Intuition.}\par
Take a piece of cardboard cut into an irregular shape, tie a string to one corner, and let it hang freely. The cardboard swings for a moment and then settles in a definite position. If you now hang a plumb line from the same point, you will find that the plumb line always passes through the centre of mass of the cardboard, no matter which point you suspend it from. This is not a coincidence: it is a direct consequence of the equilibrium of moments.

\begin{propriete}[Freely Suspended Body]
When a rigid body is freely suspended from a point $P$ and rests in equilibrium, its \cle{centre of mass} $G$ lies on the \cle{vertical line} through $P$, directly \textbf{below} $P$.
\end{propriete}

\textbf{Justification.} The body is in equilibrium under exactly two forces: the tension $T$ in the string (or the reaction at the pivot), acting at $P$ along the vertical, and the total weight $W = Mg$, acting at $G$. Two forces can only keep a body in equilibrium if they are equal in magnitude, opposite in direction, and \emph{collinear} (they act along the same straight line); otherwise they would form a couple and the body would rotate. Hence $G$ must lie on the vertical line through $P$. Since the body hangs rather than floats, $G$ is below $P$. $\blacksquare$

\begin{experience}[Locating the Centre of Mass Experimentally]
Cut an irregular lamina from cardboard. Suspend it from a pin at a point $A$ near its edge and draw the vertical line through $A$ (use a plumb line). Repeat from a second point $B$. The two lines meet at the centre of mass $G$. Check by balancing the lamina horizontally on the tip of a pencil placed at $G$.
\end{experience}

\courssub{Angle of Inclination of a Suspended Composite Lamina}

Suppose a composite lamina is suspended from a corner $P$. In equilibrium, the line $PG$ is vertical. If we know the coordinates of $G$ relative to $P$ (found by the composite-lamina method of the previous section), we can compute the angle that any edge of the lamina makes with the vertical or with the horizontal, using simple trigonometry.

\begin{methode}[Suspension Problems]
\begin{enumerate}
\item Choose axes with origin at the suspension point $P$ (or any convenient point), and compute the coordinates $(\bar{x}, \bar{y})$ of the centre of mass $G$ of the composite lamina using $\bar{x} = \dfrac{\sum m_i x_i}{\sum m_i}$, $\bar{y} = \dfrac{\sum m_i y_i}{\sum m_i}$.
\item In equilibrium, $PG$ is vertical. Draw the right-angled triangle formed by $P$, $G$, and the projection of $G$ onto the horizontal or vertical through $P$.
\item If $\theta$ is the angle between an edge of the lamina and the vertical, read off $\tan\theta = \dfrac{\text{horizontal displacement of } G \text{ from } P}{\text{vertical displacement of } G \text{ from } P}$ and solve for $\theta$.
\end{enumerate}
\end{methode}

\begin{exoresolu}[A Rectangular Lamina Suspended from a Corner]
A uniform rectangular lamina $ABCD$ has $AB = 30\ \mathrm{cm}$ and $BC = 20\ \mathrm{cm}$. It is freely suspended from the corner $A$, with $AB$ horizontal when the lamina is held flat. Find the angle that the side $AB$ makes with the vertical in the equilibrium position.
\tcblower
\textbf{Solution.} Place $A$ at the origin, with $AB$ along the $x$-axis and $AD$ along the (downward) $y$-axis. Since the lamina is uniform, $G$ is at the centre of the rectangle:
\[ G = \left(\tfrac{30}{2}, \tfrac{20}{2}\right) = (15, 10)\ \mathrm{cm}. \]
In equilibrium, $AG$ is vertical. The vector from $A$ to $G$ has horizontal component $15$ cm and vertical component $10$ cm relative to the lamina. If $\theta$ is the angle between $AB$ and the vertical line $AG$, then in the right triangle formed by $A$, $G$ and the foot of the perpendicular from $G$ to $AB$:
\[ \tan\theta = \frac{\text{distance from } G \text{ to } AB}{\text{distance from } A \text{ along } AB} = \frac{10}{15} = \frac{2}{3}. \]
Hence $\theta = \tan^{-1}\!\left(\tfrac{2}{3}\right) \approx 33.7^{\circ}$. The side $AB$ makes an angle of about $33.7^{\circ}$ with the vertical (and therefore $90^{\circ} - 33.7^{\circ} = 56.3^{\circ}$ with the horizontal).
\end{exoresolu}

\begin{popfigure}
\begin{center}
\begin{tikzpicture}[scale=0.9, >=stealth]
% Suspension point A at origin
\coordinate (A) at (0,0);
% Vertical downward direction
\draw[dashed, popmuted, thick] (A) -- (0,-4.5) node[below] {Vertical};
% String
\draw[thick, poporange] (A) -- (0,1.2);
\fill[poporange] (A) circle (0.09);
\node[poporange, above] at (A) {$P$};
% Lamina in equilibrium: diagonal AC vertical.
% Rectangle dimensions scaled: 30cm -> 3, 20cm -> 2. Diagonal length sqrt(13)~3.606.
% Angle between AB and vertical: theta = arctan(2/3) ~ 33.69 deg.
% Let theta = 33.69. sin(theta)=2/sqrt(13), cos(theta)=3/sqrt(13).
% AB direction: (sin theta, -cos theta) = (2/√13, -3/√13)
% AD direction: perpendicular to AB, pointing down-left: (-cos theta, -sin theta) = (-3/√13, -2/√13)
% Coordinates (scaled by factor 1 for drawing):
\pgfmathsetmacro{\s}{3} % scale factor for drawing
\pgfmathsetmacro{\L}{sqrt(13)} % diagonal length in cm scaled
\pgfmathsetmacro{\sinth}{2/sqrt(13)}
\pgfmathsetmacro{\costh}{3/sqrt(13)}
% Corners:
% A = (0,0)
% B = 3 * (sinθ, -cosθ) * scale? Actually AB=30cm -> 3 units. So B = 3*(sinθ, -cosθ)
% D = 2 * (-cosθ, -sinθ) (since AD=20cm -> 2 units)
% C = B + D
\pgfmathsetmacro{\Bx}{3*\sinth}
\pgfmathsetmacro{\By}{-3*\costh}
\pgfmathsetmacro{\Dx}{-2*\costh}
\pgfmathsetmacro{\Dy}{-2*\sinth}
\pgfmathsetmacro{\Cx}{\Bx+\Dx}
\pgfmathsetmacro{\Cy}{\By+\Dy}
% G = midpoint of AC = (Cx/2, Cy/2)
\pgfmathsetmacro{\Gx}{\Cx/2}
\pgfmathsetmacro{\Gy}{\Cy/2}
% Draw lamina
\draw[thick, popblue, fill=popblueL] (0,0) -- (\Bx,\By) -- (\Cx,\Cy) -- (\Dx,\Dy) -- cycle;
% Centre of mass
\fill[popdark] (\Gx,\Gy) circle (0.07);
\node[popdark, right] at (\Gx,\Gy) {$G$};
% Labels
\node[popblue, above right] at (\Bx,\By) {$B$};
\node[popblue, below left] at (\Dx,\Dy) {$D$};
\node[popblue, below right] at (\Cx,\Cy) {$C$};
% Angle arc between vertical (down) and AB
\draw[popgreen, thick] (0,-1.2) arc (-90:-90+33.69:1.2);
\node[popgreen] at (0.7,-1.0) {$\theta$};
% Vertical line through A and G (already drawn dashed)
% Show AG vertical
\draw[popgreen, thick, ->] (0,-0.2) -- (0,-1.0);
\legende{A lamina suspended from corner $A$: in equilibrium the centre of mass $G$ hangs vertically below $A$. The diagonal $AC$ is vertical; side $AB$ makes angle $\theta = \tan^{-1}(2/3) \approx 33.7^{\circ}$ with the vertical.}
\end{tikzpicture}
\end{center}
\end{popfigure}

\begin{attention}[Which Angle Is Being Asked For?]
Do not confuse the angle that an edge makes with the \textbf{vertical} and the angle it makes with the \textbf{horizontal}: they are complementary (they add to $90^{\circ}$). Always draw the diagram, mark the vertical line through the suspension point, and identify clearly which angle $\theta$ the question requires before writing $\tan\theta = \ldots$.
\end{attention}

\begin{exoresolu}[Composite Lamina Suspended from a Point]
A uniform square lamina $ABCD$ of side $4\ \mathrm{cm}$ has a square of side $2\ \mathrm{cm}$ removed from one corner. The remaining lamina is freely suspended from the corner $A$ diagonally opposite the removed square. Find the angle that the sides through $A$ make with the vertical.
\tcblower
\textbf{Solution.} Place $A$ at the origin with the sides along the axes; the large square occupies $[0,4]\times[0,4]$ and the removed square occupies the far corner $[2,4]\times[2,4]$. By symmetry about the diagonal through $A$, $\bar{x} = \bar{y}$. Using areas (uniform lamina, so mass is proportional to area):
\[ \bar{x} = \frac{(16)(2) - (4)(3)}{16 - 4} = \frac{32 - 12}{12} = \frac{20}{12} = \frac{5}{3}\ \mathrm{cm}. \]
So $G = \left(\tfrac{5}{3}, \tfrac{5}{3}\right)$, which lies on the diagonal through $A$ — as symmetry demands. In equilibrium this diagonal is vertical, and the diagonal of a square bisects the right angle at $A$. Hence each side through $A$ makes an angle of $45^{\circ}$ with the vertical.
\end{exoresolu}

\courssub{Lamina or Block Resting on an Inclined Plane}

\textbf{Intuition.}\par
Place a tall, narrow box on a board and slowly lift one end of the board. At first nothing happens. As the slope increases, one of two things will occur: either the box suddenly \emph{topples} over its lower edge, or it \emph{slides} down the slope. Which happens first depends on two things: how high the centre of mass is (compared with the width of the base), and how rough the contact surface is.

\begin{propriete}[Condition for Equilibrium on a Plane]
A body resting on a plane remains in equilibrium (it does not topple) if and only if the \cle{vertical line} through its centre of mass $G$ falls \textbf{within the base of contact} with the plane. \emph{Toppling begins} exactly when the vertical through $G$ passes through the \textbf{edge} of the base.
\end{propriete}

\textbf{Justification.} The plane can only push on the body at points of the base of contact. The resultant reaction of the plane therefore acts somewhere within the base. For equilibrium, this reaction and the weight must be collinear (two-force equilibrium, as before), so the vertical through $G$ must meet the base. At the instant of toppling, the body pivots on the lower edge $E$ of the base: the reaction acts at $E$, so the vertical through $G$ passes through $E$. If the vertical through $G$ falls beyond $E$, the weight has a moment about $E$ that no reaction can balance, and the body topples. $\blacksquare$

\begin{popfigure}
\begin{center}
\begin{tikzpicture}[scale=1.0, >=stealth]
% Inclined plane at angle alpha where tan(alpha)=0.75 -> sin=0.6, cos=0.8
\pgfmathsetmacro{\alpha}{36.86989765} % arctan(0.75)
\pgfmathsetmacro{\cosa}{0.8}
\pgfmathsetmacro{\sina}{0.6}
% Block dimensions in its own coordinates: base 1.2 along plane, height 1.6 perpendicular
% Lower corner E at (0,0). Block corners in local coords: (0,0), (1.2,0), (1.2,1.6), (0,1.6)
% Rotate by alpha about E.
% Global coordinates:
% E = (0,0)
% Base end = (1.2*cosa, 1.2*sina) = (0.96, 0.72)
% Top corner above E = (-1.6*sina, 1.6*cosa) = (-0.96, 1.28)
% Far top corner = (0.96-0.96, 0.72+1.28) = (0, 2.0)
% Centre G = (0.6, 0.8) local -> global: (0.6*cosa - 0.8*sina, 0.6*sina + 0.8*cosa) = (0, 1.0)
\coordinate (E) at (0,0);
\coordinate (BaseEnd) at (0.96,0.72);
\coordinate (TopLeft) at (-0.96,1.28);
\coordinate (TopRight) at (0,2.0);
\coordinate (G) at (0,1.0);
% Inclined plane line
\draw[thick, popdark] (-1.5, -0.5) -- (2.5, 1.5);
% Block
\draw[thick, popblue, fill=popblueL] (E) -- (BaseEnd) -- (TopRight) -- (TopLeft) -- cycle;
% Centre of mass
\fill[popdark] (G) circle (0.07);
\node[popdark, right] at (G) {$G$};
% Lower edge pivot
\fill[poporange] (E) circle (0.09);
\node[poporange, below left] at (E) {$E$};
% Vertical line through G (and E)
\draw[dashed, popmuted, thick] (G) -- (0,-0.5);
% Weight arrow
\draw[->, thick, poppink] (G) -- (0,0.2) node[midway, right] {$W$};
% Angle of plane
\draw[popgreen, thick] (0.5,0) arc (0:\alpha:1.5);
\node[popgreen] at (1.2,0.3) {$\alpha$};
% Base width and height labels (optional)
\draw[popmuted, <->] (E) -- node[midway, above, sloped, popmuted] {$2b=1.2\,\mathrm{m}$} (BaseEnd);
\draw[popmuted, <->] (E) -- node[midway, left, sloped, popmuted] {$h=1.6\,\mathrm{m}$} (TopLeft);
\legende{A block at the point of toppling: the vertical through $G$ passes exactly through the lower edge $E$ of the base. Here $\tan\alpha = 2b/h = 0.75$.}
\end{tikzpicture}
\end{center}
\end{popfigure}

\textbf{The critical angle for toppling.} Consider a rectangular block of base width $2b$ (measured along the slope) and height $h$, with its centre of mass at the geometric centre, at height $\tfrac{h}{2}$ above the base and horizontal distance $b$ from each edge. When the plane is tilted at angle $\alpha$, toppling begins when the vertical through $G$ passes through the lower edge $E$. In the right triangle formed by $E$, $G$ and the point of the base directly below $G$:
\[ \tan\alpha_{\text{top}} = \frac{b}{h/2} = \frac{2b}{h}. \]
The \textbf{taller} the body (large $h$) and the \textbf{narrower} its base (small $b$), the \emph{smaller} the critical angle — tall narrow objects topple easily.

\textbf{Sliding versus toppling.} Sliding begins when the component of weight down the slope overcomes limiting friction: $W\sin\alpha = \mu W\cos\alpha$, i.e.
\[ \tan\alpha_{\text{slide}} = \mu, \]
where $\mu$ is the coefficient of friction between the body and the plane. As the plane is gradually tilted, \textbf{whichever critical angle is smaller occurs first}:
\begin{itemize}
\item if $\mu < \dfrac{2b}{h}$, the body \textbf{slides} before it can topple;
\item if $\mu > \dfrac{2b}{h}$, the body \textbf{topples} before it can slide.
\end{itemize}

\begin{exoresolu}[Sliding or Toppling?]
A rectangular crate of base $1.2\ \mathrm{m}$ (along the slope) and height $1.6\ \mathrm{m}$ stands on a rough plane which is slowly tilted. The coefficient of friction between the crate and the plane is $\mu = 0.5$. Determine whether the crate slides or topples first, and at what angle of inclination.
\tcblower
\textbf{Solution.} Here $2b = 1.2$ m so $b = 0.6$ m, and $h = 1.6$ m.
\[ \tan\alpha_{\text{top}} = \frac{2b}{h} = \frac{1.2}{1.6} = 0.75 \quad\Rightarrow\quad \alpha_{\text{top}} \approx 36.9^{\circ}. \]
\[ \tan\alpha_{\text{slide}} = \mu = 0.5 \quad\Rightarrow\quad \alpha_{\text{slide}} \approx 26.6^{\circ}. \]
Since $\alpha_{\text{slide}} < \alpha_{\text{top}}$, the crate \textbf{slides first}, when the plane reaches an inclination of about $26.6^{\circ}$. Had the crate been taller and narrower — say $h = 2.4$ m with the same base — then $\tan\alpha_{\text{top}} = 0.5 = \mu$, and sliding and toppling would occur simultaneously; any taller and it would topple first.
\end{exoresolu}

\begin{attention}[Height of the Centre of Mass Matters]
The toppling condition uses the position of the \emph{centre of mass}, not the geometric centre. For a non-uniform body (a lorry loaded with heavy goods on top, a tank full of water), $G$ may be much higher than $\tfrac{h}{2}$, which \emph{reduces} the critical angle $\alpha_{\text{top}}$ and makes toppling more likely. Always locate $G$ first.
\end{attention}

\courssub{Balancing Problems: Position of a Support}

A composite body rests horizontally on a single support (a knife-edge, a fulcrum, a finger) exactly when the support is placed \textbf{vertically below the centre of mass}, so that the weight produces no moment about the support. This is the principle of the balance scales used in every market in Cameroon.

\begin{exoresolu}[Balancing a Loaded Rod]
A uniform rod $AB$ of length $4\ \mathrm{m}$ and mass $6\ \mathrm{kg}$ carries a mass of $10\ \mathrm{kg}$ fixed at $A$ and a mass of $4\ \mathrm{kg}$ fixed at the point $C$, where $AC = 3\ \mathrm{m}$. Find the position of the fulcrum so that the rod rests horizontally in equilibrium.
\tcblower
\textbf{Solution.} Measure distances from $A$. The three masses and their positions are:
\begin{center}
\begin{tabular}{|c|c|c|c|}
\hline
\rowcolor{popblueL} Mass (kg) & $10$ & $6$ & $4$ \\
\hline Position from $A$ (m) & $0$ & $2$ & $3$ \\
\hline
\end{tabular}
\end{center}
The centre of mass is at
\[ \bar{x} = \frac{10(0) + 6(2) + 4(3)}{10 + 6 + 4} = \frac{0 + 12 + 12}{20} = \frac{24}{20} = 1.2\ \mathrm{m}. \]
The fulcrum must be placed \textbf{$1.2$ m from $A$}. Check: taking moments about the fulcrum, the $10$ kg mass creates an anticlockwise moment $10g \times 1.2 = 12g$, balanced by $6g \times 0.8 + 4g \times 1.8 = 4.8g + 7.2g = 12g$. $\checkmark$
\end{exoresolu}

\courssub{Real Bodies: Loaded Rods, Stacked Blocks and Vehicles}

The coordinate and position vector methods of the previous sections apply directly to real three-dimensional bodies: treat each component as a particle located at its own centre of mass, then apply $\vec{r}_G = \dfrac{\sum m_i \vec{r}_i}{\sum m_i}$.

\begin{exoresolu}[Stacked Blocks]
Two uniform bricks, each a cuboid $30\ \mathrm{cm} \times 15\ \mathrm{cm} \times 10\ \mathrm{cm}$, are stacked: the first lies flat on a table, and the second (identical) is placed flat on top of it but displaced $10\ \mathrm{cm}$ horizontally in the $x$-direction. Find the centre of mass of the stack, and determine whether the stack topples off the lower brick.
\tcblower
\textbf{Solution.} Let the lower brick occupy $0 \le x \le 30$, $0 \le z \le 10$ (heights in cm); its centre of mass is at $(15, 5)$ in the $(x,z)$-plane. The upper brick occupies $10 \le x \le 40$, $10 \le z \le 20$; its centre of mass is at $(25, 15)$. The bricks have equal mass $m$, so
\[ \bar{x} = \frac{m(15) + m(25)}{2m} = 20\ \mathrm{cm}, \qquad \bar{z} = \frac{m(5) + m(15)}{2m} = 10\ \mathrm{cm}. \]
The stack rests on the table over the base $0 \le x \le 30$ of the lower brick. Since $\bar{x} = 20$ cm lies within $[0, 30]$, the vertical through $G$ falls within the base of contact with the table: \textbf{the stack does not topple off the table}. For the upper brick alone, its centre of mass at $x = 25$ cm lies within the top face of the lower brick ($0 \le x \le 30$), so the upper brick does not topple off the lower one either. The stack is stable — but displacing the upper brick by more than $15$ cm would put its centre of mass beyond the edge $x = 30$, and it would topple.
\end{exoresolu}

\begin{savaistu}[Stability on Cameroonian Roads]
An overloaded lorry carrying a high stack of goods has a centre of mass that is both \emph{high} and possibly \emph{off-centre}. On a bend, the effective "vertical" through $G$ (combining weight and the outward centrifugal effect) tilts sideways; if it passes outside the wheelbase, the lorry topples. This is why road safety authorities insist that loads be kept low and evenly distributed. The same physics explains why a \emph{bendskin} (motorcycle taxi) passenger leaning into a curve helps keep the combined centre of mass over the line of the wheels.
\end{savaistu}

\begin{savaistu}[Head-Carrying and Water Tank Stands]
A trader carrying a basin on her head walks upright with a straight neck: this keeps the combined centre of mass of her body and the load directly above her feet, so no muscular moment is needed to prevent toppling. Similarly, the elevated water tanks seen in many Cameroonian compounds stand on four widely spaced legs: the wide base ensures that the vertical through the centre of mass of the full tank stays well inside the polygon of the legs, even when the ground vibrates or the tank is slightly off-level.
\end{savaistu}

\begin{aretenir}[Key Points of This Section]
\begin{itemize}
\item A freely suspended body rests with its centre of mass \textbf{vertically below} the point of suspension; find $G$ first, then use $\tan\theta$ in the right triangle formed with the vertical.
\item A body on a plane is in equilibrium only while the vertical through $G$ falls \textbf{within the base of contact}; toppling begins when it passes through the edge.
\item Sliding begins when $\tan\alpha = \mu$; toppling begins when $\tan\alpha = \dfrac{\text{half-base}}{\text{height of }G}$. The \emph{smaller} critical angle decides what happens first.
\item A body balances horizontally on a support placed directly beneath its centre of mass.
\item Lowering the centre of mass and widening the base are the two universal recipes for stability.
\end{itemize}
\end{aretenir}

\begin{exercice}[Suspension and Toppling]
\begin{enumerate}
\item A uniform lamina consists of a square $ABCD$ of side $12\ \mathrm{cm}$ with a circular disc of radius $3\ \mathrm{cm}$ removed, the centre of the disc being at the centre of the square. The lamina is suspended from $A$. Show that it hangs with its diagonal $AC$ vertical.
\item A uniform rectangular block has base $0.8\ \mathrm{m}$ and height $1.5\ \mathrm{m}$. It stands on a plane with $\mu = 0.6$ which is slowly tilted. Determine whether it slides or topples first, and at what angle.
\item A uniform plank of mass $8\ \mathrm{kg}$ and length $5\ \mathrm{m}$ carries a $20\ \mathrm{kg}$ sack of cement $1\ \mathrm{m}$ from one end. Where must a single support be placed for the plank to balance horizontally?
\end{enumerate}
\end{exercice}

\begin{corrige}[Exercise 1]
\textbf{1.} The removed disc is centred at the centre of the square, so by symmetry the centre of mass of the remaining lamina is still at the centre of the square, which lies on the diagonal $AC$. In equilibrium the line from the suspension point $A$ to $G$ is vertical; hence the diagonal $AC$ is vertical.

\textbf{2.} Toppling: $\tan\alpha_{\text{top}} = \dfrac{0.8}{1.5} \approx 0.533$, so $\alpha_{\text{top}} \approx 28.1^{\circ}$. Sliding: $\tan\alpha_{\text{slide}} = 0.6$, so $\alpha_{\text{slide}} \approx 31.0^{\circ}$. Since $\alpha_{\text{top}} < \alpha_{\text{slide}}$, the block \textbf{topples first} at about $28.1^{\circ}$.

\textbf{3.} Measure from the end near the sack. Masses: $20$ kg at $x = 1$ m, $8$ kg (plank) at $x = 2.5$ m. Then $\bar{x} = \dfrac{20(1) + 8(2.5)}{28} = \dfrac{40}{28} \approx 1.43$ m. Place the support $1.43$ m from the loaded end.
\end{corrige}

\coursec{Centroids of Regions Bounded by Curves}

In the previous section we located the centre of mass of laminae by decomposing them into standard shapes (rectangles, triangles, semicircles) whose centres are known. But what happens when the boundary of the lamina is a \emph{curve} — a parabola, a sine wave, an arc — that cannot be built from a finite number of standard pieces? Imagine a thin metal plate cut in the shape of the region under the curve $y = x^2$ between $x = 0$ and $x = 2$. Where should you support it so that it balances horizontally? No finite decomposition will answer this exactly. The answer comes from \cle{integration}: we slice the region into infinitely many infinitely thin strips, take moments strip by strip, and let the sum become an integral. This is one of the most beautiful applications of calculus in mechanics, and it is a favourite of GCE examiners.

\courssub{From Centre of Mass to Centroid}

\begin{definition}[Centroid of a plane region]
The \cle{centroid} of a plane region is the centre of mass of a \emph{uniform lamina} having the shape of that region. Because the lamina is uniform, its surface density $\sigma$ (mass per unit area) is constant and cancels out of every calculation: the position of the centroid depends \textbf{only on the geometry} of the region, not on the material.
\end{definition}

This is why the words \emph{centroid} and \emph{centre of mass} are used interchangeably for uniform bodies. When a question says ``find the centroid of the region bounded by \ldots'', it means exactly ``find the centre of mass of a uniform lamina occupying that region''. Geometers prefer \emph{centroid}; mechanicians prefer \emph{centre of mass}; the mathematics is identical.

\begin{savaistu}[Why ``centroid''?]
The word comes from the idea of the ``centre'' of a figure in the purely geometric sense. Archimedes, more than two thousand years ago, found the centroid of a parabolic segment by a brilliant method of slicing — essentially integration before integral calculus existed! His result: the centroid of a parabolic segment lies on its axis of symmetry, at three-fifths of the way from the vertex to the base. We shall recover results of this kind below in a few lines of calculus.
\end{savaistu}

\courssub{The Thin-Strip Method: Deriving the Formulae}

Consider the region bounded by the curve $y = f(x)$ (with $f(x) \geq 0$), the $x$-axis, and the lines $x = a$ and $x = b$. We want the coordinates $(\bar{x}, \bar{y})$ of its centroid.

\textbf{Step 1: slice the region.} Cut the region into thin \emph{vertical} strips. A typical strip stands at position $x$, has width $\delta x$ and height $y = f(x)$. It is approximately a rectangle, so:
\[
\text{area of strip} \approx y\,\delta x, \qquad \text{total area } A = \int_a^b y\,dx.
\]

\textbf{Step 2: locate the strip's own centroid.} The strip is a thin rectangle, so its centroid is at its geometric centre: at horizontal distance $x$ from the $y$-axis, and at height $\dfrac{y}{2}$ above the $x$-axis.

\textbf{Step 3: take moments.} For a uniform lamina, mass is proportional to area, so we may take moments of \emph{areas} directly (the density cancels). The \cle{first moment of area} of the strip about the $y$-axis is (area) $\times$ (distance of its centroid from the $y$-axis) $= x\,y\,\delta x$. Its first moment about the $x$-axis is (area) $\times$ (height of its centroid) $= \dfrac{y}{2}\cdot y\,\delta x = \dfrac{y^2}{2}\,\delta x$.

\textbf{Step 4: integrate.} Summing over all strips and letting $\delta x \to 0$, the sums become integrals. Applying ``moment of the whole $=$ sum of the moments of the parts'', with $A\bar{x}$ and $A\bar{y}$ as the moments of the whole region, gives the fundamental result.

\begin{propriete}[Centroid of the region under a curve]
For the region bounded by $y = f(x) \geq 0$, the $x$-axis, $x = a$ and $x = b$, with area $A = \displaystyle\int_a^b y\,dx$:
\[
\bar{x} = \frac{1}{A}\int_a^b x\,y\,dx, \qquad \bar{y} = \frac{1}{2A}\int_a^b y^2\,dx.
\]
These formulae are examinable: you may be asked to \emph{derive} them from first principles using a thin strip, so learn the four steps above.
\end{propriete}

\begin{methode}[The thin-strip method]
\begin{enumerate}
\item Sketch the region and identify the bounding curve(s) and the limits $a$ and $b$.
\item Draw a typical vertical strip at position $x$, of width $\delta x$ and height $y$.
\item Write its area $y\,\delta x$ and the coordinates of its centre $\left(x, \tfrac{y}{2}\right)$.
\item Compute $A = \displaystyle\int_a^b y\,dx$.
\item Compute $\bar{x} = \dfrac{1}{A}\displaystyle\int_a^b xy\,dx$ and $\bar{y} = \dfrac{1}{2A}\displaystyle\int_a^b y^2\,dx$.
\item Check: does the answer lie \emph{inside} the region and respect its symmetry?
\end{enumerate}
\end{methode}

\begin{attention}[The strip's centroid is at height $y/2$, not $y$]
The single most common error in this topic: when computing $\bar{y}$, students take the moment of the strip as $y \cdot y\,\delta x$, using the \emph{top} of the strip. The mass of the strip is spread from height $0$ to height $y$, and its centre is halfway up, at $\dfrac{y}{2}$. Hence the moment is $\dfrac{y^2}{2}\,\delta x$ and the factor $\dfrac{1}{2}$ in the formula for $\bar{y}$. Forgetting it doubles your $\bar{y}$ — and often pushes the centroid \emph{outside} the region, which should immediately alert you.
\end{attention}

\begin{popfigure}
\begin{center}
\begin{tikzpicture}
\begin{axis}[
  width=10cm, height=7.5cm,
  axis lines=middle,
  xlabel={$x$}, ylabel={$y$},
  xmin=-0.3, xmax=2.5, ymin=-0.3, ymax=4.6,
  xtick={1,2}, ytick={1,2,3,4},
  domain=0:2, samples=80,
  clip=false
]
\addplot[draw=none, name path=curve] {x^2};
\addplot[draw=none, name path=axis] {0};
\addplot[popblueL] fill between[of=curve and axis];
\addplot[popblue, very thick] {x^2} node[pos=0.85, above left, popink] {$y=x^2$};
\draw[poporange, thick] (axis cs:1.2,0) -- (axis cs:1.2,1.44);
\draw[poporange, thick] (axis cs:1.3,0) -- (axis cs:1.3,1.69);
\draw[poporange, thick] (axis cs:1.2,1.44) -- (axis cs:1.3,1.69);
\node[poporange, font=\small] at (axis cs:1.25,-0.35) {$\delta x$};
\fill[popdark] (axis cs:1.5,1.2) circle (2.2pt) node[right, popink] {$G\left(\tfrac{3}{2},\tfrac{6}{5}\right)$};
\end{axis}
\end{tikzpicture}
\end{center}
\legende{Region under $y = x^2$ for $0 \le x \le 2$, a typical thin strip of width $\delta x$, and the centroid $G$.}
\end{popfigure}

\courssub{Worked Examples}

\begin{exoresolu}[Centroid of the region under a parabola]
Find the coordinates of the centroid of the region bounded by the curve $y = x^2$, the $x$-axis and the line $x = 2$.
\tcblower
\textbf{Solution.} Here $f(x) = x^2$, $a = 0$, $b = 2$.

\textbf{Area:}
\[
A = \int_0^2 x^2\,dx = \left[\frac{x^3}{3}\right]_0^2 = \frac{8}{3}.
\]

\textbf{Moment about the $y$-axis:}
\[
\int_0^2 x\,y\,dx = \int_0^2 x^3\,dx = \left[\frac{x^4}{4}\right]_0^2 = 4
\quad\Longrightarrow\quad
\bar{x} = \frac{4}{8/3} = \frac{3}{2}.
\]

\textbf{Moment about the $x$-axis:}
\[
\int_0^2 \frac{y^2}{2}\,dx = \frac{1}{2}\int_0^2 x^4\,dx = \frac{1}{2}\left[\frac{x^5}{5}\right]_0^2 = \frac{16}{5}
\quad\Longrightarrow\quad
\bar{y} = \frac{16/5}{8/3} = \frac{6}{5}.
\]
Hence the centroid is $G\left(\dfrac{3}{2}, \dfrac{6}{5}\right)$. \emph{Sanity check:} $\bar{x} = 1.5 > 1$, which is sensible since more area lies to the right (the parabola grows); $\bar{y} = 1.2$ is well below half the maximum height $4$, since the region is concentrated near the $x$-axis. Both coordinates lie inside the region. $\blacksquare$
\end{exoresolu}

\begin{exoresolu}[Region under a straight line: recovering the triangle]
Show, by integration, that the centroid of the region bounded by $y = \dfrac{h}{b}x$, the $x$-axis and $x = b$ is $\left(\dfrac{2b}{3}, \dfrac{h}{3}\right)$ — the known result for a triangle.
\tcblower
\textbf{Solution.} The region is a right triangle of base $b$ and height $h$.

\textbf{Area:}
\[
A = \int_0^b \frac{h}{b}x\,dx = \frac{h}{b}\cdot\frac{b^2}{2} = \frac{bh}{2}. \quad\text{(Agrees with } \tfrac{1}{2}\times\text{base}\times\text{height.)}
\]

\textbf{For $\bar{x}$:}
\[
\int_0^b x\cdot\frac{h}{b}x\,dx = \frac{h}{b}\cdot\frac{b^3}{3} = \frac{hb^2}{3}
\quad\Longrightarrow\quad
\bar{x} = \frac{hb^2/3}{bh/2} = \frac{2b}{3}.
\]

\textbf{For $\bar{y}$:}
\[
\frac{1}{2}\int_0^b \left(\frac{h}{b}x\right)^2 dx = \frac{h^2}{2b^2}\cdot\frac{b^3}{3} = \frac{h^2 b}{6}
\quad\Longrightarrow\quad
\bar{y} = \frac{h^2 b/6}{bh/2} = \frac{h}{3}.
\]
This confirms the standard result: the centroid of a triangle lies one third of the way up from the base, here at $\left(\tfrac{2b}{3}, \tfrac{h}{3}\right)$. Integration is consistent with geometry — as it must be. $\blacksquare$
\end{exoresolu}

\courssub{Centroid of the Region Between Two Curves}

Suppose the region is bounded \emph{above} by $y = f(x)$ and \emph{below} by $y = g(x)$, with $f(x) \geq g(x)$ on $[a,b]$. A vertical strip now runs from height $g(x)$ to height $f(x)$:
\begin{itemize}
\item area of strip: $\bigl(f(x) - g(x)\bigr)\delta x$;
\item height of the strip's centroid: the \emph{midpoint} of the strip, $\dfrac{f(x) + g(x)}{2}$.
\end{itemize}
The moment of the strip about the $x$-axis is therefore
\[
\frac{f(x)+g(x)}{2}\cdot\bigl(f(x)-g(x)\bigr)\,\delta x = \frac{f(x)^2 - g(x)^2}{2}\,\delta x,
\]
which explains the difference of squares in the formula below.

\begin{propriete}[Centroid of the region between two curves]
\[
A = \int_a^b \bigl(f - g\bigr)\,dx, \qquad
\bar{x} = \frac{1}{A}\int_a^b x\bigl(f - g\bigr)\,dx, \qquad
\bar{y} = \frac{1}{2A}\int_a^b \bigl(f^2 - g^2\bigr)\,dx.
\]
Note carefully that $\bar{y}$ uses the \textbf{difference of squares} $f^2 - g^2$, \emph{not} the square of the difference $(f-g)^2$.
\end{propriete}

\begin{exoresolu}[Region between a line and a parabola]
Find the centroid of the region enclosed by $y = x$ and $y = x^2$.
\tcblower
\textbf{Solution.} The curves meet where $x = x^2$, i.e. $x^2 - x = x(x-1) = 0$, so $x = 0$ and $x = 1$; on $[0,1]$, $x \geq x^2$, so $f(x) = x$ and $g(x) = x^2$.

\textbf{Area:}
\[
A = \int_0^1 (x - x^2)\,dx = \left[\frac{x^2}{2} - \frac{x^3}{3}\right]_0^1 = \frac{1}{2} - \frac{1}{3} = \frac{1}{6}.
\]

\textbf{For $\bar{x}$:}
\[
\bar{x} = \frac{1}{A}\int_0^1 x(x - x^2)\,dx = 6\left[\frac{x^3}{3} - \frac{x^4}{4}\right]_0^1 = 6\left(\frac{1}{3} - \frac{1}{4}\right) = 6\cdot\frac{1}{12} = \frac{1}{2}.
\]

\textbf{For $\bar{y}$:}
\[
\bar{y} = \frac{1}{2A}\int_0^1 \left(x^2 - x^4\right)dx = 3\left[\frac{x^3}{3} - \frac{x^5}{5}\right]_0^1 = 3\left(\frac{1}{3} - \frac{1}{5}\right) = 3\cdot\frac{2}{15} = \frac{2}{5}.
\]
The centroid is $\left(\tfrac{1}{2}, \tfrac{2}{5}\right)$.

\emph{Sanity check.} At $x = \tfrac12$, the strip extends from $y = \tfrac14$ to $y = \tfrac12$, so the point $\left(\tfrac12, \tfrac25\right)$ lies inside the region: good. Note that this region has \emph{no} axis of symmetry (interchanging $x$ and $y$ would turn the parabola $y = x^2$ into $y = \sqrt{x}$, a different curve), so neither coordinate could have been predicted by symmetry — the full computation was necessary. Lesson: only use a symmetry after checking \emph{rigorously} that it maps the region onto itself. $\blacksquare$
\end{exoresolu}

\courssub{Verifying Standard Results by Integration}

The standard results quoted for circular laminae can (and at GCE, sometimes must) be \emph{proved} by integration.

\textbf{Semicircular lamina of radius $r$.} Take the region under $y = \sqrt{r^2 - x^2}$ for $-r \le x \le r$. By symmetry about the $y$-axis, $\bar{x} = 0$ — no integral needed. The area is $A = \tfrac12\pi r^2$. For $\bar{y}$:
\[
\bar{y} = \frac{1}{2A}\int_{-r}^{r} y^2\,dx = \frac{1}{\pi r^2}\int_{-r}^{r}(r^2 - x^2)\,dx
= \frac{1}{\pi r^2}\left[r^2 x - \frac{x^3}{3}\right]_{-r}^{r}
= \frac{1}{\pi r^2}\cdot\frac{4r^3}{3} = \frac{4r}{3\pi}.
\]

\begin{aretenir}[Standard results to know and be able to prove]
\begin{itemize}
\item Semicircular lamina of radius $r$: centroid on the axis of symmetry, at distance $\dfrac{4r}{3\pi}$ from the diameter.
\item Quadrant of a circle of radius $r$: centroid at $\left(\dfrac{4r}{3\pi}, \dfrac{4r}{3\pi}\right)$ measured from the two straight edges.
\item Triangle: centroid at one third of the height from the base (intersection of the medians).
\end{itemize}
\end{aretenir}

\begin{popfigure}
\begin{center}
\begin{tikzpicture}
\begin{axis}[
  width=9cm, height=6.5cm,
  axis lines=middle,
  xlabel={$x$}, ylabel={$y$},
  xmin=-2.4, xmax=2.4, ymin=-0.4, ymax=2.6,
  xtick={-2,-1,1,2}, ytick={0.85,2},
  yticklabels={$\frac{4r}{3\pi}$,$r$},
  axis equal image,
  clip=false
]
\addplot[domain=-2:2, samples=120, draw=none, name path=semi] {sqrt(max(0,4-x^2))};
\addplot[draw=none, name path=ax] {0};
\addplot[popgreenL] fill between[of=semi and ax];
\addplot[domain=-2:2, samples=120, popgreen, very thick] {sqrt(max(0,4-x^2))};
\draw[popgreen, very thick] (axis cs:-2,0) -- (axis cs:2,0);
\fill[popdark] (axis cs:0,0.849) circle (2.2pt) node[above right, popink] {$G\left(0,\tfrac{4r}{3\pi}\right)$};
\draw[dashed, popmuted] (axis cs:0,0) -- (axis cs:0,2);
\node[popink, font=\small] at (axis cs:1.05,1.75) {$r=2$};
\end{axis}
\end{tikzpicture}
\end{center}
\legende{Semicircular lamina of radius $r$: the centroid lies on the axis of symmetry at height $\dfrac{4r}{3\pi} \approx 0.424\,r$ above the diameter.}
\end{popfigure}

\textbf{Quadrant of a circle.} For the region in the first quadrant under $y = \sqrt{r^2 - x^2}$, $0 \le x \le r$: symmetry about the line $y = x$ gives $\bar{x} = \bar{y}$, and $A = \tfrac14\pi r^2$. Then
\[
\bar{y} = \frac{1}{2A}\int_0^r (r^2 - x^2)\,dx = \frac{2}{\pi r^2}\left[r^2x - \frac{x^3}{3}\right]_0^r = \frac{2}{\pi r^2}\cdot\frac{2r^3}{3} = \frac{4r}{3\pi},
\]
so the centroid is $\left(\tfrac{4r}{3\pi}, \tfrac{4r}{3\pi}\right)$, as quoted.

\courssub{Using Symmetry to Reduce Work}

Before integrating, \emph{always} look for symmetry:
\begin{itemize}
\item If the region is symmetric about the $y$-axis, then $\bar{x} = 0$ immediately (the integrand $xy$ is an odd function of $x$, and its integral over a symmetric interval vanishes).
\item If the region is symmetric about a line $x = c$, then $\bar{x} = c$.
\item If the region is symmetric about a line $y = d$, then $\bar{y} = d$.
\end{itemize}
Symmetry can halve your work — and provides a free check on the remaining coordinate.

\begin{attention}[GCE exam technique: setting up the integrals]
Examiners report the same errors year after year:
\begin{enumerate}
\item \textbf{Wrong limits.} Find the intersections of the bounding curves \emph{first}, by solving $f(x) = g(x)$, and use those as limits.
\item \textbf{Missing factor $\tfrac12$} in $\bar{y}$ (see the warning above).
\item \textbf{Forgetting to divide by $A$}, or dividing by the wrong area.
\item \textbf{Mixing curves:} for a region between two curves, $\bar{y}$ involves $f^2 - g^2$, not $(f-g)^2$.
\item \textbf{No check:} your centroid must lie \emph{inside} the region (or on its boundary in degenerate cases). If $\bar{y}$ comes out larger than the maximum height of the region, something is wrong — go back.
\end{enumerate}
Also keep an eye on dimensions: $A$ is an area, $\int xy\,dx$ is an area $\times$ a length, so $\bar{x}$ is a length. If your $\bar{x}$ has the dimensions of an area, you have dropped an $x$ somewhere.
\end{attention}

\begin{exercice}[Practice]
\begin{enumerate}
\item Find the centroid of the region bounded by $y = 4 - x^2$ and the $x$-axis. (Use symmetry for $\bar{x}$.)
\item Find the centroid of the region bounded by $y = \sqrt{x}$, the $x$-axis and the line $x = 4$.
\item The region bounded by $y = x^2$ and $y = 2x$ is occupied by a uniform lamina. Find the coordinates of its centre of mass.
\item Prove by integration that the centroid of a quadrant of a circle of radius $r$ is at distance $\dfrac{4r}{3\pi}$ from each straight edge.
\end{enumerate}
\end{exercice}

\begin{corrige}[Exercise 1]
The region is symmetric about the $y$-axis, so $\bar{x} = 0$. The curve meets the $x$-axis at $x = \pm 2$.
$A = \displaystyle\int_{-2}^{2}(4-x^2)\,dx = \left[4x - \tfrac{x^3}{3}\right]_{-2}^{2} = \tfrac{32}{3}$.
$\bar{y} = \dfrac{1}{2A}\displaystyle\int_{-2}^{2}(4-x^2)^2\,dx = \dfrac{3}{64}\int_{-2}^{2}(16 - 8x^2 + x^4)\,dx = \dfrac{3}{64}\cdot\dfrac{512}{15} = \dfrac{8}{5}$.
Centroid: $\left(0, \tfrac{8}{5}\right)$.
\end{corrige}

\begin{corrige}[Exercise 2]
$A = \displaystyle\int_0^4 \sqrt{x}\,dx = \left[\tfrac{2}{3}x^{3/2}\right]_0^4 = \tfrac{16}{3}$.
$\bar{x} = \dfrac{3}{16}\displaystyle\int_0^4 x^{3/2}\,dx = \dfrac{3}{16}\cdot\dfrac{2}{5}\cdot 32 = \dfrac{12}{5}$.
$\bar{y} = \dfrac{1}{2A}\displaystyle\int_0^4 x\,dx = \dfrac{3}{32}\cdot 8 = \dfrac{3}{4}$.
Centroid: $\left(\tfrac{12}{5}, \tfrac{3}{4}\right)$.
\end{corrige}

\begin{corrige}[Exercise 3]
Intersections: $x^2 = 2x \Rightarrow x = 0, 2$; on $[0,2]$, $2x \ge x^2$.
$A = \displaystyle\int_0^2 (2x - x^2)\,dx = \left[x^2 - \tfrac{x^3}{3}\right]_0^2 = \tfrac{4}{3}$.
$\bar{x} = \dfrac{3}{4}\displaystyle\int_0^2 x(2x - x^2)\,dx = \dfrac{3}{4}\left[\tfrac{2x^3}{3} - \tfrac{x^4}{4}\right]_0^2 = \dfrac{3}{4}\cdot\dfrac{4}{3} = 1$.
$\bar{y} = \dfrac{1}{2A}\displaystyle\int_0^2 (4x^2 - x^4)\,dx = \dfrac{3}{8}\left[\tfrac{4x^3}{3} - \tfrac{x^5}{5}\right]_0^2 = \dfrac{3}{8}\cdot\dfrac{64}{15} = \dfrac{8}{5}$.
Centre of mass: $\left(1, \tfrac{8}{5}\right)$.
\end{corrige}

\begin{corrige}[Exercise 4]
Region: under $y = \sqrt{r^2 - x^2}$, $0 \le x \le r$; $A = \tfrac14\pi r^2$. By symmetry about the line $y = x$, $\bar{x} = \bar{y}$. Then $\bar{y} = \dfrac{1}{2A}\displaystyle\int_0^r (r^2 - x^2)\,dx = \dfrac{2}{\pi r^2}\cdot\dfrac{2r^3}{3} = \dfrac{4r}{3\pi}$. Hence the centroid is $\left(\tfrac{4r}{3\pi}, \tfrac{4r}{3\pi}\right)$, i.e. at distance $\tfrac{4r}{3\pi}$ from each straight edge.
\end{corrige}

With the thin-strip method mastered for regions bounded by curves, we are ready for the final step: combining these curved pieces with the standard shapes of earlier sections to find centroids of \emph{composite areas} — the subject of the next and last section of this chapter.

\coursec{Centroids of Composite Areas and Applications of Centroids}

In the previous section we learnt how to locate the centroid of a region bounded by curves using integration: we sliced the region into thin strips, computed the moment of each strip about the axes, and summed by integrating. That method is powerful, but in many real and examination problems the region is \emph{not} bounded by a single curve. Think of the metal plate of a road sign on the Douala--Yaound\'e highway: it may be a rectangle with a semicircular top, or a triangular panel with a circular hole drilled through it. For such shapes, integrating over the whole boundary would be clumsy. The intelligent strategy is to \textbf{cut the region into simple pieces whose centroids we already know}, and then recombine. That is the subject of this final section: \cle{composite areas}, the \cle{addition and subtraction of moments}, and the practical \cle{applications of centroids}.

\courssub{Centroid of a Composite Area: the Weighted-Average Principle}

\textbf{Intuition first.} Suppose a uniform plate is made of two rectangular parts. Each part, taken alone, would balance at its own centre. The whole plate balances at a point $G$ which lies \emph{between} the two centres, but \emph{closer to the bigger part}. The bigger a part is, the more it ``pulls'' $G$ towards itself. This is exactly the same weighted-average idea we used for systems of particles in Section 4: each part plays the role of a particle whose ``mass'' is proportional to its area (the lamina being uniform, mass and area are proportional).

\begin{propriete}[Centroid of a composite area]
Let a plane region of total area $A$ be divided into parts of areas $A_1, A_2, \dots, A_n$, with centroids $(\bar{x}_1,\bar{y}_1), (\bar{x}_2,\bar{y}_2), \dots, (\bar{x}_n,\bar{y}_n)$ respectively. Then the centroid $(\bar{x},\bar{y})$ of the whole region is
\[
\bar{x}=\frac{\sum_{i=1}^{n} A_i\,\bar{x}_i}{\sum_{i=1}^{n} A_i}=\frac{\sum A_i\bar{x}_i}{A},
\qquad
\bar{y}=\frac{\sum_{i=1}^{n} A_i\,\bar{y}_i}{\sum_{i=1}^{n} A_i}=\frac{\sum A_i\bar{y}_i}{A}.
\]
The quantities $A_i\bar{x}_i$ and $A_i\bar{y}_i$ are called the \cle{first moments} of the parts about the $y$-axis and $x$-axis respectively. This result is examinable and follows directly from the additive property of integrals: $\iint_R x\,\mathrm{d}A = \sum_i \iint_{R_i} x\,\mathrm{d}A$ when $R=\bigcup_i R_i$ with disjoint interiors.
\end{propriete}

\begin{attention}[Subtracting areas means subtracting moments]
When a region is obtained by \emph{removing} a part (a hole, a cut-out), treat the removed part as having a \textbf{negative area}. Its moments must be subtracted too:
\[
\bar{x}=\frac{A_1\bar{x}_1-A_2\bar{x}_2}{A_1-A_2},
\qquad
\bar{y}=\frac{A_1\bar{y}_1-A_2\bar{y}_2}{A_1-A_2}.
\]
A very common GCE error is to subtract the areas in the denominator but \emph{add} the moments in the numerator. Remember: \textbf{the same sign applies everywhere} --- a removed piece contributes negatively to the area, to the moment about the $x$-axis, and to the moment about the $y$-axis.
\end{attention}

\courssub{Worked Examples}

\begin{exoresolu}[Rectangle with a parabolic cap]
A uniform lamina consists of the rectangle with vertices $(0,0)$, $(4,0)$, $(4,2)$, $(0,2)$, together with the region bounded by the parabola $y=2+\dfrac{x(4-x)}{2}$ and the line $y=2$ for $0\le x\le 4$. Find the coordinates of its centroid.
\tcblower
\textbf{Solution.} By symmetry about $x=2$, we have $\bar{x}=2$ immediately.

\emph{Part 1: the rectangle.} Area $A_1=4\times 2=8$; centroid at $(2,1)$, so $\bar{y}_1=1$ and $A_1\bar{y}_1=8$.

\emph{Part 2: the parabolic cap.} Its area is
\[
A_2=\int_0^4 \frac{x(4-x)}{2}\,\mathrm{d}x
=\frac{1}{2}\int_0^4(4x-x^2)\,\mathrm{d}x
=\frac{1}{2}\left[2x^2-\frac{x^3}{3}\right]_0^4
=\frac{1}{2}\left(32-\frac{64}{3}\right)=\frac{16}{3}.
\]
For the moment of the cap about the $x$-axis, take a vertical strip at position $x$. The strip runs from $y=2$ to $y=2+\tfrac{x(4-x)}{2}$; its area is $\tfrac{x(4-x)}{2}\,\mathrm{d}x$ and its centroid is at the mid-height $2+\tfrac{x(4-x)}{4}$. Hence
\begin{align*}
A_2\bar{y}_2&=\int_0^4 \frac{x(4-x)}{2}\left(2+\frac{x(4-x)}{4}\right)\mathrm{d}x
=2\int_0^4\frac{x(4-x)}{2}\,\mathrm{d}x+\frac{1}{8}\int_0^4 x^2(4-x)^2\,\mathrm{d}x\\
&=2A_2+\frac{1}{8}\int_0^4 x^2(4-x)^2\,\mathrm{d}x.
\end{align*}
Now
\[
\int_0^4 x^2(4-x)^2\,\mathrm{d}x=\int_0^4\left(16x^2-8x^3+x^4\right)\mathrm{d}x
=\left[\frac{16x^3}{3}-2x^4+\frac{x^5}{5}\right]_0^4
=\frac{1024}{3}-512+\frac{1024}{5}=\frac{256}{15}.
\]
Therefore
\[
A_2\bar{y}_2=2\times\frac{16}{3}+\frac{1}{8}\times\frac{256}{15}
=\frac{32}{3}+\frac{32}{15}=\frac{160+32}{15}=\frac{192}{15}=\frac{64}{5}.
\]

\emph{Recombination.}
\[
\bar{y}=\frac{A_1\bar{y}_1+A_2\bar{y}_2}{A_1+A_2}
=\frac{8+\tfrac{64}{5}}{8+\tfrac{16}{3}}
=\frac{\tfrac{104}{5}}{\tfrac{40}{3}}
=\frac{104}{5}\times\frac{3}{40}=\frac{39}{25}=1.56.
\]
The centroid is $\left(2,\;1.56\right)$. Note that $1.56<2$: although the cap adds area above, the rectangle dominates, so $G$ stays inside the rectangle --- a useful sanity check.
\end{exoresolu}

\begin{popfigure}
\begin{center}
\begin{tikzpicture}
\begin{axis}[
  width=11cm, height=7.5cm,
  axis lines=middle,
  xlabel={$x$}, ylabel={$y$},
  xmin=-0.5, xmax=4.8, ymin=-0.5, ymax=4.6,
  xtick={0,1,2,3,4}, ytick={1,2,3,4},
  axis equal image,
  every axis x label/.style={at={(ticklabel* cs:1.02)}, anchor=west},
  every axis y label/.style={at={(ticklabel* cs:1.02)}, anchor=south}
]
\addplot[fill=popblueL, draw=none, domain=0:4] {2} \closedcycle;
\addplot[fill=popgreenL, draw=none, domain=0:4] {2+x*(4-x)/2} \closedcycle;
\addplot[draw=popblue, thick, domain=0:4] {2+x*(4-x)/2};
\draw[draw=popblue, thick] (axis cs:0,0) -- (axis cs:4,0) -- (axis cs:4,2);
\draw[draw=popblue, thick] (axis cs:0,0) -- (axis cs:0,2);
\addplot[only marks, mark=*, mark size=2.5pt, color=poporange] coordinates {(2,1.56)};
\node[color=poporange, anchor=west, font=\small] at (axis cs:2.15,1.45) {$G(2,\ 1.56)$};
\node[color=popdark, font=\small] at (axis cs:2,0.9) {rectangle};
\node[color=popdark, font=\small] at (axis cs:2,3.1) {parabolic cap};
\draw[dashed, gray] (axis cs:2,0) -- (axis cs:2,4.2);
\node[color=gray, font=\footnotesize, anchor=south] at (axis cs:2,4.2) {axis of symmetry};
\end{axis}
\end{tikzpicture}
\end{center}
\legende{Composite region: a rectangle topped by a parabolic cap. The axis of symmetry gives $\bar{x}=2$ at once; only $\bar{y}$ needs computation.}
\end{popfigure}

\begin{exoresolu}[A region with a part removed]
A uniform square lamina of side $6\ \mathrm{cm}$ has vertices $(0,0),(6,0),(6,6),(0,6)$. A circular disc of radius $1.5\ \mathrm{cm}$ centred at $(4,4)$ is cut out. Find the centroid of the remaining plate.
\tcblower
\textbf{Solution.} Square: $A_1=36$, centroid $(3,3)$. Removed disc: $A_2=\pi(1.5)^2=2.25\pi$, centroid $(4,4)$, taken \emph{negative}.
\[
\bar{x}=\frac{36(3)-2.25\pi(4)}{36-2.25\pi}
=\frac{108-9\pi}{36-2.25\pi}
\approx\frac{108-28.27}{36-7.07}
\approx\frac{79.73}{28.93}\approx 2.76.
\]
The configuration is symmetric in the line $y=x$ (the square and the centre of the hole both lie on it), so $\bar{y}=\bar{x}\approx 2.76\ \mathrm{cm}$. The centroid has shifted \emph{away} from the hole, from $(3,3)$ towards the opposite corner --- exactly as intuition predicts: removing mass on one side moves the balance point to the other side.
\end{exoresolu}

\courssub{Choosing a Strategy: Symmetry, Decomposition or Integration?}

Faced with an examination question, do not start integrating blindly. Follow this decision order.

\begin{methode}[Decision strategy for locating a centroid]
\begin{enumerate}
\item \textbf{Symmetry first.} Look for axes of symmetry. Every axis of symmetry contains the centroid, so each one gives one coordinate (or a relation between them) for free.
\item \textbf{Decomposition next.} Can the region be built from, or reduced by, standard shapes (rectangles, triangles, circles, semicircles, circular sectors, standard parabolic segments)? If yes, use $\bar{x}=\dfrac{\sum A_i\bar{x}_i}{\sum A_i}$, $\bar{y}=\dfrac{\sum A_i\bar{y}_i}{\sum A_i}$, subtracting moments of removed parts.
\item \textbf{Integration last.} If a boundary is a genuine curve that is not a standard shape, integrate: $\bar{x}=\dfrac{1}{A}\displaystyle\int x\,y\,\mathrm{d}x$, $\bar{y}=\dfrac{1}{2A}\displaystyle\int y^2\,\mathrm{d}x$ for a region under $y=f(x)$, and combine the result with any standard parts as in the worked example above.
\item \textbf{Check.} Verify that your answer lies inside (or at least on the convex hull of) the region and on every axis of symmetry; estimate roughly whether the position is plausible.
\end{enumerate}
\end{methode}

\begin{popfigure}
\begin{center}
\begin{tikzpicture}[
  box/.style={draw=popblue, fill=popblueL, rounded corners, text width=4.6cm, align=center, font=\small, inner sep=5pt},
  dec/.style={draw=poporange, fill=white, rounded corners, text width=4.6cm, align=center, font=\small, inner sep=5pt},
  arr/.style={-{Stealth}, thick, color=popdark}
]
\node[dec, align=center] (q) at (0,0){Region given.\\ Axis of symmetry?};
\node[box, align=center] (sym) at (6.5,0){Use symmetry:\\ $G$ lies on each axis};
\node[dec, align=center] (d) at (0,-2.6){Standard parts\\ (add or remove)?};
\node[box, align=center] (dec) at (6.5,-2.6){Weighted average:\\ $\bar{x}=\frac{\sum A_i\bar{x}_i}{\sum A_i}$};
\node[box, align=center] (int) at (0,-5.2){Integrate:\\ $\bar{x}=\frac{1}{A}\int xy\,\mathrm{d}x$,\quad $\bar{y}=\frac{1}{2A}\int y^2\,\mathrm{d}x$};
\node[box, align=center] (chk) at (6.5,-5.2){Check plausibility\\ and symmetry of $G$};
\draw[arr] (q) -- node[above, font=\footnotesize]{yes} (sym);
\draw[arr] (q) -- node[left, font=\footnotesize]{no / partial} (d);
\draw[arr] (d) -- node[above, font=\footnotesize]{yes} (dec);
\draw[arr] (d) -- node[left, font=\footnotesize]{no} (int);
\draw[arr] (sym) |- (chk);
\draw[arr] (dec) -- (chk);
\draw[arr] (int) -- (chk);
\end{tikzpicture}
\end{center}
\legende{Decision flowchart: symmetry, then decomposition into standard parts, then integration; always finish with a plausibility check.}
\end{popfigure}

\courssub{Applications of Centroids}

\textbf{Balancing plates.} A uniform plate balances on a knife edge or on a fingertip placed exactly under its centroid. A craftsman cutting a decorative metal plate in Bamenda must know where the centroid lies to drill the suspension hole: if the hole is at $G$, the plate hangs horizontally.

\textbf{Suspension of curved shapes.} When a lamina is suspended freely from a point $P$, it comes to rest with $G$ \emph{vertically below} $P$ (this is the equilibrium condition from the chapter on moments: the weight acts through $G$). Conversely, suspending a shape from two different points and drawing the vertical line each time is a practical experimental way to \emph{find} $G$: the two lines meet at the centroid.

\begin{experience}[Finding a centroid experimentally]
Cut an irregular shape from cardboard. Pin it at a point $A$ near its edge so it swings freely, hang a plumb line from $A$, and draw the line. Repeat from a second point $B$. The intersection of the two lines is the centroid $G$. Verify by balancing the cardboard on a pencil tip at $G$.
\end{experience}

\textbf{Design of stable objects.} An object topples when the vertical line through its centre of mass falls outside its base of support. Low centroids and wide bases give stability: this is why a lorry loaded with heavy goods on the steep roads of the North-West Region must keep the load low, why traditional water pots have wide bases, and why a racing car is built low to the ground. For a curved lamina standing on an edge, computing the centroid tells the designer exactly how far the shape can be tilted before it topples.

\begin{savaistu}[Pappus' theorem: from centroids to volumes (beyond the syllabus)]
If a plane region of area $A$, lying entirely on one side of an axis, is rotated through one full turn about that axis, the volume generated is
\[
V = A \times 2\pi\,\bar{d},
\]
where $\bar{d}$ is the distance of the centroid from the axis: the volume equals the area times the distance travelled by the centroid. For example, rotating a disc of radius $r$ whose centre is at distance $R>r$ from the axis gives a torus of volume $V=(\pi r^2)(2\pi R)=2\pi^2 R r^2$. This \cle{theorem of Pappus} is not required at GCE Advanced Level, but it shows beautifully why the centroid is the ``average position'' of a region, and it offers a quick check on volume computations you will meet in the chapter on volumes of revolution.
\end{savaistu}

\courssub{Chapter Synthesis: Choosing Among All the Methods}

Throughout this chapter we have built four complementary tools for locating a centre of mass or centroid. The table below gathers them and says when each is the right choice.

\begin{center}
\begin{tabularx}{\linewidth}{|l|X|X|}
\hline
\rowcolor{popblueL}
\textbf{Method} & \textbf{Key formula / idea} & \textbf{Use when} \\
\hline
Symmetry & $G$ lies on every axis of symmetry & The body or lamina has one or more axes of symmetry (rods, discs, rectangles, isosceles triangles). \\
\hline
Coordinates (particles) & $\bar{x}=\dfrac{\sum m_i x_i}{\sum m_i}$, \ $\bar{y}=\dfrac{\sum m_i y_i}{\sum m_i}$ & A system of discrete particles or point masses given by coordinates. \\
\hline
Position vectors & $\bar{\vec{r}}=\dfrac{\sum m_i\vec{r}_i}{\sum m_i}$ & Particles in 2-D or 3-D, especially when the answer is wanted as a single vector or in 3-D geometry. \\
\hline
Decomposition (composite bodies) & $\bar{x}=\dfrac{\sum A_i\bar{x}_i}{\sum A_i}$ with negative areas for holes & Uniform laminae or rigid bodies built from standard shapes, possibly with parts removed. \\
\hline
Integration & $\bar{x}=\dfrac{1}{A}\displaystyle\int x\,\mathrm{d}A$, \ $\bar{y}=\dfrac{1}{A}\displaystyle\int y\,\mathrm{d}A$ & Regions bounded by genuine curves; also the source of all standard results (semicircle, sector, triangle). \\
\hline
\end{tabularx}
\end{center}

\begin{aretenir}[What you must remember from this chapter]
\begin{itemize}
\item The centre of mass is the unique point where the total weight of a body may be considered to act; a body balances when supported at or below this point.
\item For a uniform body, mass is proportional to length, area or volume, so the centre of mass coincides with the geometric \cle{centroid}.
\item Always exploit symmetry first; it can halve the work.
\item For composite areas, moments add: $\bar{x}=\dfrac{\sum A_i\bar{x}_i}{\sum A_i}$, and a removed part contributes \emph{negative} area and \emph{negative} moments.
\item For regions bounded by curves, fall back on integration, and combine the result with standard parts when the region is mixed.
\item Sanity checks: $G$ lies on all axes of symmetry, inside the convex hull of the region, and shifts away from any hole.
\end{itemize}
\end{aretenir}

\begin{exercice}[Composite area with a triangular cut-out]
A uniform lamina is the rectangle with vertices $(0,0)$, $(8,0)$, $(8,5)$, $(0,5)$, from which the triangle with vertices $(2,0)$, $(6,0)$, $(4,3)$ has been removed. Find the coordinates of the centroid of the remaining lamina.
\end{exercice}

\begin{corrige}[Exercise 1]
Rectangle: $A_1=40$, centroid $(4,2.5)$. Triangle removed: base $4$, height $3$, so $A_2=\tfrac12(4)(3)=6$; its centroid is at the mean of the vertices, $\left(\tfrac{2+6+4}{3},\tfrac{0+0+3}{3}\right)=(4,1)$. By symmetry about $x=4$, $\bar{x}=4$.
\[
\bar{y}=\frac{40(2.5)-6(1)}{40-6}=\frac{100-6}{34}=\frac{94}{34}=\frac{47}{17}\approx 2.76.
\]
The centroid is $\left(4,\ \tfrac{47}{17}\right)$. The hole is at the bottom, so the centroid has risen above $2.5$, as expected.
\end{corrige}

\begin{exercice}[Curve plus standard shape]
A lamina consists of the square $[0,2]\times[0,2]$ together with the region bounded by $y=x^2$, the $x$-axis and the lines $x=2$, $x=3$. Find the centroid of the whole lamina. (Given: $\int_2^3 x^2\,\mathrm{d}x=\tfrac{19}{3}$, $\int_2^3 x^3\,\mathrm{d}x=\tfrac{65}{4}$, $\int_2^3 x^4\,\mathrm{d}x=\tfrac{211}{5}$.)
\end{exercice}

\begin{corrige}[Exercise 2]
Square: $A_1=4$, centroid $(1,1)$, moments $A_1\bar{x}_1=4$, $A_1\bar{y}_1=4$.

Curved part: $A_2=\int_2^3 x^2\,\mathrm{d}x=\tfrac{19}{3}$.
\[
A_2\bar{x}_2=\int_2^3 x\cdot x^2\,\mathrm{d}x=\frac{65}{4},\qquad
A_2\bar{y}_2=\frac12\int_2^3 (x^2)^2\,\mathrm{d}x=\frac{1}{2}\cdot\frac{211}{5}=\frac{211}{10}.
\]
Total area $A=4+\tfrac{19}{3}=\tfrac{31}{3}$. Hence
\[
\bar{x}=\frac{4+\tfrac{65}{4}}{\tfrac{31}{3}}=\frac{\tfrac{81}{4}}{\tfrac{31}{3}}=\frac{243}{124}\approx 1.96,
\qquad
\bar{y}=\frac{4+\tfrac{211}{10}}{\tfrac{31}{3}}=\frac{\tfrac{251}{10}}{\tfrac{31}{3}}=\frac{753}{310}\approx 2.43.
\]
The centroid is approximately $(1.96,\ 2.43)$; it lies above the square because the added region sits on the right and extends high --- a plausible result.
\end{corrige}

\newpage

\coursec{Integration activity}

\coursec{Application: Centre of Mass and Stability of a Water Tank Stand}

\courssub{Aims of the Activity}
\begin{itemize}
    \item Mobilise the concepts of centre of mass for uniform rigid bodies, composite laminae, and systems of particles in a single authentic engineering task.
    \item Apply symmetry arguments, the addition/subtraction method for composite laminae, and the coordinate formula for systems of bodies.
    \item Determine the stability of a structure by analysing the toppling condition when the ground tilts.
    \item Evaluate a design modification (widening the base or lowering the tank) and quantify its effect on stability.
    \item Develop collaborative problem‑solving, technical drawing, and communication skills through a group poster presentation.
\end{itemize}

\courssub{Situation}
The Government Technical High School in Bamenda plans to install a \SI{5000}{L} cylindrical water tank on a wooden stand to ensure water supply during dry seasons. The tank is a uniform right circular cylinder of radius $R = \SI{0.6}{m}$ and height $H = \SI{1.8}{m}$. When full, its mass is $M_{\text{tank}} = \SI{5000}{kg}$ (water density $\SI{1000}{kg/m^3}$; tank wall mass is negligible). The tank sits with its circular base on the horizontal tops of the vertical posts at the top of the stand.

The stand is modelled as a \textbf{uniform lamina} of thickness $t = \SI{0.05}{m}$ and density $\rho = \SI{600}{kg/m^3}$ (typical for dried mahogany). It consists of:
\begin{itemize}
    \item A rectangular base frame: outer dimensions $\SI{1.2}{m} \times \SI{1.2}{m}$, with a uniform width of $\SI{0.15}{m}$ (like a square picture frame).
    \item Four vertical posts at the corners, each a rectangle $\SI{0.15}{m} \times \SI{1.5}{m}$ (width $\times$ height) and thickness $t = \SI{0.05}{m}$.
    \item A triangular brace on each side: right‑angled triangles of legs $\SI{1.5}{m}$ (vertical) and $\SI{0.4}{m}$ (horizontal), attached to the outer face of the posts.
    \item A circular hole of radius $\SI{0.05}{m}$ cut through the centre of the base frame to allow a pipe to pass through.
\end{itemize}
The coordinate system is chosen with origin $O$ at the centre of the bottom of the base frame (ground level), $x$-axis along one side of the square base, $y$-axis along the adjacent side, and $z$-axis vertically upward. The tank is centred on the top platform, so its central axis coincides with the $z$-axis.

\begin{popfigure}
\begin{tikzpicture}[scale=0.7, every node/.style={font=\small}]
% Axes
\draw[->] (-1,0) -- (3,0) node[right] {$x$};
\draw[->] (0,-0.5) -- (0,4) node[above] {$z$};
% Ground
\draw[popline, thick] (-1,0) -- (3,0);
% Base frame (horizontal lamina, thickness 0.05 m)
\draw[popblue, thick, fill=popblueL] (0,0) rectangle (1.2,0.05);
% Left post (vertical lamina, width 0.15 m, height 1.5 m)
\draw[popblue, thick, fill=popblueL] (0,0.05) rectangle (0.15,1.55);
% Right post
\draw[popblue, thick, fill=popblueL] (1.05,0.05) rectangle (1.2,1.55);
% Triangular brace on left side (attached to outer face of left post, protruding leftwards)
\draw[popgreen, thick, fill=popgreenL] (0,0.05) -- (-0.4,0.05) -- (0,1.55) -- cycle;
% Triangular brace on right side (attached to outer face of right post, protruding rightwards)
\draw[popgreen, thick, fill=popgreenL] (1.2,0.05) -- (1.6,0.05) -- (1.2,1.55) -- cycle;
% Tank on top: cylinder represented as rectangle of width 1.2 (diameter) and height 1.8
\draw[poporange, thick, fill=poporange!30] (-0.3,1.55) rectangle (1.5,3.35);
% Centre of mass markers
\fill[poppink] (0.6,0.025) circle (2pt) node[below right] {$G_{\text{base}}$};
\fill[poppink] (0.075,0.8) circle (2pt) node[left] {$G_{\text{post}}$};
\fill[poppink] (-0.133,0.55) circle (2pt) node[left] {$G_{\text{brace}}$};
\fill[poppink] (0.6,2.45) circle (2pt) node[right] {$G_{\text{tank}}$};
% Dimensions
\draw[<->, popmuted] (0,-0.3) -- (1.2,-0.3) node[midway, below] {1.2 m};
\draw[<->, popmuted] (1.3,0.05) -- (1.3,1.55) node[midway, right] {1.5 m};
\draw[<->, popmuted] (1.3,1.55) -- (1.3,3.35) node[midway, right] {1.8 m};
\draw[<->, popmuted] (-0.3,1.55) -- (-0.3,3.35) node[midway, left] {1.8 m};
\draw[<->, popmuted] (-0.3,3.35) -- (1.5,3.35) node[midway, above] {1.2 m (diameter)};
\draw[<->, popmuted] (0,-0.15) -- (0,0.05) node[midway, left] {0.05 m};
\end{tikzpicture}
\legende{Side view of the water‑tank stand (along $x$-axis). The tank (orange) sits on the tops of the posts. The base frame (blue) is a horizontal lamina of thickness 0.05 m. The posts (blue) are vertical laminae of width 0.15 m and height 1.5 m. The triangular braces (green) are attached to the outer faces of the posts. The circular hole in the base is not visible in this view.}
\end{popfigure}

\courssub{Method: Centre of Mass of Composite Laminae}
\begin{methode}[Addition/Subtraction Method for Composite Laminae]
\begin{enumerate}
    \item \textbf{Decompose} the lamina into simple shapes (rectangles, triangles, circles) whose centres of mass are known.
    \item \textbf{Assign} a mass to each part: $m_i = \sigma \times A_i$, where $\sigma = \rho t$ is the surface density (mass per unit area). For holes, treat the missing material as a \emph{negative mass} ($m_i < 0$).
    \item \textbf{Locate} the centre of mass $(x_i, y_i, z_i)$ of each part using symmetry or standard formulas (e.g., rectangle at geometric centre, triangle at centroid $1/3$ from the right angle).
    \item \textbf{Compute} the overall centre of mass coordinates:
    \[
    x_G = \frac{\sum m_i x_i}{\sum m_i},\quad
    y_G = \frac{\sum m_i y_i}{\sum m_i},\quad
    z_G = \frac{\sum m_i z_i}{\sum m_i}.
    \]
    \item \textbf{Check} symmetry: if the whole body is symmetric about an axis, the centre of mass lies on that axis.
\end{enumerate}
\end{methode}

\courssub{Property: Toppling Condition on a Tilted Plane}
\begin{propriete}[Toppling Angle]
For a rigid body with a flat base resting on a rough inclined plane (tilt angle $\theta$), toppling occurs when the vertical line through the centre of mass $G$ passes through the edge of the base contact area. If the base is symmetric and the centre of mass lies on the central vertical axis at height $z_G$ above the base, and the half‑width of the base in the direction of tilt is $b$, then the critical angle $\theta$ satisfies
\[
\tan\theta = \frac{b}{z_G}.
\]
\emph{Derivation:} In the tilted position, the base plane makes angle $\theta$ with the horizontal. The horizontal offset of the vertical line through $G$ from the centre of the base is $z_G \tan\theta$. Toppling impends when this offset equals $b$.
\end{propriete}

\courssub{Tasks}
Work in groups of 3–4. Present your calculations, a labelled diagram, and your conclusions on a poster.

\begin{enumerate}
    \item \textbf{Centre of mass of the stand alone.}
    \begin{enumerate}
        \item Decompose the stand into simple uniform laminae: the base frame (with hole), the four posts, and the four triangular braces. Use symmetry to locate the centre of mass of each component in the $(x,y,z)$ coordinate system.
        \item Compute the mass of each component.
        \item Using the addition/subtraction method (treating the hole as negative mass), find the coordinates $(x_S, y_S, z_S)$ of the centre of mass of the complete stand.
    \end{enumerate}

    \item \textbf{Centre of mass of the full structure (stand + full tank).}
    \begin{enumerate}
        \item Determine the centre of mass of the full cylindrical tank. (Hint: for a uniform cylinder, the centre of mass is at its geometric centre.)
        \item Using the coordinate formula for a system of particles, find the coordinates $(x_G, y_G, z_G)$ of the centre of mass of the combined structure (stand + full tank).
    \end{enumerate}

    \item \textbf{Toppling angle.}
    \begin{enumerate}
        \item The stand rests on a horizontal ground. If the ground tilts slowly about a line parallel to the $y$-axis (i.e., the stand tips in the $x$-direction), at what angle $\theta$ does the structure topple? Assume no sliding occurs.
        \item The toppling condition is that the vertical line through the combined centre of mass passes through the edge of the base. Derive the formula for $\theta$ in terms of $x_G$, $z_G$, and the half‑width of the base $b = \SI{0.6}{m}$.
        \item Calculate $\theta$ (in degrees) for the given dimensions.
    \end{enumerate}

    \item \textbf{Design modification for improved stability.}
    \begin{enumerate}
        \item Propose \textbf{one} modification: either (a) widen the base frame to $\SI{1.6}{m} \times \SI{1.6}{m}$ (keeping the same frame width $\SI{0.15}{m}$ and hole radius), or (b) lower the tank by reducing the post height to $\SI{1.0}{m}$ (keeping the base unchanged).
        \item Recalculate the combined centre of mass and the new toppling angle for your chosen modification.
        \item Compare the two modifications (if time permits, analyse both) and recommend which is more effective and practical for the school.
    \end{enumerate}
\end{enumerate}

\courssub{Model Answers}

\begin{exoresolu}[Task 1: Centre of Mass of the Stand Alone]
\noindent\textbf{Step 1: Decomposition and symmetry.}
The stand is symmetric about the $z$-axis (the $x$ and $y$ axes are axes of symmetry for the whole structure). Therefore, the centre of mass of the stand lies on the $z$-axis:
\[
x_S = 0,\qquad y_S = 0.
\]
We only need to compute $z_S$.

All wooden parts are cut from the same uniform lamina of thickness $t = \SI{0.05}{m}$ and density $\rho = \SI{600}{kg/m^3}$. The surface density (mass per unit area) is
\[
\sigma = \rho t = 600 \times 0.05 = \SI{30}{kg/m^2}.
\]

Components (all uniform laminae):
\begin{enumerate}
    \item \textbf{Base frame (with hole):} The base frame is a square annulus. Outer square side $L_{\text{out}} = 1.2$ m, inner square side $L_{\text{in}} = 1.2 - 2\times0.15 = 0.9$ m. A circular hole of radius $r = 0.05$ m is cut at the centre. We treat the base as:
    \begin{itemize}
        \item Outer square (positive mass).
        \item Inner square (negative mass, representing the removed central part of the frame).
        \item Circular hole (negative mass, but note the inner square already removes the central area; the hole is an \emph{additional} cut inside the inner square? Actually the base frame is the region between the outer and inner squares. The hole is cut through the frame, so it removes material from the frame itself. The frame area is $L_{\text{out}}^2 - L_{\text{in}}^2$. The hole is centred, so it lies entirely within the frame? The inner square has side 0.9 m, so its half-side is 0.45 m. The hole radius is 0.05 m, so the hole fits inside the inner square. But the frame is the region between outer and inner squares; the hole is cut through the frame, meaning it removes material from the frame region. However, the frame region does not include the inner square area. The hole is at the centre, so it would be in the inner square region if the inner square were present. Since the inner square is already removed (it's the empty space inside the frame), the hole is actually cut through the frame material? Wait: The base frame is a square annulus. The hole is at the centre of the base frame. The centre of the base frame is the centre of the inner square. The hole is entirely within the inner square region, which is already empty. So the hole does not remove any additional material from the frame! The problem says: "A circular hole of radius 0.05 m cut through the centre of the base frame to allow a pipe to pass through." If the base frame is a square annulus, the centre is empty. Cutting a hole there would not remove any wood because there is no wood at the centre. This is a contradiction. Perhaps the base frame is a solid square of side 1.2 m with a square hole of side 0.9 m? That is the annulus. The centre of the annulus is the centre of the square hole. A circular hole at the centre would be inside the square hole, so it doesn't cut the frame. To make sense, the base frame might be a solid square of side 1.2 m, and the "frame width 0.15 m" means the frame is the border of width 0.15 m, but the inner part is not completely removed? Actually "rectangular base frame: outer dimensions 1.2 m × 1.2 m, with a uniform width of 0.15 m (like a square picture frame)." That implies the frame is just the border, i.e., a square annulus. The hole is at the centre of the base frame. The centre of the base frame is the centre of the annulus, which is empty. So the hole is in the empty region. That seems odd. Perhaps the base frame is a solid square plate of side 1.2 m, and the "frame width" refers to the width of the border that is reinforced? But it says "like a square picture frame", which suggests an annulus. The hole is "cut through the centre of the base frame". If the base frame is an annulus, its centre is empty. Maybe the base frame is a solid square, and the hole is cut through it. Then the "frame width" might be the width of the border that remains after cutting the hole? That doesn't match "like a square picture frame". Let's re-read: "A rectangular base frame: outer dimensions 1.2 m × 1.2 m, with a uniform width of 0.15 m (like a square picture frame)." This describes a square annulus of outer side 1.2 m and inner side 0.9 m. "A circular hole of radius 0.05 m cut through the centre of the base frame". The centre of the base frame is the centre of the annulus. The hole would be in the empty middle. That is physically meaningless. Perhaps the base frame is a solid square plate, and the "frame width" is the width of the border that is left after cutting a square hole? But then it's not a picture frame, it's a plate with a square hole. The phrase "like a square picture frame" strongly suggests an annulus. The hole might be cut through the frame material itself, i.e., the hole is centred on the frame's centreline? But the frame's centreline is at a distance of 0.525 m from the centre? No. I think the intended interpretation is that the base frame is a solid square plate of side 1.2 m, and a square hole of side 0.9 m is cut out to form a frame of width 0.15 m. Then an additional circular hole of radius 0.05 m is cut at the centre of the plate (which is also the centre of the square hole). But the square hole already removes the centre. So the circular hole is redundant. Unless the square hole is not cut? The description: "A rectangular base frame: outer dimensions 1.2 m × 1.2 m, with a uniform width of 0.15 m (like a square picture frame)." That means the frame itself is the border. The hole is cut through the centre of the base frame. If the base frame is just the border, its centre is empty. So the hole must be cut through the frame material, i.e., the hole is centred on the frame's centreline? That would be a hole in the middle of the border, not at the geometric centre of the square. The phrase "centre of the base frame" likely means the geometric centre of the square (the centre of the outer square). In a picture frame, the centre of the frame is the centre of the picture, which is empty. So the hole is in the empty space. This is confusing. In many such problems, the base frame is a solid square plate with a square hole (making a frame), and then a circular hole is cut at the centre of the plate (which is in the middle of the square hole). That would mean the circular hole removes no material. That can't be. Perhaps the base frame is a solid square plate (no square hole), and the "uniform width of 0.15 m" refers to the width of the frame that is actually the border of the plate? That doesn't make sense. Let's look at the model answer provided by the colleague: they treated the base frame as outer square minus inner square (annulus) and then subtracted a circular hole. They computed outer area 1.44, inner area 0.81, frame area 0.63. Then they subtracted the circular hole area (0.007854) from the frame area. That implies the circular hole is cut through the frame material, i.e., the hole is centred at the centre of the square, but the frame material exists at that centre? But the frame material is only in the annulus between radius 0.45 and 0.6 from the centre. The centre (radius <0.45) is empty. A circular hole of radius 0.05 at the centre would be entirely in the empty region. So subtracting it is wrong. However, the colleague's model answer did exactly that. To stay consistent with the colleague's interpretation (which might be the intended one, even if physically odd), we will follow the same approach: treat the base frame as a square annulus (outer square minus inner square) and then subtract a circular disc at the centre. This effectively means the base frame is a solid square with a square hole and a circular hole at the centre, but the circular hole overlaps the square hole. The net material is the annulus minus the circular disc. Since the circular disc lies entirely inside the square hole, subtracting it from the annulus is equivalent to subtracting it from the solid square? Actually, if we start with a solid square of side 1.2, cut out a square hole of side 0.9, we get the annulus. Then cutting a circular hole at the centre removes material that was already removed (the square hole). So the net material is just the annulus. The colleague's calculation: frame area = outer - inner = 0.63. Then they subtract the circular hole area (0.007854) to get net base area = 0.622146. That means they are treating the base frame as a solid square of side 1.2 with a square hole of side 0.9 and a circular hole of radius 0.05 at the centre, but the circular hole is cut through the remaining material? But the remaining material is the annulus, which does not include the centre. The only way this makes sense is if the "base frame" is a solid square plate of side 1.2 m, and the "uniform width of 0.15 m" is not a hole but a description of the frame's border? That is, the frame is the entire plate, but it has a border of width 0.15 m that is reinforced? No. Given the colleague's model answer is part of the input, we must correct it if it's wrong, but we are to review and improve. The colleague's model answer might have an error. We need to decide on the correct interpretation. Let's think: The stand is a lamina of thickness 0.05 m. The base frame is a rectangular base frame: outer dimensions 1.2×1.2, uniform width 0.15 m (like a square picture frame). That is a square annulus. The circular hole is cut through the centre of the base frame. If the base frame is an annulus, its centre is empty. So the hole must be cut through the frame material, meaning the hole is not at the geometric centre of the square, but at the centre of the frame's width? That is, the hole is centred on the midline of the frame border. But the problem says "centre of the base frame", which usually means the centre of the square. I suspect the problem intends the base frame to be a solid square plate, and the "frame width" is the width of the border that is left after cutting a square hole? But then it's not a picture frame, it's a plate with a square hole. The phrase "like a square picture frame" means it looks like a picture frame, i.e., a border. So it is an annulus. The hole is cut through the centre of the base frame. In a picture frame, the centre is the empty middle. So the hole is in the empty middle. That is a contradiction. Perhaps the base frame is a solid square, and the "uniform width of 0.15 m" means the frame is the border of width 0.15 m, but the inner part is not cut out? That would be a solid square with a border of width 0.15 m that is somehow distinct? No. I'll assume the intended geometry is: the base frame is a square annulus (outer 1.2, inner 0.9). The circular hole is cut through the frame material, i.e., it is centred on the frame's centreline. But the frame's centreline is at a distance of 0.525 m from the centre? That would be four holes, one on each side? The problem says "a circular hole", singular. So it's one hole at the centre of the square. The only way to have a hole at the centre of the square in an annulus is if the annulus has no hole? That is, the base frame is a solid square plate, and the "frame width" is the width of the frame that is left after cutting a square hole? But then the frame width is 0.15 m, so the inner square side is 1.2 - 2*0.15 = 0.9 m. That is exactly the annulus. Then the circular hole is cut at the centre of the plate, which is the centre of the inner square. Since the inner square is cut out, the circular hole falls in the empty region. So it removes no material. The colleague's model answer subtracts it anyway. That is an error. However, the mass of the hole is very small (0.2356 kg) compared to the total stand mass (~80 kg). The effect on the final answer is negligible. For the sake of consistency with the colleague's model answer (which might be the official expected answer), we will keep the subtraction of the circular hole, but we should note the ambiguity. In a rigorous solution, we should state our interpretation. I will interpret the base frame as a solid square plate of side 1.2 m with a square hole of side 0.9 m (making a frame of width 0.15 m) and an additional circular hole of radius 0.05 m at the centre. But since the circular hole lies entirely within the square hole, it removes no material. Therefore, the net base frame is just the annulus. The colleague subtracted the circular hole from the annulus, which is incorrect. I will correct this: the base frame mass is just the annulus mass. The hole is in the empty centre, so it does not affect the mass. However, the problem might have intended the base frame to be a solid square plate (no square hole) with a circular hole at the centre, and the "uniform width of 0.15 m" refers to something else? That seems unlikely. Let's check the phrase: "A rectangular base frame: outer dimensions 1.2 m × 1.2 m, with a uniform width of 0.15 m (like a square picture frame)." This is a classic description of a square annulus. The circular hole is an additional feature. In many engineering contexts, a "frame" might have a hole in the middle of the frame material for a pipe. But the centre of the frame is the centre of the square. If the frame is an annulus, the centre is empty. So the hole must be in the frame material, meaning it is not at the geometric centre of the square, but at the centre of the frame's width? That would be four holes (one on each side). The problem says "a circular hole", singular. So it's one hole at the centre of the square. I think the problem has a slight flaw. The colleague's model answer treats the base frame as an annulus and then subtracts a central circular hole. That is mathematically inconsistent but yields a number. I will follow the colleague's approach for consistency, but add a note about the ambiguity. Actually, as an expert reviewer, I should correct the error. The correct interpretation: the base frame is a square annulus. The circular hole is at the centre of the square, which is in the empty region. Therefore, the hole does not remove any material from the frame. The mass of the base frame is simply the mass of the annulus. I will use this correct interpretation. The colleague's model answer will be corrected accordingly. This will slightly change the numbers. Let's recompute with the correct interpretation.

    \textbf{Correct interpretation:} The base frame is a square annulus (outer square 1.2 m, inner square 0.9 m). The circular hole at the centre lies in the empty inner square, so it does not remove any wood. Therefore, the base frame mass is just the annulus mass. The hole is irrelevant for mass calculation (it's just a hole in the empty space). However, if the hole were cut through the frame material, it would have to be centred on the frame's midline, but then there would be four holes. Since the problem says "a circular hole", we assume it's at the geometric centre and does not affect the frame mass. We'll proceed with the annulus only.
    \end{itemize}
    \item \textbf{Four vertical posts:} Each post is a rectangular lamina of width $w = 0.15$ m and height $h = 1.5$ m. They are positioned at the four corners of the outer square. Their centres in the $xy$-plane are at $(\pm0.525, \pm0.525)$ (since the outer square extends from $-0.6$ to $0.6$, the posts occupy the corners: $x \in [0.45,0.6]$ and $y \in [0.45,0.6]$ for the first quadrant post, so its centre is at $(0.525, 0.525)$). By symmetry, the combined centre of mass of the four posts is at $(0,0,z_{\text{post}})$. The bottom of the posts sits on the top of the base frame at $z = 0.05$ m. The centre of mass of each post is at half its height: $z_{\text{post}} = 0.05 + 1.5/2 = 0.8$ m.
    \item \textbf{Four triangular braces:} Each brace is a right‑angled triangular lamina with vertical leg $1.5$ m and horizontal leg $0.4$ m. They are attached to the outer faces of the posts. For the brace on the $+x$ face, it lies in the plane $x = 0.6$ (the outer face of the posts at $x=0.6$), spanning vertically from $z=0.05$ to $z=1.55$ and horizontally (along $y$) from $y=0$ to $y=0.4$? The problem says "on each side", so there is one brace per side. With four sides, there are four braces. Each brace is attached to the outer face of the posts on that side. Since the posts are at the corners, a brace on the $+x$ side would connect the two posts on that side? But the legs are 1.5 m (vertical) and 0.4 m (horizontal). The horizontal leg of 0.4 m is much smaller than the distance between posts (1.05 m). So each brace is a small triangle at the corner of the post? "Attached to the outer face of the posts" suggests each brace is attached to one post. There are four posts and four braces, so one brace per post. Each brace is a right triangle with the right angle at the bottom outer corner of the post. The vertical leg goes up the post (1.5 m), the horizontal leg goes outward from the post (0.4 m). So the brace lies in a vertical plane perpendicular to the face of the post. For the post at $(0.525,0.525)$, the outer faces are at $x=0.6$ and $y=0.6$. A brace on the $+x$ face would be in the plane $y=0.525$? Actually, the post is a lamina of width 0.15 m in $x$ and 0.15 m in $y$? Wait, the post is a rectangle 0.15 m × 1.5 m. Which dimension is width? The cross-section is 0.15 m × 0.15 m? The problem says "each a rectangle 0.15 m × 0.15 m in cross‑section and height 1.5 m." That suggests the post is a square tube? But it's a lamina of thickness 0.05 m. A lamina cannot have a cross-section of 0.15×0.15. This is a contradiction. The colleague treated the post as a lamina of area 0.15×1.5. That means the post is a flat plate of width 0.15 m (in the horizontal plane) and height 1.5 m. Its thickness is 0.05 m. So its cross-section (horizontal) is 0.15 m × 0.05 m. The problem's "cross‑section 0.15 m × 0.15 m" is likely a mistake; it should be "width 0.15 m, height 1.5 m". We'll go with the lamina interpretation: each post is a vertical rectangle of width 0.15 m (horizontal) and height 1.5 m. The four posts are at the corners, so their widths are along $x$ and $y$. For the post at the $+x,+y$ corner, it occupies $x \in [0.45,0.6]$, $y \in [0.45,0.6]$. Its outer faces are at $x=0.6$ and $y=0.6$. A brace on the $+x$ side attached to this post would be on the face $x=0.6$, spanning $y \in [0.45,0.6]$? But the brace has horizontal leg 0.4 m. If it's attached to the outer face, the horizontal leg could be along $y$ (parallel to the base) or along $x$ (protruding outward). Typically a brace protrudes outward, so the horizontal leg is along $x$ (from $x=0.6$ to $x=1.0$). Then the brace lies in a plane of constant $y$ (the $y$-coordinate of the post's centre? or the post's face?). The post has width 0.15 m in $y$, so its faces are at $y=0.45$ and $y=0.6$. The brace could be attached to the centre of the face? The problem is ambiguous. However, due to symmetry, the exact $x,y$ positions of the braces do not affect the overall $x_S,y_S$ (they will cancel). We only need the $z$-coordinate of each brace's centre of mass. For a right triangle with vertical leg 1.5 m, the centroid is at $1/3$ of the height from the base. The base of the brace is at the bottom of the post, $z=0.05$ m. So $z_{\text{brace}} = 0.05 + 1.5/3 = 0.55$ m. By symmetry, the four braces together have centre of mass at $(0,0,0.55)$.
\end{enumerate}

\textbf{Step 2: Masses of components.}
Surface density $\sigma = \SI{30}{kg/m^2}$.
\begin{itemize}
    \item \textbf{Base frame (annulus only):} Area $A_{\text{base}} = 1.2^2 - 0.9^2 = 1.44 - 0.81 = \SI{0.63}{m^2}$. Mass $M_{\text{base}} = 0.63 \times 30 = \SI{18.9}{kg}$. Centre of mass at $z_{\text{base}} = t/2 = \SI{0.025}{m}$.
    \item \textbf{Four posts:} Area per post $A_{\text{post}} = 0.15 \times 1.5 = \SI{0.225}{m^2}$. Mass per post $m_{\text{post}} = 0.225 \times 30 = \SI{6.75}{kg}$. Total $M_{\text{posts}} = 4 \times 6.75 = \SI{27}{kg}$. $z_{\text{posts}} = \SI{0.8}{m}$.
    \item \textbf{Four braces:} Area per brace $A_{\text{brace}} = \frac{1}{2} \times 1.5 \times 0.4 = \SI{0.3}{m^2}$. Mass per brace $m_{\text{brace}} = 0.3 \times 30 = \SI{9}{kg}$. Total $M_{\text{braces}} = 4 \times 9 = \SI{36}{kg}$. $z_{\text{braces}} = \SI{0.55}{m}$.
\end{itemize}
Total mass of stand:
\[
M_S = M_{\text{base}} + M_{\text{posts}} + M_{\text{braces}} = 18.9 + 27 + 36 = \SI{81.9}{kg}.
\]

\textbf{Step 3: $z$-coordinate of stand's centre of mass.}
\[
z_S = \frac{M_{\text{base}} z_{\text{base}} + M_{\text{posts}} z_{\text{posts}} + M_{\text{braces}} z_{\text{braces}}}{M_S}
\]
\[
= \frac{18.9 \times 0.025 + 27 \times 0.8 + 36 \times 0.55}{81.9}
= \frac{0.4725 + 21.6 + 19.8}{81.9} = \frac{41.8725}{81.9} \approx \SI{0.5113}{m}.
\]
Thus the stand's centre of mass is at $(0, 0, \SI{0.511}{m})$.
\end{exoresolu}

\begin{exoresolu}[Task 2: Centre of Mass of the Full Structure]
\noindent\textbf{Tank:} The tank is a uniform right circular cylinder of radius $R = \SI{0.6}{m}$ and height $H = \SI{1.8}{m}$. Its centre of mass is at its geometric centre. The tank sits on the tops of the posts. The top of the posts is at $z = 0.05 + 1.5 = \SI{1.55}{m}$. The tank height is $1.8$ m, so its centre is at
\[
z_T = 1.55 + \frac{1.8}{2} = \SI{2.45}{m}.
\]
The tank mass is $M_T = \SI{5000}{kg}$ (given). By symmetry, $x_T = y_T = 0$.

\textbf{Combined structure (stand + tank):} Total mass $M_{\text{tot}} = M_S + M_T = 81.9 + 5000 = \SI{5081.9}{kg}$.
The overall centre of mass coordinates are
\[
x_G = y_G = 0,
\]
\[
z_G = \frac{M_S z_S + M_T z_T}{M_{\text{tot}}}
= \frac{81.9 \times 0.5113 + 5000 \times 2.45}{5081.9}
= \frac{41.88 + 12250}{5081.9} = \frac{12291.88}{5081.9} \approx \SI{2.419}{m}.
\]
So the combined centre of mass is at $(0, 0, \SI{2.42}{m})$.
\end{exoresolu}

\begin{exoresolu}[Task 3: Toppling Angle]
The stand rests on the ground. The base frame outer dimension is $\SI{1.2}{m}$, so the contact area with the ground is the outer square of the base frame (the outer edges are at $x = \pm0.6$, $y = \pm0.6$). The half‑width in the $x$-direction is $b = \SI{0.6}{m}$.

When the ground tilts about a line parallel to the $y$-axis by an angle $\theta$, the base plane becomes inclined. The structure will topple when the vertical line through the combined centre of mass $G$ passes through the edge of the base at $x = +b$ (for tilt in the $+x$ direction).

In the body frame (attached to the stand), the centre of mass is at $(0, 0, z_G)$. When the stand tilts by $\theta$, the vertical direction makes an angle $\theta$ with the body's $z$-axis. The horizontal offset of the vertical line through $G$ from the centre of the base is $z_G \tan\theta$. Toppling impends when this offset equals the half‑width $b$:
\[
z_G \tan\theta = b \quad\Longrightarrow\quad \tan\theta = \frac{b}{z_G}.
\]

Substituting $b = \SI{0.6}{m}$ and $z_G = \SI{2.419}{m}$:
\[
\tan\theta = \frac{0.6}{2.419} = 0.2480 \quad\Longrightarrow\quad \theta = \arctan(0.2480) \approx \SI{13.9}{\ensuremath{{}^{\circ}}}.
\]
The structure topples at a tilt of about $\SI{13.9}{\ensuremath{{}^{\circ}}}$.
\end{exoresolu}

\begin{exoresolu}[Task 4: Design Modification for Improved Stability]
We analyse both proposed modifications.

\textbf{Modification A: Widen base to $\SI{1.6}{m} \times \SI{1.6}{m}$.}
Outer side $L = 1.6$ m, frame width $w = 0.15$ m, so inner side $= 1.6 - 2\times0.15 = 1.3$ m.
Base frame area (annulus) $= 1.6^2 - 1.3^2 = 2.56 - 1.69 = \SI{0.87}{m^2}$.
Mass $M_{\text{base}} = 0.87 \times 30 = \SI{26.1}{kg}$. $z_{\text{base}} = \SI{0.025}{m}$.
Posts: unchanged dimensions and height, so $M_{\text{posts}} = \SI{27}{kg}$, $z_{\text{posts}} = \SI{0.8}{m}$.
Braces: unchanged, $M_{\text{braces}} = \SI{36}{kg}$, $z_{\text{braces}} = \SI{0.55}{m}$.
Stand mass $M_S' = 26.1 + 27 + 36 = \SI{89.1}{kg}$.
$z_S' = \frac{26.1\times0.025 + 27\times0.8 + 36\times0.55}{89.1} = \frac{0.6525 + 21.6 + 19.8}{89.1} = \frac{42.0525}{89.1} \approx \SI{0.472}{m}$.
Tank unchanged: $M_T = \SI{5000}{kg}$, $z_T = \SI{2.45}{m}$.
Combined $z_G' = \frac{89.1\times0.472 + 5000\times2.45}{5089.1} = \frac{42.06 + 12250}{5089.1} = \frac{12292.06}{5089.1} \approx \SI{2.415}{m}$.
New half-width $b' = 1.6/2 = \SI{0.8}{m}$.
$\tan\theta_A = \frac{0.8}{2.415} = 0.3313 \Rightarrow \theta_A \approx \SI{18.3}{\ensuremath{{}^{\circ}}}$.

\textbf{Modification B: Lower tank by reducing post height to $\SI{1.0}{m}$.}
Base unchanged: $M_{\text{base}} = \SI{18.9}{kg}$, $z_{\text{base}} = \SI{0.025}{m}$.
Posts: height $1.0$ m, area per post $= 0.15 \times 1.0 = \SI{0.15}{m^2}$, mass per post $= 4.5$ kg, total $M_{\text{posts}} = \SI{18}{kg}$. $z_{\text{posts}} = 0.05 + 1.0/2 = \SI{0.55}{m}$.
Braces: vertical leg now $1.0$ m, area per brace $= \frac{1}{2} \times 1.0 \times 0.4 = \SI{0.2}{m^2}$, mass per brace $= 6$ kg, total $M_{\text{braces}} = \SI{24}{kg}$. $z_{\text{braces}} = 0.05 + 1.0/3 = \SI{0.3833}{m}$.
Stand mass $M_S'' = 18.9 + 18 + 24 = \SI{60.9}{kg}$.
$z_S'' = \frac{18.9\times0.025 + 18\times0.55 + 24\times0.3833}{60.9} = \frac{0.4725 + 9.9 + 9.2}{60.9} = \frac{19.5725}{60.9} \approx \SI{0.321}{m}$.
Tank bottom now at $z = 0.05 + 1.0 = \SI{1.05}{m}$. Tank centre $z_T'' = 1.05 + 0.9 = \SI{1.95}{m}$.
Combined $z_G'' = \frac{60.9\times0.321 + 5000\times1.95}{5060.9} = \frac{19.55 + 9750}{5060.9} = \frac{9769.55}{5060.9} \approx \SI{1.930}{m}$.
Half-width $b = \SI{0.6}{m}$ (base unchanged).
$\tan\theta_B = \frac{0.6}{1.930} = 0.3109 \Rightarrow \theta_B \approx \SI{17.3}{\ensuremath{{}^{\circ}}}$.

\textbf{Comparison:}
\begin{center}
\begin{tabular}{|l|c|c|c|}
\hline
\textbf{Design} & \textbf{$z_G$ (m)} & \textbf{Half-width $b$ (m)} & \textbf{Toppling angle $\theta$} \\
\hline
Original & 2.419 & 0.6 & 13.9° \\
Modification A (wider base) & 2.415 & 0.8 & 18.3° \\
Modification B (lower tank) & 1.930 & 0.6 & 17.3° \\
\hline
\end{tabular}
\end{center}
Both modifications significantly increase the toppling angle. Widening the base (Modification A) gives a slightly higher angle (18.3° vs 17.3°). However, widening the base requires more materials and space, while lowering the tank reduces the overall height, which may be easier to construct and also reduces wind loading. The school should consider cost, available space, and ease of construction. A combined approach (moderately wider base and moderately lower tank) could yield even better stability.
\end{exoresolu}

\courssub{Common Mistakes}
\begin{attention}[Pitfalls to Avoid]
\begin{itemize}
    \item \textbf{Forgetting the lamina thickness:} The stand is a lamina of thickness $t=0.05$ m. Masses are computed using surface density $\sigma = \rho t$, not volume density $\rho$ times volume. Using $\rho$ directly would give masses 20 times too large.
    \item \textbf{Misplacing the base of the posts:} The posts sit on top of the base frame, so their bottom is at $z = t = 0.05$ m, not at $z=0$.
    \item \textbf{Incorrect centroid of a triangle:} The centroid of a right triangle is at $1/3$ of the legs from the right angle, not at $1/2$.
    \item \textbf{Ignoring the hole in the base frame:} If the hole is in the frame material, it must be subtracted as negative mass. If it lies in the empty central region (as in our corrected interpretation), it has no effect. Always check geometry carefully.
    \item \textbf{Toppling formula:} The formula $\tan\theta = b/z_G$ assumes the centre of mass is exactly above the centre of the base ($x_G=0$). If $x_G \neq 0$, the formula becomes $\tan\theta = (b - x_G)/z_G$ (for tilt in the $+x$ direction).
    \item \textbf{Units:} Always carry units through calculations. The final angle must be in degrees (or radians) as required.
\end{itemize}
\end{attention}

\courssub{Key Learning Points}
\begin{aretenir}[Summary]
\begin{itemize}
    \item The centre of mass of a composite lamina can be found by treating each part (including holes as negative mass) and using the weighted average formula $\vec{r}_G = \frac{\sum m_i \vec{r}_i}{\sum m_i}$.
    \item Symmetry greatly simplifies calculations: if a body is symmetric about an axis, its centre of mass lies on that axis.
    \item For a system of separate bodies (stand + tank), the overall centre of mass is the mass‑weighted average of the individual centres of mass.
    \item Toppling occurs when the vertical line through the combined centre of mass falls outside the base support area. The critical tilt angle $\theta$ satisfies $\tan\theta = \frac{\text{half‑base width}}{z_G}$ for a symmetric body.
    \item Lowering the centre of mass or widening the base both increase the toppling angle, but the quantitative effect must be calculated for the specific design.
    \item In engineering design, stability improvements must be balanced against cost, space, and practicality.
\end{itemize}
\end{aretenir}

\newpage

\coursec{Methods and worked examples}

\coursec{Centre of Mass of Uniform Rigid Bodies}

\begin{definition}[Centre of Mass]
The \cle{centre of mass} of a system of particles or a rigid body is the unique point $G$ where the total mass of the system can be considered to be concentrated for the purpose of analysing translational motion. If a single external force $\vec{F}$ acts on the system, the centre of mass moves as if it were a particle of mass $M$ (the total mass) subjected to that force: $\vec{F}=M\vec{a}_G$.
\end{definition}

\begin{propriete}[Motion of the Centre of Mass]
For any system of particles (or a rigid body) with total mass $M$, the position vector $\vec{r}_G$ of the centre of mass satisfies
\[
M\vec{r}_G = \sum_{i=1}^n m_i\vec{r}_i,
\]
where $m_i$ and $\vec{r}_i$ are the mass and position vector of the $i$-th particle. Differentiating with respect to time gives the velocity and acceleration of the centre of mass:
\[
M\vec{v}_G = \sum m_i\vec{v}_i,\qquad M\vec{a}_G = \sum m_i\vec{a}_i = \sum \vec{F}_i^{\text{ext}}.
\]
The internal forces cancel in pairs (Newton's third law), so the motion of the centre of mass depends \emph{only} on the external forces.
\end{propriete}

\begin{aretenir}[Key Idea]
The centre of mass is a \emph{mathematical point} — it may lie inside the body (e.g. a solid sphere) or outside it (e.g. a ring). For a \cle{uniform rigid body} (constant density), the centre of mass coincides with the \cle{centroid} of its volume (3D), area (lamina, 2D) or length (wire, 1D).
\end{aretenir}

\courssub{Centre of Mass of a System of Particles}

\begin{methode}[Centre of mass by coordinates (discrete system)]
\begin{enumerate}
\item Choose a convenient origin $O$ and coordinate axes.
\item List each particle: mass $m_i$, coordinates $(x_i,y_i,z_i)$.
\item Compute the total mass $M=\sum m_i$.
\item The coordinates of the centre of mass $G(\bar{x},\bar{y},\bar{z})$ are the weighted means:
\[
\bar{x}=\frac{\sum m_i x_i}{M},\quad
\bar{y}=\frac{\sum m_i y_i}{M},\quad
\bar{z}=\frac{\sum m_i z_i}{M}.
\]
\item If a particle is \emph{removed}, treat it as a particle of \cle{negative mass} (same coordinates, mass $-m_i$).
\item \textbf{Check:} $\bar{x}$ must lie between $\min x_i$ and $\max x_i$; similarly for $\bar{y},\bar{z}$.
\end{enumerate}
\end{methode}

\begin{exoresolu}[Particles on a light rod]
Particles of masses $2$~kg, $3$~kg and $5$~kg are attached to a light rod $AB$ of length $1.2$~m at distances $0$, $0.5$~m and $1.2$~m from $A$ respectively. Find the distance of the centre of mass of the system from $A$.
\tcblower
Take $A$ as origin, the $x$-axis along the rod. The rod is light (massless), so only the particles contribute.
Total mass: $M=2+3+5=10$~kg.
Moment about $A$: $\sum m_i x_i = 2(0)+3(0.5)+5(1.2)=0+1.5+6=7.5$~kg\,m.
\[
\bar{x}=\frac{7.5}{10}=0.75\ \text{m}.
\]
The centre of mass lies $0.75$~m from $A$. Check: $0\le 0.75\le 1.2$, and it is pulled towards the heavier $5$~kg mass, as expected.
\end{exoresolu}

\begin{exoresolu}[Particles in a plane]
Four particles of masses $1$, $2$, $3$ and $4$~kg are placed at the vertices of a square of side $2$~m. Find the coordinates of the centre of mass if the $1$~kg mass is at the origin and the square lies in the first quadrant with sides along the axes.
\tcblower
Coordinates: $m_1=1$ at $(0,0)$; $m_2=2$ at $(2,0)$; $m_3=3$ at $(2,2)$; $m_4=4$ at $(0,2)$.
$M=1+2+3+4=10$~kg.
\[
\bar{x}=\frac{1(0)+2(2)+3(2)+4(0)}{10}=\frac{10}{10}=1,\qquad
\bar{y}=\frac{1(0)+2(0)+3(2)+4(2)}{10}=\frac{14}{10}=1.4.
\]
Centre of mass: $G(1,\,1.4)$. Check: $0\le\bar{x}\le2$, $0\le\bar{y}\le2$.
\end{exoresolu}

\courssub{Centre of Mass Using Position Vectors}

\begin{propriete}[Vector Formula for Centre of Mass]
Let $O$ be any origin. The position vector of the centre of mass $G$ of a system of particles is
\[
\vec{OG}=\bar{\vec{r}}=\frac{\sum m_i\vec{r}_i}{\sum m_i}.
\]
This single vector equation contains the three coordinate equations. It is independent of the choice of origin: if $O'$ is another origin, $\vec{O'G}=\frac{\sum m_i\vec{O'r}_i}{\sum m_i}$.
\end{propriete}

\begin{methode}[Centre of mass via position vectors]
\begin{enumerate}
\item Write the position vector $\vec{r}_i$ of each particle relative to a chosen origin $O$.
\item Apply the vector formula: $\vec{OG}=\frac{\sum m_i\vec{r}_i}{\sum m_i}$.
\item Work component by component: the $x$-, $y$- (and $z$-) coordinates of $G$ are weighted means of the coordinates of the particles.
\item For a \cle{lamina} or rigid body, replace each part by a particle of equal mass placed at its own centre of mass, then apply the same formula.
\end{enumerate}
\end{methode}

\begin{exoresolu}[Three particles in space]
Particles of masses $1$~kg, $2$~kg and $3$~kg are placed at the points $P(1,0,2)$, $Q(2,3,-1)$ and $R(-1,1,4)$ respectively. Find the position vector of the centre of mass of the system.
\tcblower
Total mass $M=1+2+3=6$~kg. Position vectors:
\[
\vec{OP}=\begin{pmatrix}1\\0\\2\end{pmatrix},\quad
\vec{OQ}=\begin{pmatrix}2\\3\\-1\end{pmatrix},\quad
\vec{OR}=\begin{pmatrix}-1\\1\\4\end{pmatrix}.
\]
\[
\vec{OG}=\frac{1}{6}\left[1\begin{pmatrix}1\\0\\2\end{pmatrix}+2\begin{pmatrix}2\\3\\-1\end{pmatrix}+3\begin{pmatrix}-1\\1\\4\end{pmatrix}\right]
=\frac{1}{6}\begin{pmatrix}1+4-3\\0+6+3\\2-2+12\end{pmatrix}
=\frac{1}{6}\begin{pmatrix}2\\9\\12\end{pmatrix}
=\begin{pmatrix}1/3\\3/2\\2\end{pmatrix}.
\]
Hence $G=\left(\tfrac13,\,\tfrac32,\,2\right)$. Each coordinate is the weighted mean: e.g. $\bar{x}=\frac{1(1)+2(2)+3(-1)}{6}=\frac{2}{6}=\frac13$.
\end{exoresolu}

\begin{experience}[Centre of mass of a system of particles — vector approach]
Consider three particles of masses $m_1,m_2,m_3$ at positions $\vec{r}_1,\vec{r}_2,\vec{r}_3$. Show that the centre of mass of the whole system can be found by first combining $m_1$ and $m_2$ into a single particle at their centre of mass $G_{12}$, then combining that with $m_3$. This illustrates the \emph{associativity} of the centre of mass operation.
\end{experience}

\courssub{Uniform Laminae, Symmetry and Standard Results}

\begin{definition}[Uniform Lamina]
A \cle{uniform lamina} is a thin flat sheet of constant area density (mass per unit area) $\sigma$. Its mass is proportional to its area: $m=\sigma A$. For a uniform lamina, we can replace masses by areas in all centre-of-mass formulas.
\end{definition}

\begin{propriete}[Symmetry Principle]
If a uniform lamina has an axis of symmetry, its centre of mass lies on that axis. If it has two (or more) non-parallel axes of symmetry, the centre of mass is at their intersection.
\end{propriete}

\begin{aretenir}[Standard Centres of Mass for Uniform Laminae]
\begin{center}
\begin{tabular}{|l|c|c|}\hline
\textbf{Shape} & \textbf{Centre of Mass} & \textbf{Notes} \\ \hline
Uniform rod (length $L$) & Midpoint & On the rod \\ \hline
Uniform rectangular lamina & Intersection of diagonals & $(\frac{L}{2},\frac{W}{2})$ \\ \hline
Uniform triangular lamina & Centroid (intersection of medians) & $\bar{x}=\frac{x_1+x_2+x_3}{3},\ \bar{y}=\frac{y_1+y_2+y_3}{3}$; $\frac{2}{3}$ along median from vertex \\ \hline
Uniform circular disc & Centre of circle & \\ \hline
Uniform semicircular lamina (radius $r$) & On axis of symmetry, distance $\frac{4r}{3\pi}$ from diameter & $\approx 0.424\,r$ \\ \hline
Uniform quadrant of a disc (radius $r$) & $\left(\frac{4r}{3\pi},\,\frac{4r}{3\pi}\right)$ from the two straight edges & \\ \hline
Uniform circular arc (radius $r$, angle $2\alpha$) & On axis of symmetry, distance $\frac{r\sin\alpha}{\alpha}$ from centre & Arc only, not lamina \\ \hline
\end{tabular}
\end{center}
\end{aretenir}

\begin{exoresolu}[Centre of mass of a triangular lamina]
A uniform triangular lamina has vertices $A(0,0)$, $B(6,0)$ and $C(2,6)$. Find the coordinates of its centre of mass.
\tcblower
For a uniform triangular lamina the centre of mass coincides with the centroid of the vertices:
\[
\bar{x}=\frac{0+6+2}{3}=\frac{8}{3},\qquad \bar{y}=\frac{0+0+6}{3}=2.
\]
Hence $G=\left(\dfrac{8}{3},\,2\right)$.

\emph{Verification by medians.} The midpoint of $AB$ is $M(3,0)$. The median from $C$ joins $C(2,6)$ to $M(3,0)$. The point two-thirds of the way from $C$ to $M$ is
\[
G=C+\tfrac23(M-C)=\left(2+\tfrac23(1),\ 6+\tfrac23(-6)\right)=\left(\tfrac83,\ 2\right),
\]
which agrees with the formula.
\end{exoresolu}

\begin{exoresolu}[Semicircular lamina]
A uniform lamina is a semicircle of radius $6$~cm. Find the distance of its centre of mass from the straight edge.
\tcblower
For a semicircular lamina, the centre of mass lies on the axis of symmetry at a distance $\frac{4r}{3\pi}$ from the diameter (straight edge).
\[
d=\frac{4\times 6}{3\pi}=\frac{24}{3\pi}=\frac{8}{\pi}\approx 2.55\ \text{cm}.
\]
\end{exoresolu}

\courssub{Composite Laminae (Addition and Subtraction)}

\begin{methode}[Centre of mass of a composite lamina]
\begin{enumerate}
\item Split the lamina into \cle{standard parts} (rectangles, triangles, discs, semicircles, quadrants) whose centres of mass are known.
\item Since the lamina is uniform, masses are proportional to \cle{areas}: use areas as "masses".
\item Tabulate: area $A_i$, coordinates $(x_i,y_i)$ of the centre of each part, and moments $A_i x_i$, $A_i y_i$.
\item A hole or cut-out part counts with a \cle{negative area}.
\item Compute $\bar{x}=\dfrac{\sum A_i x_i}{\sum A_i}$, $\bar{y}=\dfrac{\sum A_i y_i}{\sum A_i}$.
\end{enumerate}
\end{methode}

\begin{exoresolu}[Rectangle with a circular hole]
A uniform rectangular lamina $ABCD$ has $AB=12$~cm and $BC=8$~cm. A circular hole of radius $2$~cm is drilled, its centre being $4$~cm from $AB$ and $4$~cm from $AD$. Find the distance of the centre of mass of the remaining lamina from $AB$ and from $AD$.
\tcblower
Set axes with $A$ as origin, $AD$ along the $x$-axis (length $12$) and $AB$ along the $y$-axis (height $8$).

\textbf{Areas.}
Rectangle: $A_1=12\times 8=96$~cm$^2$, centre $(6,4)$.
Hole: $A_2=-\pi(2)^2=-4\pi$~cm$^2$, centre $(4,4)$.

\textbf{Moment about $AB$ (i.e.\ for $\bar{x}$):}
\[
\sum A_i x_i = 96(6)+(-4\pi)(4)=576-16\pi.
\]
Total area: $\sum A_i=96-4\pi$. Hence
\[
\bar{x}=\frac{576-16\pi}{96-4\pi}=\frac{144-4\pi}{24-\pi}\approx\frac{131.43}{20.86}\approx 6.30\ \text{cm}.
\]

\textbf{Moment about $AD$ (i.e.\ for $\bar{y}$):}
\[
\sum A_i y_i=96(4)+(-4\pi)(4)=384-16\pi,
\]
\[
\bar{y}=\frac{384-16\pi}{96-4\pi}=\frac{96-4\pi}{24-\pi}\approx\frac{83.43}{20.86}\approx 4.00\ \text{cm}.
\]
The centre of mass is about $6.30$~cm from $AB$ and $4.00$~cm from $AD$. Note $\bar{y}$ is unchanged because the hole is symmetrically placed with respect to the horizontal midline $y=4$: the lamina retains $y=4$ as an axis of symmetry, so $\bar{y}=4$ exactly — a useful symmetry check.
\end{exoresolu}

\begin{exoresolu}[L-shaped lamina]
A uniform lamina consists of a $10\times 10$~cm square with a $4\times 4$~cm square removed from one corner. Find the centre of mass of the remaining lamina.
\tcblower
Place the large square with vertices $(0,0)$, $(10,0)$, $(10,10)$, $(0,10)$. Remove the small square at the top-right corner: vertices $(6,6)$, $(10,6)$, $(10,10)$, $(6,10)$.

Areas: $A_1=100$ (large square), centre $(5,5)$; $A_2=-16$ (removed square), centre $(8,8)$.
Total area $=84$~cm$^2$.
\[
\bar{x}=\frac{100(5)+(-16)(8)}{84}=\frac{500-128}{84}=\frac{372}{84}=\frac{31}{7}\approx 4.43\ \text{cm},
\]
\[
\bar{y}=\frac{100(5)+(-16)(8)}{84}=\frac{31}{7}\approx 4.43\ \text{cm}.
\]
By symmetry about $y=x$, $\bar{x}=\bar{y}$.
\end{exoresolu}

\begin{exoresolu}[Composite lamina: rectangle + triangle]
A uniform lamina consists of a rectangle $ABCD$ with $AB=8$~cm, $BC=6$~cm, and an isosceles triangle $CDE$ of base $CD=8$~cm and height $4$~cm attached externally on side $CD$. Find the distance of the centre of mass from $AB$.
\tcblower
Place $A(0,0)$, $B(8,0)$, $C(8,6)$, $D(0,6)$. Triangle base $CD$ from $(0,6)$ to $(8,6)$, apex $E(4,10)$.

Rectangle: area $A_1=48$, centre $(4,3)$.
Triangle: area $A_2=\frac12\times 8\times 4=16$, centroid at $\left(\frac{0+8+4}{3},\frac{6+6+10}{3}\right)=(4,\frac{22}{3})$.

Total area $=64$.
\[
\bar{y}=\frac{48(3)+16(\frac{22}{3})}{64}=\frac{144+\frac{352}{3}}{64}=\frac{\frac{432+352}{3}}{64}=\frac{784}{192}=\frac{49}{12}\approx 4.083\ \text{cm}.
\]
Distance from $AB$ (the $x$-axis) is $\bar{y}=\frac{49}{12}$~cm.
\end{exoresolu}

\courssub{Applications — Equilibrium, Suspension and Toppling}

\begin{methode}[Solving suspension and toppling problems]
\begin{enumerate}
\item \cle{Suspension}: When a lamina hangs freely from a point $P$, the centre of mass $G$ lies \emph{vertically below} $P$. Draw the vertical line through $P$ and use trigonometry to find the angle an edge makes with the vertical (or horizontal).
\item \cle{Toppling}: A body on an inclined plane topples when the vertical through $G$ passes through the lowest edge (corner) of the base. At the limiting position, $\tan\theta=\dfrac{\text{half-width of base}}{\text{height of }G}$ (for a rectangular block).
\item Always locate $G$ first (composite-lamina method), then draw a clear diagram with the vertical through $G$.
\end{enumerate}
\end{methode}

\begin{exoresolu}[Lamina suspended from a vertex]
A uniform lamina consists of a square $ABCD$ of side $10$~cm with an isosceles triangle of base $CD$ and height $6$~cm attached externally on the side $CD$. The lamina is suspended freely from $A$. Find the angle that $AB$ makes with the vertical.
\tcblower
Place $A(0,0)$, $B(10,0)$, $C(10,10)$, $D(0,10)$; the triangle has base $CD$ and apex $E(5,16)$.

\textbf{Parts.}
Square: area $100$, centre $(5,5)$.
Triangle: area $\tfrac12(10)(6)=30$, centroid at $\bar y=\dfrac{10+10+16}{3}=12$, $\bar x=5$, i.e.\ $(5,12)$.

\textbf{Centre of mass.} By symmetry about $x=5$, $\bar{x}=5$.
\[
\bar{y}=\frac{100(5)+30(12)}{100+30}=\frac{500+360}{130}=\frac{860}{130}=\frac{86}{13}\approx 6.615\ \text{cm}.
\]

\textbf{Suspension.} When hanging from $A$, $AG$ is vertical. The vector $\vec{AG}=(5,\,86/13)$. The angle $\alpha$ that $AG$ makes with $AB$ (the $x$-axis) satisfies
\[
\tan\alpha=\frac{\bar y}{\bar x}=\frac{86/13}{5}=\frac{86}{65}\approx 1.3231\quad\Rightarrow\quad \alpha\approx 52.9^{\circ}.
\]
Since $AG$ is vertical, $AB$ makes an angle
\[
90^{\circ}-\alpha\approx 90^{\circ}-52.9^{\circ}=37.1^{\circ}
\]
with the vertical. So $AB$ makes approximately $37.1^{\circ}$ with the downward vertical.
\end{exoresolu}

\begin{exoresolu}[Toppling of a composite lamina on an incline]
A uniform lamina is made by attaching a right-angled triangle of legs $6$~cm and $8$~cm to a rectangle $8$~cm $\times$ $4$~cm along the $8$~cm side. The lamina is placed on a rough inclined plane with the rectangle's $8$~cm side on the plane. Find the maximum angle of inclination before the lamina topples.
\tcblower
Place rectangle on $x$-axis from $x=0$ to $x=8$, height $4$~cm. Triangle attached on top: vertices $(0,4)$, $(8,4)$, $(0,10)$ (right angle at $(0,4)$).

Rectangle: area $32$, centre $(4,2)$.
Triangle: area $\frac12\times 8\times 6=24$, centroid $\left(\frac{0+8+0}{3},\frac{4+4+10}{3}\right)=\left(\frac{8}{3},6\right)$.

Total area $=56$.
\[
\bar{x}=\frac{32(4)+24(\frac{8}{3})}{56}=\frac{128+64}{56}=\frac{192}{56}=\frac{24}{7}\approx 3.43\ \text{cm},
\]
\[
\bar{y}=\frac{32(2)+24(6)}{56}=\frac{64+144}{56}=\frac{208}{56}=\frac{26}{7}\approx 3.71\ \text{cm}.
\]

The base is from $x=0$ to $x=8$. The lamina topples when the vertical through $G$ passes through the lower edge $x=0$ (assuming incline rises to the right). At limiting equilibrium, $\tan\theta = \frac{\bar{x}}{\bar{y}} = \frac{24/7}{26/7} = \frac{24}{26} = \frac{12}{13}$.
\[
\theta = \arctan\left(\frac{12}{13}\right) \approx 42.7^{\circ}.
\]
\end{exoresolu}

\courssub{Centroid of a Region Bounded by Curves (Integration)}

\begin{definition}[Centroid of a Plane Region]
The \cle{centroid} of a plane region is the centre of mass of a uniform lamina of that shape. For a region of area $A$, the coordinates $(\bar{x},\bar{y})$ are given by
\[
\bar{x}=\frac{1}{A}\iint_R x\,dA,\qquad \bar{y}=\frac{1}{A}\iint_R y\,dA.
\]
For a region bounded by $y=f(x)$, the $x$-axis, and $x=a$, $x=b$ ($f(x)\ge 0$), using vertical strips:
\[
A=\int_a^b f(x)\,dx,\quad
\bar{x}=\frac{1}{A}\int_a^b x f(x)\,dx,\quad
\bar{y}=\frac{1}{2A}\int_a^b [f(x)]^2\,dx.
\]
The factor $\frac12$ in $\bar{y}$ arises because the centroid of a thin vertical strip of height $y$ is at its midpoint, $y/2$.
\end{definition}

\begin{methode}[Centroid by integration (vertical strips)]
\begin{enumerate}
\item Sketch the region bounded by $y=f(x)$, the $x$-axis and the lines $x=a$, $x=b$.
\item Compute the area $A=\displaystyle\int_a^b f(x)\,dx$.
\item Use a thin vertical strip of width $dx$: its area is $y\,dx$ and its centroid is at $(x,\,y/2)$.
\item Then
\[
\bar{x}=\frac{1}{A}\int_a^b x\,f(x)\,dx,\qquad \bar{y}=\frac{1}{2A}\int_a^b [f(x)]^2\,dx.
\]
\item For a region between two curves $y=f(x)\ge g(x)$ on $[a,b]$, use
$A=\int_a^b(f-g)\,dx$,
$\bar{x}=\dfrac1A\int_a^b x(f-g)\,dx$,
$\bar{y}=\dfrac{1}{2A}\int_a^b (f^2-g^2)\,dx$.
\end{enumerate}
\end{methode}

\begin{exoresolu}[Centroid under a parabola]
Find the coordinates of the centroid of the region bounded by the curve $y=4-x^2$, the $x$-axis and the $y$-axis in the first quadrant.
\tcblower
The curve meets the $x$-axis at $x=2$ (since $4-x^2=0$ gives $x=2$ in the first quadrant).

\textbf{Area.}
\[
A=\int_0^2 (4-x^2)\,dx=\left[4x-\frac{x^3}{3}\right]_0^2=8-\frac83=\frac{16}{3}.
\]

\textbf{Moment about the $y$-axis.}
\[
\int_0^2 x(4-x^2)\,dx=\int_0^2(4x-x^3)\,dx=\left[2x^2-\frac{x^4}{4}\right]_0^2=8-4=4.
\]
\[
\bar{x}=\frac{4}{16/3}=\frac{12}{16}=\frac34.
\]

\textbf{Moment about the $x$-axis.}
\[
\frac12\int_0^2 (4-x^2)^2\,dx=\frac12\int_0^2(16-8x^2+x^4)\,dx=\frac12\left[16x-\frac{8x^3}{3}+\frac{x^5}{5}\right]_0^2
\]
\[
=\frac12\left(32-\frac{64}{3}+\frac{32}{5}\right)=\frac12\cdot\frac{480-320+96}{15}=\frac12\cdot\frac{256}{15}=\frac{128}{15}.
\]
\[
\bar{y}=\frac{128/15}{16/3}=\frac{128}{15}\times\frac{3}{16}=\frac{8}{5}.
\]
Hence the centroid is $G=\left(\dfrac34,\ \dfrac85\right)$.
\end{exoresolu}

\begin{exoresolu}[Centroid between two curves]
Find the centroid of the region bounded by $y=x^2$ and $y=x$ in the first quadrant.
\tcblower
Curves intersect at $x=0$ and $x=1$. On $[0,1]$, $x\ge x^2$.
Area: $A=\int_0^1 (x-x^2)\,dx=\left[\frac{x^2}{2}-\frac{x^3}{3}\right]_0^1=\frac12-\frac13=\frac16$.
\[
\bar{x}=\frac{1}{1/6}\int_0^1 x(x-x^2)\,dx=6\int_0^1 (x^2-x^3)\,dx=6\left[\frac{x^3}{3}-\frac{x^4}{4}\right]_0^1=6\left(\frac13-\frac14\right)=6\cdot\frac1{12}=\frac12.
\]
\[
\bar{y}=\frac{1}{2(1/6)}\int_0^1 (x^2-x^4)\,dx=3\left[\frac{x^3}{3}-\frac{x^5}{5}\right]_0^1=3\left(\frac13-\frac15\right)=3\cdot\frac{2}{15}=\frac{2}{5}.
\]
Centroid: $G\left(\frac12,\,\frac25\right)$.
\end{exoresolu}

\begin{savaistu}[Pappus–Guldin Theorems]
The centroid of a plane curve or region can be used to find surface areas and volumes of solids of revolution without integration:
\begin{itemize}
\item \textbf{First theorem:} Area of surface of revolution = (arc length of curve) $\times$ (distance travelled by centroid of curve).
\item \textbf{Second theorem:} Volume of solid of revolution = (area of region) $\times$ (distance travelled by centroid of region).
\end{itemize}
These are powerful shortcuts in mechanics and engineering.
\end{savaistu}

\courssub{Centroids of Composite Areas}

\begin{methode}[Centroid of a composite area]
\begin{enumerate}
\item Divide the area into simple shapes (rectangles, triangles, circles, semicircles, etc.) whose centroids are known.
\item Choose a reference axis system.
\item For each part, note its area $A_i$ and the coordinates $(\bar{x}_i,\bar{y}_i)$ of its centroid.
\item If a part is a hole, assign it a \cle{negative area}.
\item Apply the weighted mean formulas:
\[
\bar{x}=\frac{\sum A_i\bar{x}_i}{\sum A_i},\qquad \bar{y}=\frac{\sum A_i\bar{y}_i}{\sum A_i}.
\]
\end{enumerate}
\end{methode}

\begin{exoresolu}[Composite area: rectangle with semicircular cutout]
A rectangular plate $10$~cm $\times$ $6$~cm has a semicircular notch of radius $2$~cm removed from the middle of one $10$~cm side. Find the centroid of the remaining plate.
\tcblower
Place rectangle with corners $(0,0)$, $(10,0)$, $(10,6)$, $(0,6)$. Semicircle removed from bottom side ($y=0$), centre at $(5,0)$, radius $2$, lying below $y=0$ (so it's a cutout from the rectangle).

Rectangle: $A_1=60$, centroid $(5,3)$.
Semicircle (hole): $A_2=-\frac12\pi(2)^2=-2\pi$, centroid at $(5, -\frac{4r}{3\pi})=(5,-\frac{8}{3\pi})$.

Total area $=60-2\pi$.
\[
\bar{x}=\frac{60(5)+(-2\pi)(5)}{60-2\pi}=5 \quad\text{(by symmetry)}.
\]
\[
\bar{y}=\frac{60(3)+(-2\pi)(-\frac{8}{3\pi})}{60-2\pi}=\frac{180+\frac{16}{3}}{60-2\pi}=\frac{\frac{540+16}{3}}{60-2\pi}=\frac{556}{3(60-2\pi)}=\frac{278}{3(30-\pi)}\approx \frac{278}{3(26.858)}\approx 3.45\ \text{cm}.
\]
The centroid is raised slightly above the rectangle's centre ($y=3$) because mass is removed from the bottom.
\end{exoresolu}

\courssub{Centre of Mass of Uniform Solids (Three Dimensions)}

\begin{propriete}[Standard Results for Uniform Solids]
\begin{center}
\begin{tabular}{|l|c|}\hline
\textbf{Solid} & \textbf{Centre of Mass} \\ \hline
Solid cone or pyramid (height $h$) & On axis, $\frac{h}{4}$ above base ($\frac{3h}{4}$ from vertex) \\ \hline
Solid hemisphere (radius $r$) & On axis, $\frac{3r}{8}$ from centre of flat face \\ \hline
Hemispherical shell (radius $r$) & On axis, $\frac{r}{2}$ from centre of flat face \\ \hline
Conical shell (height $h$) & On axis, $\frac{h}{3}$ above base \\ \hline
Solid cylinder / prism & At centre (midpoint of axis) \\ \hline
Solid sphere & At centre \\ \hline
\end{tabular}
\end{center}
\end{propriete}

\begin{methode}[Centre of mass of composite solids]
\begin{enumerate}
\item Use \cle{symmetry} first: $G$ lies on every plane or axis of symmetry.
\item For a \cle{composite solid}, replace each part by a particle of mass $\rho V_i$ at its centre of mass and use $\bar{x}=\dfrac{\sum V_i x_i}{\sum V_i}$ (density $\rho$ cancels if uniform).
\item For a \cle{solid of revolution} about the $x$-axis, $\bar{x}=\dfrac{\displaystyle\int_a^b \pi x y^2\,dx}{\displaystyle\int_a^b \pi y^2\,dx}$, with $\bar y=\bar z=0$ by symmetry.
\end{enumerate}
\end{methode}

\begin{exoresolu}[Cylinder topped by a hemisphere]
A uniform solid consists of a cylinder of radius $3$~cm and height $8$~cm, surmounted by a hemisphere of radius $3$~cm on one circular face. Find the height of the centre of mass of the solid above the base of the cylinder.
\tcblower
By symmetry $G$ lies on the common axis. Measure heights from the base of the cylinder.

\textbf{Cylinder.} Volume $V_1=\pi r^2 h=\pi(9)(8)=72\pi$; centre of mass at height $x_1=4$.

\textbf{Hemisphere.} Volume $V_2=\tfrac23\pi r^3=\tfrac23\pi(27)=18\pi$; its centre of mass is $\tfrac{3r}{8}=\tfrac98$ above its flat face, i.e.\ at height
\[
x_2=8+\frac98=\frac{73}{8}.
\]

\textbf{Composite.} With common density $\rho$ (which cancels):
\[
\bar{x}=\frac{V_1 x_1+V_2 x_2}{V_1+V_2}=\frac{72\pi(4)+18\pi\left(\tfrac{73}{8}\right)}{72\pi+18\pi}
=\frac{288+\tfrac{1314}{8}}{90}=\frac{288+164.25}{90}=\frac{452.25}{90}.
\]
\[
\bar{x}=\frac{1809}{360}=\frac{201}{40}=5.025\ \text{cm}.
\]
The centre of mass is $5.025$~cm above the base — above the mid-height of the cylinder, as expected since the hemisphere adds mass at the top.
\end{exoresolu}

\begin{exoresolu}[Solid of revolution]
Find the centre of mass of the solid formed by rotating the region bounded by $y=\sqrt{x}$, $x=4$ and the $x$-axis about the $x$-axis.
\tcblower
Volume: $V=\pi\int_0^4 y^2\,dx=\pi\int_0^4 x\,dx=\pi\left[\frac{x^2}{2}\right]_0^4=8\pi$.
Moment about $yz$-plane: $\int_0^4 \pi x y^2\,dx=\pi\int_0^4 x^2\,dx=\pi\left[\frac{x^3}{3}\right]_0^4=\frac{64\pi}{3}$.
\[
\bar{x}=\frac{64\pi/3}{8\pi}=\frac{8}{3}\approx 2.667.
\]
By symmetry $\bar{y}=\bar{z}=0$. Centre of mass: $\left(\frac{8}{3},0,0\right)$.
\end{exoresolu}

\courssub{Applications of Centre of Mass in Mechanics}

\begin{exoresolu}[Equilibrium of a composite body on an inclined plane]
A uniform solid consists of a cylinder of radius $r$ and height $h$ with a cone of the same radius and height $h$ attached to one end. The solid is placed on a rough inclined plane with the cylindrical face on the plane. Find the maximum angle of inclination for which the solid does not topple.
\tcblower
By symmetry, $G$ lies on the axis. Measure from the base of the cylinder.
Cylinder: volume $V_1=\pi r^2 h$, COM at $h/2$.
Cone: volume $V_2=\frac13\pi r^2 h$, COM at $h + h/4 = 5h/4$ from cylinder base.
Total volume $V=\frac43\pi r^2 h$.
\[
\bar{x}=\frac{\pi r^2 h(h/2) + \frac13\pi r^2 h(5h/4)}{\frac43\pi r^2 h}
=\frac{h^2/2 + 5h^2/12}{4/3}
=\frac{6h^2/12 + 5h^2/12}{4/3}
=\frac{11h^2/12}{4/3}=\frac{11h}{16}.
\]
The base is a circle of radius $r$. Toppling occurs when the vertical through $G$ passes through the edge of the base. At limiting equilibrium, $\tan\theta = \frac{r}{\bar{x}} = \frac{r}{11h/16} = \frac{16r}{11h}$.
\[
\theta_{\max}=\arctan\left(\frac{16r}{11h}\right).
\]
\end{exoresolu}

\begin{exoresolu}[Suspension of a 3D solid]
A uniform solid is a hemisphere of radius $r$ with a cone of the same base radius and height $r$ attached to its flat face. The solid is suspended from a point on the rim of the hemisphere. Find the angle the axis makes with the vertical.
\tcblower
Place hemisphere flat face on $xy$-plane, centre at origin, dome in $z>0$. Cone attached on flat face, vertex at $(0,0,-r)$. Axis is $z$-axis.
Hemisphere: volume $\frac23\pi r^3$, COM at $z=\frac{3r}{8}$.
Cone: volume $\frac13\pi r^3$, COM at $z=-\frac{r}{4}$ (since $\frac14$ height above base, base at $z=0$).
Total volume $=\pi r^3$.
\[
\bar{z}=\frac{\frac23\pi r^3(\frac{3r}{8}) + \frac13\pi r^3(-\frac{r}{4})}{\pi r^3}
=\frac{\frac{r}{4} - \frac{r}{12}}{1}=\frac{3r-r}{12}=\frac{r}{6}.
\]
Suspension point: on rim of hemisphere, e.g. $(r,0,0)$. Vector from suspension point to $G$: $(-r, 0, r/6)$.
This vector is vertical. The axis is the $z$-axis $(0,0,1)$. Angle $\phi$ between axis and vertical:
\[
\cos\phi = \frac{(0,0,1)\cdot(-r,0,r/6)}{\sqrt{r^2+(r/6)^2}\cdot 1} = \frac{r/6}{r\sqrt{1+1/36}} = \frac{1/6}{\sqrt{37/36}} = \frac{1}{\sqrt{37}}.
\]
\[
\phi = \arccos\left(\frac{1}{\sqrt{37}}\right) \approx 80.5^{\circ}.
\]
\end{exoresolu}

\begin{attention}[Common Traps to Avoid]
\begin{itemize}
\item Forgetting that a hole must be counted with a \cle{negative} area or mass.
\item Mixing up the factor $\tfrac12$ in $\bar{y}=\dfrac{1}{2A}\int y^2\,dx$ for centroids of regions.
\item Quoting $\tfrac{4r}{3\pi}$ (lamina) where $\tfrac{3r}{8}$ (solid hemisphere) is required — learn which result applies to which object.
\item In suspension problems, drawing the vertical through the suspension point \emph{before} locating $G$: always compute $G$ first, then join it to the point of suspension.
\item Using masses instead of areas for a \emph{uniform} lamina is fine only if the same surface density is used throughout; areas are simpler and equivalent.
\item For solids of revolution, forgetting that $\bar{y}=\bar{z}=0$ by symmetry and only computing $\bar{x}$.
\item Confusing centroid of a \emph{curve} (arc) with centroid of a \emph{lamina} (area) or \emph{solid} (volume).
\end{itemize}
\end{attention}

\begin{exercice}[Practice Problems]
\begin{enumerate}
\item Particles of masses $2$, $3$, $4$, $5$~kg are placed at the vertices of a rectangle of sides $3$~m and $4$~m. Find the distance of the centre of mass from the $3$~m side.
\item A uniform lamina is a square of side $a$ with a quadrant of a circle of radius $a/2$ removed from one corner. Find the centre of mass of the remaining lamina.
\item Find the centroid of the region bounded by $y=\sin x$, $y=0$, $x=0$, $x=\pi$.
\item A uniform solid consists of a cylinder of radius $r$ and height $2r$ with a hemisphere of radius $r$ attached to each end. Find the centre of mass of the solid.
\item A uniform lamina in the shape of an equilateral triangle of side $6$~cm has a circular hole of radius $1$~cm drilled at its centroid. Find the new centre of mass.
\item A solid is formed by rotating the region bounded by $y=x^2$, $y=4$ about the $y$-axis. Find the height of its centre of mass above the base.
\item A composite lamina consists of a rectangle $8\times 4$ and a semicircle of radius $2$ attached to one $8$~cm side. The lamina is suspended from the midpoint of the opposite $8$~cm side. Find the angle the base makes with the horizontal.
\item A uniform wire is bent into the shape of three sides of a square of side $L$. Find the distance of its centre of mass from the open side.
\end{enumerate}
\end{exercice}

\begin{corrige}[Exercise 1]
Place rectangle with vertices $(0,0)$, $(4,0)$, $(4,3)$, $(0,3)$. Masses: $2$ at $(0,0)$, $3$ at $(4,0)$, $4$ at $(4,3)$, $5$ at $(0,3)$.
$M=14$. $\bar{x}=\frac{2(0)+3(4)+4(4)+5(0)}{14}=\frac{28}{14}=2$. $\bar{y}=\frac{2(0)+3(0)+4(3)+5(3)}{14}=\frac{27}{14}\approx 1.93$~m from the $3$~m side (the $x$-axis).
\end{corrige}

\begin{corrige}[Exercise 2]
Square: area $a^2$, centre $(a/2,a/2)$. Quadrant removed: area $-\frac{\pi a^2}{16}$, centroid at $(\frac{4(a/2)}{3\pi},\frac{4(a/2)}{3\pi})=(\frac{2a}{3\pi},\frac{2a}{3\pi})$ from the corner. If corner is at origin, quadrant centroid is $(\frac{2a}{3\pi},\frac{2a}{3\pi})$.
Total area $=a^2(1-\pi/16)$.
$\bar{x}=\frac{a^2(a/2) - \frac{\pi a^2}{16}(\frac{2a}{3\pi})}{a^2(1-\pi/16)} = \frac{a/2 - a/24}{1-\pi/16} = \frac{11a/24}{(16-\pi)/16} = \frac{11a}{24}\cdot\frac{16}{16-\pi} = \frac{22a}{3(16-\pi)}$.
By symmetry $\bar{y}=\bar{x}$.
\end{corrige}

\begin{corrige}[Exercise 3]
Area $A=\int_0^\pi \sin x\,dx = 2$.
$\bar{x}=\frac{1}{2}\int_0^\pi x\sin x\,dx = \frac{1}{2}[-x\cos x+\sin x]_0^\pi = \frac{\pi}{2}$.
$\bar{y}=\frac{1}{4}\int_0^\pi \sin^2 x\,dx = \frac{1}{4}\int_0^\pi \frac{1-\cos 2x}{2}\,dx = \frac{1}{8}[\pi-0] = \frac{\pi}{8}$.
Centroid: $(\pi/2, \pi/8)$.
\end{corrige}

\begin{corrige}[Exercise 4]
By symmetry, COM is at the centre of the cylinder (midpoint of axis). Cylinder COM at centre. Each hemisphere COM at $3r/8$ from its flat face. The two hemispheres are symmetric about the cylinder centre, so their combined COM is at the cylinder centre. Overall COM at the centre of the cylinder.
\end{corrige}

\begin{corrige}[Exercise 5]
Equilateral triangle side $6$: area $9\sqrt{3}$, centroid at centre. Hole: area $-\pi$, centroid at same point (centroid of triangle). Since hole is at the COM of the triangle, the COM of the remaining lamina is unchanged — still at the centroid of the triangle.
\end{corrige}

\begin{corrige}[Exercise 6]
Region: $x$ from $-2$ to $2$, $y$ from $x^2$ to $4$. Rotate about $y$-axis.
Use horizontal strips (washers). At height $y$, radius $x=\sqrt{y}$. Area of cross-section $\pi y$.
Volume $V=\pi\int_0^4 y\,dy = 8\pi$.
Moment about base ($y=0$): $\int_0^4 \pi y \cdot y\,dy = \pi\int_0^4 y^2\,dy = \frac{64\pi}{3}$.
$\bar{y} = \frac{64\pi/3}{8\pi} = \frac{8}{3}$.
\end{corrige}

\begin{corrige}[Exercise 7]
Rectangle: area $32$, centre $(4,2)$ if base on $x$-axis from $0$ to $8$, height $4$.
Semicircle on top: radius $2$, area $2\pi$, centroid at $(4, 4+\frac{8}{3\pi})$.
Total area $=32+2\pi$.
$\bar{x}=4$ (symmetry).
$\bar{y}=\frac{32(2)+2\pi(4+8/(3\pi))}{32+2\pi} = \frac{64+8\pi+16/3}{32+2\pi} = \frac{208/3+8\pi}{32+2\pi}$.
Suspended from midpoint of bottom side: $(4,0)$. Vector to $G$: $(0, \bar{y})$. This is vertical, so the base (horizontal) makes $0^\circ$ with horizontal? Wait: suspension from midpoint of bottom side means the point is $(4,0)$. $G$ is at $(4,\bar{y})$. The line joining them is vertical. The base is horizontal. So the base remains horizontal. Angle with horizontal = $0^\circ$. (The lamina hangs with base horizontal because suspension point is directly below COM horizontally.)
\end{corrige}

\begin{corrige}[Exercise 8]
Wire: three sides of square. Total length $3L$. By symmetry, COM lies on midline $x=L/2$.
Bottom side: length $L$, centre $(L/2, 0)$.
Left side: length $L$, centre $(0, L/2)$.
Right side: length $L$, centre $(L, L/2)$.
$\bar{x}=\frac{L(L/2)+L(0)+L(L)}{3L} = \frac{L/2+0+L}{3} = \frac{3L/2}{3} = L/2$.
$\bar{y}=\frac{L(0)+L(L/2)+L(L/2)}{3L} = \frac{L}{3}$.
Distance from open side (top, $y=L$) is $L - L/3 = 2L/3$.
\end{corrige}

\newpage

\coursec{Exercises}

\courssub{I check my knowledge}

\begin{exercice}[MCQ: Centre of mass of particles]
Three particles of masses \SI{2}{kg}, \SI{3}{kg} and \SI{5}{kg} are placed at the points with coordinates $(1,2)$, $(4,-1)$ and $(-2,3)$ respectively.  
Which of the following gives the coordinates of the centre of mass of the system?
\begin{enumerate}
\item $(0.5,\; 1.5)$
\item $(0.7,\; 1.3)$
\item $(0.9,\; 1.1)$
\item $(1.1,\; 0.9)$
\end{enumerate}
\end{exercice}

\begin{exercice}[True/False: Properties of centre of mass]
State whether each of the following statements is true or false. Justify your answer briefly.
\begin{enumerate}
\item The centre of mass of a uniform rigid body always lies within the material of the body.
\item For a system of particles, the position vector of the centre of mass is the weighted average of the position vectors of the particles, with the masses as weights.
\item If a uniform lamina is symmetric about a line, its centre of mass lies on that line.
\item The centre of mass of a composite body can be found by treating each component as a particle of its total mass placed at its own centre of mass.
\end{enumerate}
\end{exercice}

\begin{exercice}[Definition and direct calculation]
A uniform lamina consists of a rectangle $ABCD$ with $AB = \SI{8}{cm}$, $BC = \SI{6}{cm}$, and a semicircle of radius \SI{3}{cm} attached to the side $CD$ (the semicircle's diameter coincides with $CD$). The lamina is made of the same material throughout.  
Find the distance of the centre of mass of the lamina from the side $AB$.
\end{exercice}

\begin{exercice}[Position vector method]
Three particles have position vectors $\vec{r}_1 = 2\vec{i} + \vec{j}$, $\vec{r}_2 = -\vec{i} + 4\vec{j}$ and $\vec{r}_3 = 3\vec{i} - 2\vec{k}$ with masses $m$, $2m$ and $3m$ respectively.  
\begin{enumerate}
\item Find the position vector of the centre of mass of the system.
\item A fourth particle of mass $M$ is placed at the origin. Determine the value of $M$ (in terms of $m$) such that the new centre of mass lies on the $x$-axis.
\end{enumerate}
\end{exercice}

\courssub{I practise}

\begin{exercice}[Composite lamina: rectangle and semicircle]
A uniform lamina is made of a rectangle $8\ \mathrm{cm}$ by $6\ \mathrm{cm}$ with a semicircle of radius $3\ \mathrm{cm}$ attached to one of the shorter sides (the diameter of the semicircle coincides with that side). The lamina is freely suspended from the midpoint of the opposite shorter side. Find the angle that this side makes with the vertical when the lamina hangs in equilibrium.
\end{exercice}

\begin{exercice}[Circular disc with a hole]
A uniform circular disc of radius \SI{10}{cm} has a circular hole of radius \SI{4}{cm} drilled in it. The centre of the hole is \SI{3}{cm} from the centre of the disc. Find the distance of the centre of mass of the remaining plate from the centre of the original disc, and state on which side of the centre it lies.
\end{exercice}

\begin{exercice}[Suspended triangular lamina]
A uniform triangular lamina $ABC$ has side lengths $AB = \SI{13}{cm}$, $BC = \SI{14}{cm}$, $CA = \SI{15}{cm}$. The lamina is freely suspended from vertex $A$. Find the angle that side $AB$ makes with the vertical when the lamina hangs in equilibrium.
\end{exercice}

\begin{exercice}[Centroid by integration]
Find, by integration, the coordinates of the centroid of the region bounded by the curve $y = 4 - x^2$ and the $x$-axis. Verify that the centroid lies on the axis of symmetry of the region.
\end{exercice}

\courssub{I prepare for the GCE}

\begin{exercice}[Composite lamina suspended from a corner]
A uniform lamina is formed from a square of side \SI{20}{cm} by removing a quadrant of a circle of radius \SI{10}{cm} from one corner (the quadrant's centre is at that corner). The lamina is freely suspended from the opposite corner of the square. Calculate the angle of inclination of the diagonal through the suspension point to the vertical when the lamina hangs in equilibrium. \hfill [12 marks]
\end{exercice}

\begin{exercice}[Solid cylinder with cone]
A uniform solid consists of a right circular cylinder of height \SI{12}{cm} and radius \SI{5}{cm}, surmounted by a right circular cone of the same radius and height \SI{9}{cm}. The base of the cylinder is on a horizontal table. Find the height of the centre of mass of the solid above the base. \hfill [10 marks]
\end{exercice}

\begin{exercice}[Block on inclined plane: slide or topple]
A uniform rectangular block of base \SI{40}{cm} and height \SI{90}{cm} stands on a plane inclined at an angle $\theta$ to the horizontal. The coefficient of friction between the block and the plane is $\mu = 0.5$. Determine whether the block will slide or topple first as $\theta$ is gradually increased from zero. Find the critical angle for each mode of failure. \hfill [13 marks]
\end{exercice}

\newpage

\coursec{Solutions to the exercises}

\coursec{Corrections of Exercises — Centre of Mass of Uniform Rigid Bodies}

\begin{corrige}[Exercise 1 — MCQ: Centre of mass of particles]
The total mass is $M = 2+3+5 = \SI{10}{kg}$. Applying the centre of mass formulae $\bar{x}=\dfrac{\sum m_i x_i}{\sum m_i}$ and $\bar{y}=\dfrac{\sum m_i y_i}{\sum m_i}$:
\begin{align*}
\bar{x} &= \frac{2(1)+3(4)+5(-2)}{10} = \frac{2+12-10}{10} = \frac{4}{10} = 0.4,\\[2pt]
\bar{y} &= \frac{2(2)+3(-1)+5(3)}{10} = \frac{4-3+15}{10} = \frac{16}{10} = 1.6.
\end{align*}
The centre of mass is therefore at $(0.4,\; 1.6)$.

\textbf{Conclusion:} \emph{none} of the proposed options (a)--(d) is correct; the true answer is $(0.4,\,1.6)$. This is a classic MCQ trap: options (a) and (b) are obtained by slips in the signs of the negative coordinates, and options (c), (d) by dividing by the wrong total mass. At the GCE, always recompute rather than picking the ``closest'' answer.
\end{corrige}

\begin{corrige}[Exercise 2 — True/False: Properties of centre of mass]
\begin{enumerate}
\item \textbf{False.} The centre of mass of a uniform body need not lie within the material of the body. Counter-examples: a uniform ring or annulus has its centre of mass at its geometric centre, which is in the hole; a uniform boomerang-shaped or L-shaped lamina can also have its centre of mass outside the material.
\item \textbf{True.} By definition, $\vec{r}_G = \dfrac{\sum_i m_i \vec{r}_i}{\sum_i m_i} = \sum_i \dfrac{m_i}{M}\vec{r}_i$, which is exactly a weighted average of the position vectors with the masses as weights.
\item \textbf{True.} If a uniform lamina is symmetric about a line $\ell$, then for every element of mass at a point $P$ there is an identical element at the mirror point $P'$. The contributions perpendicular to $\ell$ cancel in pairs, so the centre of mass lies on $\ell$.
\item \textbf{True.} This is the composite-body principle: each component of mass $m_i$ and centre of mass $G_i$ may be replaced by a single particle of mass $m_i$ at $G_i$, because the first moment $\sum m_i \vec{r}_{G_i}$ of the reduced system equals that of the original body. The centre of mass of the whole is then found from these ``equivalent particles''.
\end{enumerate}
\end{corrige}

\begin{corrige}[Exercise 3 — Definition and direct calculation]
\textbf{Setting up.} Place the lamina with $AB$ along the $y$-axis, the rectangle occupying $0 \le x \le 8$, $0 \le y \le 6$. Then $CD$ is the line $x = 8$ and the semicircle of radius $r = \SI{3}{cm}$ bulges out to the right of $CD$. By symmetry about the horizontal mid-line, we only need the $x$-coordinate of the centre of mass, which gives the distance from $AB$.

\textbf{Components.} Since the lamina is uniform, masses are proportional to areas; we work with areas directly.
\begin{itemize}
\item Rectangle: area $A_1 = 8 \times 6 = \SI{48}{cm^2}$, centre of mass at $x_1 = \SI{4}{cm}$ from $AB$.
\item Semicircle: area $A_2 = \tfrac12 \pi r^2 = \tfrac{9\pi}{2}\ \mathrm{cm^2}$. The centroid of a semicircle lies at distance $\dfrac{4r}{3\pi} = \dfrac{4}{\pi}$ from its diameter, so $x_2 = 8 + \dfrac{4}{\pi}$ cm from $AB$.
\end{itemize}

\textbf{Composite formula.}
\begin{align*}
\bar{x} &= \frac{A_1 x_1 + A_2 x_2}{A_1 + A_2}
= \frac{48(4) + \dfrac{9\pi}{2}\left(8 + \dfrac{4}{\pi}\right)}{48 + \dfrac{9\pi}{2}}\\[3pt]
&= \frac{192 + 36\pi + 18}{48 + \frac{9\pi}{2}}
= \frac{210 + 36\pi}{48 + \frac{9\pi}{2}}
= \frac{420 + 72\pi}{96 + 9\pi}
= \frac{140 + 24\pi}{32 + 3\pi}\ \mathrm{cm}.
\end{align*}
Numerically, $\bar{x} = \dfrac{140 + 24\pi}{32 + 3\pi} \approx \SI{5.20}{cm}$.

\textbf{Answer:} the centre of mass is approximately $\SI{5.2}{cm}$ from the side $AB$ (exact value $\dfrac{140+24\pi}{32+3\pi}\ \mathrm{cm}$).
\end{corrige}

\begin{corrige}[Exercise 4 — Position vector method]
\textbf{(a) Position vector of the centre of mass.} The total mass is $m + 2m + 3m = 6m$. Then
\begin{align*}
\vec{r}_G &= \frac{m\vec{r}_1 + 2m\vec{r}_2 + 3m\vec{r}_3}{6m}\\
&= \frac{1}{6}\Big[(2\vec{i}+\vec{j}) + 2(-\vec{i}+4\vec{j}) + 3(3\vec{i}-2\vec{k})\Big]\\
&= \frac{1}{6}\Big[(2-2+9)\vec{i} + (1+8)\vec{j} + (-6)\vec{k}\Big]\\
&= \frac{1}{6}\big(9\vec{i} + 9\vec{j} - 6\vec{k}\big)
= 1.5\vec{i} + 1.5\vec{j} - \vec{k}.
\end{align*}

\textbf{(b) Adding a particle at the origin.} With a fourth particle of mass $M$ at $\vec{r}_4 = \vec{0}$, the new centre of mass is
\[
\vec{r}_{G'} = \frac{6m\,\vec{r}_G + M\vec{0}}{6m + M} = \frac{9m\vec{i} + 9m\vec{j} - 6m\vec{k}}{6m+M}.
\]
For $G'$ to lie on the $x$-axis we need both $\bar{y} = \dfrac{9m}{6m+M} = 0$ and $\bar{z} = \dfrac{-6m}{6m+M} = 0$. Since the numerators $9m$ and $-6m$ are fixed and non-zero, no finite value of $M$ can make either coordinate vanish.

\textbf{Conclusion:} there is \emph{no} value of $M$ for which the new centre of mass lies on the $x$-axis. Adding mass at the origin only pulls $G'$ towards the origin along the line $OG$; it can never bring it onto the $x$-axis. (The question illustrates that the moments $\sum m_i y_i$ and $\sum m_i z_i$ are unchanged by a mass placed at the origin.)
\end{corrige}

\begin{corrige}[Exercise 5 — Composite lamina: rectangle and semicircle]
\textbf{Locating the centre of mass.} This is the same lamina as in Exercise 3: a rectangle $8 \times 6$ cm with a semicircle of radius \SI{3}{cm} on one of the \SI{6}{cm} sides. Let $AB$ be the opposite shorter side (the suspension side). From Exercise 3, the centre of mass $G$ lies on the axis of symmetry of the lamina, at distance
\[
\bar{x} = \frac{140 + 24\pi}{32 + 3\pi} \approx \SI{5.20}{cm} \quad \text{from } AB,
\]
measured along the perpendicular bisector of $AB$.

\textbf{Equilibrium of the suspended lamina.} The lamina is suspended from $O$, the midpoint of $AB$. In equilibrium, $G$ hangs vertically below $O$. But $G$ lies on the perpendicular bisector of $AB$, so the line $OG$ is perpendicular to $AB$. Hence $OG$ vertical forces $AB$ to be \emph{horizontal}.

\textbf{Answer:} the side $AB$ makes an angle of $\boxed{0^\circ}$ with the vertical (it hangs horizontally). The computation of $\bar{x}$ is still essential: it confirms that $G$ is inside the lamina on the symmetry axis, but because the suspension point is \emph{on} that same axis of symmetry, the lamina hangs with its axis vertical and $AB$ horizontal.
\end{corrige}

\begin{corrige}[Exercise 6 — Circular disc with a hole]
\textbf{Method: negative mass.} Treat the plate as the full disc plus a ``negative'' disc for the hole. Let $O$ be the centre of the original disc and let the centre of the hole be at $x = +\SI{3}{cm}$ on the $x$-axis. Since the plate is uniform, masses are proportional to areas.

\begin{itemize}
\item Full disc: area $A_1 = \pi(10)^2 = 100\pi\ \mathrm{cm^2}$, centre of mass at $x_1 = 0$.
\item Hole (removed): area $A_2 = \pi(4)^2 = 16\pi\ \mathrm{cm^2}$, at $x_2 = \SI{3}{cm}$.
\end{itemize}

\textbf{Computation.}
\[
\bar{x} = \frac{A_1 x_1 - A_2 x_2}{A_1 - A_2} = \frac{100\pi(0) - 16\pi(3)}{100\pi - 16\pi} = \frac{-48}{84} = -\frac{4}{7}\ \mathrm{cm} \approx \SI{-0.57}{cm}.
\]

\textbf{Answer:} the centre of mass of the remaining plate is $\dfrac{4}{7}\ \mathrm{cm} \approx \SI{0.57}{cm}$ from the centre of the original disc, on the side \emph{opposite} to the hole (the negative sign indicates the side away from the hole, as expected physically: removing material shifts the balance point away from the hole).
\end{corrige}

\begin{corrige}[Exercise 7 — Suspended triangular lamina]
\textbf{Key fact.} The centre of mass of a uniform triangular lamina is its centroid $G$, which lies on each median, two-thirds of the way from a vertex to the midpoint of the opposite side. Suspended from $A$, the lamina hangs with $G$ vertically below $A$; hence the \emph{median} from $A$ is vertical, and the required angle is the angle between $AB$ and the median from $A$.

\textbf{Length of the median from $A$.} With $a = BC = 14$, $b = CA = 15$, $c = AB = 13$, the median formula gives
\[
m_a = \tfrac12\sqrt{2b^2 + 2c^2 - a^2} = \tfrac12\sqrt{2(225) + 2(169) - 196} = \tfrac12\sqrt{592} = \sqrt{148} = 2\sqrt{37}\ \mathrm{cm} \approx \SI{12.17}{cm}.
\]

\textbf{Angle between $AB$ and the median.} Let $M$ be the midpoint of $BC$, so $BM = \SI{7}{cm}$. In triangle $ABM$, the cosine rule at $A$ gives
\[
\cos\widehat{BAM} = \frac{AB^2 + AM^2 - BM^2}{2\cdot AB \cdot AM} = \frac{169 + 148 - 49}{2(13)(\sqrt{148})} = \frac{268}{26\sqrt{148}} = \frac{67}{13\sqrt{37}} \approx 0.8473.
\]
\[
\widehat{BAM} \approx 32.1^\circ.
\]

\textbf{Answer:} side $AB$ makes an angle of approximately $32.1^\circ$ with the vertical.
\end{corrige}

\begin{corrige}[Exercise 8 — Centroid by integration]
The curve $y = 4 - x^2$ meets the $x$-axis where $4 - x^2 = 0$, i.e. $x = \pm 2$. The region is symmetric about the $y$-axis.

\textbf{Area.}
\[
A = \int_{-2}^{2} (4 - x^2)\,dx = \left[4x - \frac{x^3}{3}\right]_{-2}^{2} = \left(8 - \tfrac{8}{3}\right) - \left(-8 + \tfrac{8}{3}\right) = \frac{32}{3}.
\]

\textbf{First coordinate.} By symmetry (or by computing $\int_{-2}^{2} x(4-x^2)\,dx = 0$, an odd integrand over a symmetric interval),
\[
\bar{x} = 0.
\]

\textbf{Second coordinate.} For a region under a curve, $\bar{y} = \dfrac{1}{A}\displaystyle\int_{-2}^{2} \tfrac12 y^2\,dx$:
\begin{align*}
\int_{-2}^{2} \tfrac12 (4-x^2)^2\,dx &= \tfrac12\int_{-2}^{2}\left(16 - 8x^2 + x^4\right)dx\\
&= \tfrac12\left[16x - \frac{8x^3}{3} + \frac{x^5}{5}\right]_{-2}^{2}\\
&= \tfrac12 \times 2\left(32 - \frac{64}{3} + \frac{32}{5}\right) = \frac{256}{15}.
\end{align*}
Hence
\[
\bar{y} = \frac{256/15}{32/3} = \frac{256 \times 3}{15 \times 32} = \frac{8}{5} = 1.6.
\]

\textbf{Answer and verification:} the centroid is $\left(0,\; \dfrac{8}{5}\right)$. Since $\bar{x} = 0$, the centroid lies on the $y$-axis, which is precisely the axis of symmetry of the parabolic region — as required. $\blacksquare$
\end{corrige}

\begin{corrige}[Exercise 9 — Composite lamina suspended from a corner (GCE type, 12 marks)]
\textbf{Model.} Place the square with the suspension corner at the origin $O=(0,0)$, sides along the axes, so the square is $[0,20]\times[0,20]$. The quadrant of radius \SI{10}{cm} is removed from the opposite corner $(20,20)$.

\textbf{Component data} (uniform lamina, so masses $\propto$ areas):
\begin{itemize}
\item Square: $A_1 = 400\ \mathrm{cm^2}$, centroid $(10,10)$.
\item Quadrant removed: $A_2 = \tfrac14\pi(10)^2 = 25\pi\ \mathrm{cm^2}$. The centroid of a quadrant of radius $r$ is at distance $\dfrac{4r}{3\pi} = \dfrac{40}{3\pi}$ cm from each straight edge, so its centroid is at $\left(20 - \dfrac{40}{3\pi},\; 20 - \dfrac{40}{3\pi}\right) \approx (15.756,\; 15.756)$.
\end{itemize}

\textbf{Centre of mass of the remainder} (negative-area method):
\begin{align*}
\bar{x} &= \frac{400(10) - 25\pi\left(20 - \frac{40}{3\pi}\right)}{400 - 25\pi}
= \frac{4000 - 500\pi + \frac{1000}{3}}{400 - 25\pi}\\
&= \frac{\frac{13000 - 1500\pi}{3}}{400 - 25\pi}
= \frac{13000 - 1500\pi}{1200 - 75\pi}
= \frac{520 - 60\pi}{48 - 3\pi}\ \mathrm{cm} \approx \SI{8.59}{cm}.
\end{align*}
By the symmetry of the figure about the diagonal $y = x$, $\bar{y} = \bar{x} \approx \SI{8.59}{cm}$.

\textbf{Hanging equilibrium.} The lamina is suspended from $O$; in equilibrium $G$ is vertically below $O$. Since $\bar{x} = \bar{y}$, the point $G$ lies exactly on the diagonal through $O$ (the line $y = x$). Therefore this diagonal hangs vertically.

\textbf{Answer:} the diagonal through the suspension point makes an angle of $\boxed{0^\circ}$ with the vertical — it is itself vertical in equilibrium.

\textbf{Indicative mark scheme (12 marks):}
\begin{itemize}
\item M1: correct areas of square and quadrant; A1: both correct ($400$, $25\pi$).
\item M1: centroid of quadrant at $\frac{4r}{3\pi}$ from each edge; A1: coordinates $(15.76, 15.76)$.
\item M1: negative-moment equation for $\bar{x}$ (or $\bar{y}$); A1: correct unsimplified expression; A1: $\bar{x} \approx 8.59$ cm.
\item M1: use of symmetry to state $\bar{y} = \bar{x}$; A1: $G$ on the diagonal.
\item M1: equilibrium condition ($G$ vertically below point of suspension); A1: diagonal vertical.
\item E1: conclusion $\theta = 0^\circ$ with justification.
\end{itemize}
\textbf{Examiner expectations:} candidates must quote the quadrant centroid formula correctly, handle the removed part with a \emph{negative} moment, and recognise the symmetry. A common error is to suspend from the wrong corner or to subtract the quadrant's moment with the wrong sign.
\end{corrige}

\begin{corrige}[Exercise 10 — Solid cylinder with cone (GCE type, 10 marks)]
\textbf{Component data.} Measure heights from the base of the cylinder. Volumes (uniform solid, masses $\propto$ volumes):
\begin{itemize}
\item Cylinder: $V_1 = \pi r^2 h = \pi(5^2)(12) = 300\pi\ \mathrm{cm^3}$, centre of mass at $y_1 = \SI{6}{cm}$.
\item Cone: $V_2 = \tfrac13\pi r^2 h = \tfrac13\pi(25)(9) = 75\pi\ \mathrm{cm^3}$. The centre of mass of a solid cone is $\tfrac14$ of its height above its base, so $y_2 = 12 + \tfrac{9}{4} = \SI{14.25}{cm}$.
\end{itemize}

\textbf{Composite formula.}
\begin{align*}
\bar{y} &= \frac{V_1 y_1 + V_2 y_2}{V_1 + V_2}
= \frac{300\pi(6) + 75\pi(14.25)}{300\pi + 75\pi}\\[2pt]
&= \frac{1800 + 1068.75}{375} = \frac{2868.75}{375} = 7.65.
\end{align*}

\textbf{Answer:} the centre of mass of the solid is $\SI{7.65}{cm}$ above the base.

\textbf{Indicative mark scheme (10 marks):}
\begin{itemize}
\item M1 A1: volume of cylinder and position of its centre of mass.
\item M1 A1: volume of cone; B1: cone centre of mass at $h/4$ above its base, i.e. \SI{14.25}{cm}.
\item M1: correct weighted-mean equation; A1: correct unsimplified numerator and denominator.
\item A1: $\bar{y} = 7.65$ cm; B1: statement of height above the base with units.
\item E1: consistency remark (e.g. $6 < 7.65 < 12$: the heavy cylinder pulls $G$ below the join, as expected).
\end{itemize}
\textbf{Examiner expectations:} quote the standard result for the cone correctly (a frequent error is to use $h/3$ instead of $h/4$), and remember to add the cylinder height when locating the cone's centre of mass.
\end{corrige}

\begin{corrige}[Exercise 11 — Block on inclined plane: slide or topple (GCE type, 13 marks)]
\textbf{Setting up.} The block has base $b = \SI{40}{cm}$ and height $h = \SI{90}{cm}$; its centre of mass $G$ is at its geometric centre, i.e. \SI{20}{cm} horizontally from each vertical face and \SI{45}{cm} above the base. Let $\theta$ be the angle of inclination, increased slowly from $0$.

\textbf{Condition for sliding.} Resolve the weight $W = mg$ along and normal to the plane: driving force $W\sin\theta$, normal reaction $R = W\cos\theta$. Sliding begins when the driving force reaches limiting friction:
\[
W\sin\theta = \mu R = \mu W\cos\theta \quad\Longrightarrow\quad \tan\theta = \mu = 0.5.
\]
\[
\theta_{\text{slide}} = \arctan(0.5) \approx 26.6^\circ.
\]

\textbf{Condition for toppling.} Toppling begins when the vertical line through $G$ passes through the lower downhill edge (corner) $P$ of the base — equivalently, when the reaction at the upper edge becomes zero and the block is on the point of rotating about $P$. In the limiting position, the line $PG$ makes an angle $\theta$ with the normal to the plane, and from the geometry of the rectangle
\[
\tan\theta = \frac{\text{half base}}{\text{half height}} = \frac{b/2}{h/2} = \frac{b}{h} = \frac{40}{90} = \frac{4}{9}.
\]
\[
\theta_{\text{topple}} = \arctan\!\left(\tfrac{4}{9}\right) \approx 24.0^\circ.
\]

\textbf{Comparison and conclusion.} Since
\[
\theta_{\text{topple}} \approx 24.0^\circ < \theta_{\text{slide}} \approx 26.6^\circ,
\]
the toppling condition is reached first as $\theta$ increases. \textbf{The block topples before it slides}, at the critical angle $\theta \approx 24.0^\circ$.

\textbf{Indicative mark scheme (13 marks):}
\begin{itemize}
\item B1: $G$ at the centre of the block; diagram with forces ($W$, $R$, friction $F$).
\item M1: resolving parallel and perpendicular to the plane; A1: $F = W\sin\theta$, $R = W\cos\theta$.
\item M1: sliding condition $F = \mu R$; A1: $\tan\theta = 0.5$; A1: $\theta_{\text{slide}} \approx 26.6^\circ$.
\item M1: toppling condition — vertical through $G$ passes through the lower corner; M1: correct geometry $\tan\theta = \dfrac{b/2}{h/2}$; A1: $\tan\theta = \tfrac{4}{9}$; A1: $\theta_{\text{topple}} \approx 24.0^\circ$.
\item M1: comparison of the two critical angles; A1: correct conclusion (topples first); E1: physical comment, e.g. a tall narrow block with moderate friction topples rather than slides.
\end{itemize}
\textbf{Examiner expectations:} a clear force diagram; the toppling condition derived from moments about the lower corner (or from the line of action of the weight); an explicit comparison. Common errors: using $\tan\theta = h/b$ (inverted ratio), or comparing $\mu$ with $h/b$ instead of with $b/h$.
\end{corrige}

\newpage

\coursec{Summary sheet}

\begin{popfigure}
\begin{center}\tcbox[colback=popdark,colframe=popdark,coltext=white,arc=9pt,boxsep=4pt,left=6pt,right=6pt]{\bfseries\small Centre of Mass}\end{center}
\vspace{-2pt}
\begin{tcbraster}[raster columns=3,raster equal height=rows,raster column skip=6pt,raster row skip=6pt,enhanced,blankest]
\begin{tcolorbox}[enhanced,colback=white,colframe=popblue,boxrule=0.9pt,arc=3pt,title={Concept \& motion},fonttitle=\bfseries\scriptsize,coltitle=white,colbacktitle=popblue,left=4pt,right=4pt,top=2pt,bottom=2pt,boxsep=1pt,toptitle=1.5pt,bottomtitle=1.5pt,halign=left]
\begin{itemize}[nosep,leftmargin=*,itemsep=1pt,font=\scriptsize]
\item Balance point
\item $M\vec{a}_G=\sum\vec{F}_{\mathrm{ext}}$
\end{itemize}
\end{tcolorbox}
\begin{tcolorbox}[enhanced,colback=white,colframe=popgreen,boxrule=0.9pt,arc=3pt,title={Coordinates},fonttitle=\bfseries\scriptsize,coltitle=white,colbacktitle=popgreen,left=4pt,right=4pt,top=2pt,bottom=2pt,boxsep=1pt,toptitle=1.5pt,bottomtitle=1.5pt,halign=left]
\begin{itemize}[nosep,leftmargin=*,itemsep=1pt,font=\scriptsize]
\item $\bar{x}=\dfrac{\sum m_ix_i}{\sum m_i}$
\item 2D \& 3D
\end{itemize}
\end{tcolorbox}
\begin{tcolorbox}[enhanced,colback=white,colframe=poporange,boxrule=0.9pt,arc=3pt,title={Position vectors},fonttitle=\bfseries\scriptsize,coltitle=white,colbacktitle=poporange,left=4pt,right=4pt,top=2pt,bottom=2pt,boxsep=1pt,toptitle=1.5pt,bottomtitle=1.5pt,halign=left]
\begin{itemize}[nosep,leftmargin=*,itemsep=1pt,font=\scriptsize]
\item $\vec{r}_G=\dfrac{\sum m_i\vec{r}_i}{M}$
\item Particle systems
\end{itemize}
\end{tcolorbox}
\begin{tcolorbox}[enhanced,colback=white,colframe=poppurple,boxrule=0.9pt,arc=3pt,title={Uniform laminae},fonttitle=\bfseries\scriptsize,coltitle=white,colbacktitle=poppurple,left=4pt,right=4pt,top=2pt,bottom=2pt,boxsep=1pt,toptitle=1.5pt,bottomtitle=1.5pt,halign=left]
\begin{itemize}[nosep,leftmargin=*,itemsep=1pt,font=\scriptsize]
\item Symmetry axes
\item Composite: add \& subtract
\end{itemize}
\end{tcolorbox}
\begin{tcolorbox}[enhanced,colback=white,colframe=poppink,boxrule=0.9pt,arc=3pt,title={Rigid bodies},fonttitle=\bfseries\scriptsize,coltitle=white,colbacktitle=poppink,left=4pt,right=4pt,top=2pt,bottom=2pt,boxsep=1pt,toptitle=1.5pt,bottomtitle=1.5pt,halign=left]
\begin{itemize}[nosep,leftmargin=*,itemsep=1pt,font=\scriptsize]
\item Rod: midpoint
\item Standard results
\end{itemize}
\end{tcolorbox}
\begin{tcolorbox}[enhanced,colback=white,colframe=popgold,boxrule=0.9pt,arc=3pt,title={Centroids},fonttitle=\bfseries\scriptsize,coltitle=white,colbacktitle=popgold,left=4pt,right=4pt,top=2pt,bottom=2pt,boxsep=1pt,toptitle=1.5pt,bottomtitle=1.5pt,halign=left]
\begin{itemize}[nosep,leftmargin=*,itemsep=1pt,font=\scriptsize]
\item $\bar{x}=\dfrac{1}{A}\displaystyle\int_a^b x\,y\,dx$
\item Composite areas
\end{itemize}
\end{tcolorbox}
\end{tcbraster}

\legende{Mind map of the chapter: centre of mass of uniform rigid bodies.}
\end{popfigure}

\begin{bilan}
\textbf{Key definitions}
\begin{itemize}
\item The \cle{centre of mass} $G$ of a system is the point where the whole mass may be supposed concentrated for translational motion; it is the balance point of the body.
\item A \cle{uniform} body has constant density, so its mass is proportional to its length, area or volume.
\item A \cle{lamina} is a flat body of negligible thickness; its centre of mass lies in its plane.
\item The \cle{centroid} is the geometric centre of an area (or volume); for a uniform lamina, centroid $=$ centre of mass.
\end{itemize}

\textbf{Formulas to know}
\begin{itemize}
\item Discrete system (coordinates): $\bar{x}=\dfrac{\sum m_ix_i}{\sum m_i}$, $\bar{y}=\dfrac{\sum m_iy_i}{\sum m_i}$, $\bar{z}=\dfrac{\sum m_iz_i}{\sum m_i}$.
\item Position-vector form: $\vec{r}_G=\dfrac{\sum_{i} m_i\vec{r}_i}{\sum_i m_i}=\dfrac{1}{M}\sum_i m_i\vec{r}_i$.
\item Motion of the mass centre: $M\vec{a}_G=\sum \vec{F}_{\mathrm{ext}}$ (internal forces cancel).
\item Uniform lamina bounded by $y=f(x)$, $a\le x\le b$:
$\bar{x}=\dfrac{\displaystyle\int_a^b x\,f(x)\,dx}{\displaystyle\int_a^b f(x)\,dx}$, \quad $\bar{y}=\dfrac{\dfrac12\displaystyle\int_a^b [f(x)]^2\,dx}{\displaystyle\int_a^b f(x)\,dx}$.
\item Standard results: rod $\to$ midpoint; rectangle $\to$ intersection of diagonals; triangle $\to$ centroid at $\tfrac13$ of the height from the base; semicircular lamina of radius $r$ $\to$ $\dfrac{4r}{3\pi}$ from the centre; circular arc (semicircle) $\to$ $\dfrac{2r}{\pi}$.
\end{itemize}

\textbf{Methods}
\begin{itemize}
\item \emph{Symmetry:} if a body has an axis (or plane) of symmetry, $G$ lies on it; two axes of symmetry meet at $G$.
\item \emph{Composite laminae:} split into simple parts, tabulate $m_i$, $x_i$, $y_i$, then apply $\bar{x}=\dfrac{\sum m_ix_i}{\sum m_i}$. For a hole, treat the removed part with \textbf{negative} mass.
\item \emph{Position vectors:} write each $\vec{r}_i$ from the same origin, compute $\sum m_i\vec{r}_i$, divide by $M$.
\item \emph{Integration:} for regions bounded by curves, use strips of area $y\,dx$ and the centroid formulas above.
\end{itemize}

\textbf{Common mistakes}
\begin{itemize}
\item Forgetting that uniform laminae use \emph{areas} (not masses) as weights; masses are $\rho\times$ area.
\item Mixing origins: all coordinates and position vectors must be measured from the \textbf{same} origin.
\item Sign errors when subtracting a cut-out part.
\item Confusing the centroid of a triangle ($\tfrac13$ height from base) with the midpoint.
\item In the $\bar{y}$ formula, forgetting the factor $\tfrac12$ and the square $[f(x)]^2$.
\end{itemize}

\textbf{GCE tips}
\begin{itemize}
\item Draw a clear diagram, mark the axes and the origin before computing.
\item Present composite-body work in a table: part, mass/area, coordinates of $G_i$, moments $m_ix_i$, $m_iy_i$.
\item Always state the final position of $G$ relative to a named point (e.g.\ ``$2.3$ cm from $O$ along $OA$'').
\item Check plausibility: $G$ must lie inside (or on the boundary of) a convex body, and on every axis of symmetry.
\end{itemize}
\end{bilan}

\findecours

\end{document}
